\documentclass[]{report}

\usepackage{import}

\import{../}{settings.tex}

\title{Analyse et topologie}

\author{Cours de Claire Debord, transcrit par SADR TAHOURI Luca}
\date{}

\begin{document}

    \maketitle
    \tableofcontents

    \chapter{Espaces Topologiques}

    \section{Définition d'une topologie}

    \subsection{Espaces à produits sclalaires, espaces normés, espaces métriques, espaces topologiques}

    \begin{df}{}{}
        Un produit scalaire sur un $\K$-\ev($\K=\R$ ou $\C$) est une application \[\funto{\langle\cdot , \cdot \rangle}{E^2}{\K}\]
        qui vérifie: \begin{itemize}
            \item Symétrie: $\forall x,y\in E, \overline{\sca{y,x}}=\sca{x,y}$
            \item Homogénéité:$\forall x,y,z\in E,\forall \lambda\in \K, \sca{x+\lambda y,z}=\sca{x,z}+\lambda\sca{y,z}$
            \item Définie positive: $\forall x\in E, \sca{x,x}\ge 0$ et $\sca{x,x}=0\iff x=0$
        \end{itemize}
        On dit que $(E,\sca{\cdot,\cdot})$ est un espace à produit scalaire.\\
        On dit que $\fun{(x,y)}{\sca{x,y}}$ est une forme bilinéaire lorsque $\K=\R$ et une forme sesquilinéaire lorsque $\K=\C$
    \end{df}

    \begin{ex}{}{}
        \begin{enumerate}
            \item sur $\K^n, \sca{x,y}=x_1 y_1+\cdots +x_n y_n$ où $x=(x_1,\cdots,x_n),y=(y_1,\cdots,y_n)$
            \item Sur $\mathcal{C}([0,1],\K)=\{\funto{f}{[0,1]}{\K},f \text{ continue} \}$ muni de $\sca{f,g}=\int_{0}^{1} f(x)\overline{g(x)}\, dx$
        \end{enumerate}
    \end{ex}

    \begin{re}{}{}
        Soit $(E,\sca{\cdot,\cdot})$ un espace à produit scalaire. Notons pour $x\in E : \norme{x}=\sqrt{\sca{x,x}}$ la longueur ou la norme de x.
    \end{re}

    \begin{prop}{Propriété de Pythagore et identité du parallélogramme}{}
        \[\forall x,y\in E^2,\norme{x+y}^2=\norme{x}^2+\norme{y}^2+2\Re (\sca{x,y})\]
        C'est une conséquence de la linéarité et de la symétrie et on en déduit l'identité du parallélogramme :
        \[\norme{x+y}^2+\norme{x-y}^2=2\norme{x}^2+2\norme{y}^2\]
        Si $\sca{x,y}=0$, on dit que x et y sont orthogonnaux et on a $\norme{x+y}^2=\norme{x}^2+\norme{y}^2$
    \end{prop}

    \begin{prop}{Inégalité de Cauchy-Schwartz}{}
        \[\forall x,y\in E^2, |\sca{x,y}|^2\le \sca{x,x}\sca{y,y}\]
        \begin{center}
            avec égalité \ssi x et y sont colinéaires
        \end{center}
    \end{prop}

    \begin{df}{}{}
        Une norme sur un $\K$-\ev E est une application $\funto{\norme{\cdot}}{E}{\R^+}$ qui vérifie :
        \begin{itemize}
            \item Séparation: $\forall x\in E, \norme{x}=0$ \ssi $x=0$
            \item Homogénéité: $\forall x \in E,\forall \lambda\in \K, \norme{\lambda x}=|\lambda|\norme{x}$
            \item Sous-additivité: $\forall x,y\in E, \norme{x+y}\le \norme{x}+\norme{y}$
        \end{itemize}
        On dit alors que $(E,\norme{\cdot})$ est un \evn. La distance induite entre x et y est $d(x,y)=\norme{x-y}$
    \end{df}

    \begin{ex}{}{}
        \begin{enumerate}
            \item L'inégalité de Cauchy-Schwartz garantie que la norme issue du produit scalaire est sous-additive et est bien une norme
            \item Il existe des normes qui ne sont pas issues de produits scalaires. Par exemple $\norme{\cdot}_p$ sur $\K^n$ pour $p\neq 2$ qui ne vérifie pas l'identité du parallélogramme
        \end{enumerate}
    \end{ex}

    On rennonce maintenant à la structure d'\ev.

    \begin{df}{}{}
        Une distance ou métrique sur un ensemble X  est une application $\funto{d}{X^2}{\R^+}$ qui vérifie :
        \begin{itemize}
            \item Séparation: $\forall x,y \in X, d(x,y)=0 \iff x=y$
            \item Symétrie: $\forall x,y\in X, d(x,y)=d(y,x)$
            \item Inégalité triangulaire : $\forall x,y,z \in X, d(x,z)\le d(x,y)+d(y,z)$
        \end{itemize}
        On dit alors que $(X,d)$ est un espace métrique.
    \end{df}

    \begin{ex}{}{}
        \begin{enumerate}
            \item Si $(E,\norme{\cdot})$ est un \evn, $d(x,y)=\norme{x-y}$ est une métrique sur $E$.
            \item Il existe des espaces métriques qui ne sont pas des \evs. Par exemple, 
            \[S^1=\{z_\theta=(\cos(\theta),\sin(\theta)),\theta\in [0,2\pi[\}\] 
            pour la distance d'arc:
            \[ d(z_\theta,z_\alpha)=\min \{|\alpha-\theta|,2\pi-|\alpha-\theta|\}\in [0,\pi]\] 
        \end{enumerate}
    \end{ex}

    \begin{df}{}{}
        Si $(X,d)$ est un espace métrique, pour $x\in X$ et $r>0$, la boule ouverte de centre x et de rayon r est définie par:
        \[B(x,r)=\{y\in X, d(x,y)<r\}\]
        On dira que $O\subset X$ est ouvert si $O=\emptyset$ ou $O$ est réunion (finie ou infinie) de boules ouvertes. En particulier, une boule ouverte est ouverte.
    \end{df}

    \begin{re}{}{}
        Une intersection finie de boules ouvertes est toujours ouverte.
    \end{re}

    On renonce à la notion de métrique pour ne garder que la notion d'ouvert.

    \begin{df}{}{}
        Soit $X$ un ensemble. On appelle topologie sur $X$ la donnée d'un sous-ensemble $\Or \in \Par(X)$ vérifiant:
        \begin{itemize}
            \item $\emptyset\in \Or, X\in \Or$
            \item Stabilité par réunion quelconque: pour toute famille $(O_i)_{i\in I}$ d'éléments de $\Or$, $\bigcup_{i\in I} O_i\in \Or$
            \item Stabilité par intersection finie: Pour toute famille finie $(O_n)_{1\le n\le N}$ d'éléments de $\Or$, $\bigcap_{n=1}^{N}O_n\in \Or$
        \end{itemize}
        On dit alors que $(X,\Or)$ est un espace topologique.\\
        Les éléments de $\Or$ sont appelés les ouverts et le complémentaire d'un ouvert dans $X$ est un fermé.
    \end{df}

    \begin{re}{}{}
        La condition de stabilité par intersection finie peut être remplacée par: Si $O_1,O_2\in \Or$, alors $O_1\cap O_2\in \Or$.
    \end{re}

    \begin{ex}{}{}
        \begin{enumerate}
            \item Si $(X,d)$ est un espace métrique, la topologie métrique est celle pour laquelle $O\subset X$ est ouvert \ssi $O=\emptyset$ ou $O$ est réunion de boules ouvertes. Cette condition est équivalente à:
            \[\forall x\in O,\exists \epsilon_x>0, B(x,\epsilon_x)\subset O\]
            \item $\Or=\{\emptyset,X\}$ est appelée la topologie grossière.
            \item $\Or=\Par(X)$ est appelée la topologie discrète. Notons que c'est la topologie métrique pour la distance donnée par $d(x,y)=1\iff x\neq y$
        \end{enumerate}
    \end{ex}
    
    \begin{re}{}{}
        Soit $(X,\Or)$ un espace topologique $F\subset X$ est fermé \ssi $X\wt F\in \Or$.\\
        On en déduit, en passant au complémentaire, qu'une réunion finie de fermés est fermée, et qu'une intersection quelconque de fermés est un fermé de $X$. Enfin $\emptyset$ et $X$ sont fermés.
    \end{re}

    \begin{prop}{}{}
        Soit $(X,\Or)$ un espace topologique et $A\subset X$. Notons $\Or_A=\{O\cap A, O\in \Or\}\subset \Par(A)$.\\
        On définit une topologie sur A dite topologie induite par $(X,\Or)$ sur A et $(A,\Or_A)$ est un sous-espace topologique de $(X,\Or)$
    \end{prop}

    \begin{proof}
        $\emptyset=\emptyset\cap A$ et $A=A\cap X$ sont dans $\Or_A$. Pour une famille $(O_i)_{i\in I}$ d'éléments de $\Or$: $\bigcup_{i\in I} (O_i\cap A)=(\bigcup_{i\in I} O_i)\cap A \in \Or_A$
    \end{proof}

    \begin{dfprop}{}{}
        Soient $(X,\Or_X)$ et $(Y,\Or_Y)$ deux espaces topologiques. Une application $\funto{f}{X}{Y}$ est continue si elle vérifie l'une des conditions équivalentes suivantes:
        \begin{enumerate}
            \item L'image réciproque par $f$ de tout ouvert de $Y$ est un ouvert de $X$: $\forall O\in \Or, f^{-1}(O)\in\Or_X$
            \item L'image réciproque par $f$ de tout fermé de $Y$ est un fermé de $X$.
        \end{enumerate}
    \end{dfprop}

    \begin{proof}
        \begin{align*}
            A\subset Y, f^{-1}(Y\wt A)&=\{x\in X, f(x)\in Y\wt A\}\\
            &=\{x\in X, f(x)\notin A \}\\
            &=X\wt f^{-1}(A)
        \end{align*}
    \end{proof}

    \begin{ex}{}{}
        Si $(X,d)$ est un espace métrique, la définition de continuité coïncide avec la définition de continuité entre espace métrique.
    \end{ex}

    \begin{re}{}{}
        \begin{itemize}
            \item Si $\Or$ et $\Tr$ sont deux topologies sur $X$, on dit que $\Or$ est moins fine que $\Tr$ si $\Or\subset \Tr$
            \item La topologie grossière est la topologie la moins fine sur $X$
            \item La topologie discrète est la topologie la plus fine sur $X$
        \end{itemize}
    \end{re}

    \begin{ex}{}{}
        Soit $(X,\Or)$ un espace topologique et $A$ une partie de $X$. La topologie induite sur A $\Or_A=\{O\cap A,\forall O\in\Or\}$ est la topologie la moins fine rendant l'injection $A\hookrightarrow X$ continue
    \end{ex}

    \begin{prop}{}{}
        \begin{enumerate}
            \item Soient $(X,\Or_X)$ et $(Y,\Or_Y)$ deux espaces topologiques et $\funto{f}{X}{Y}$ une application:
            \begin{enumerate}
                \item Si $(X,\Or_X)$ est un espace direct, $f$ est toujours continue
                \item Si $\Or_Y=\{\emptyset,Y\},f$ est toujours continue
            \end{enumerate}
            \item La composée de deux applications continues est continue
        \end{enumerate}
    \end{prop}

    \begin{proof}
        \begin{enumerate}
            \item \begin{enumerate}
                \item Si $O$ est un ouvert de $Y$, alors $f^{-1}(O)\in\Par(X)=\Or_X$
                \item $\Or_Y=\{\emptyset,Y\}$ et $f^{-1}(\emptyset)=\emptyset\in\Or_X$ et $f^{-1}(Y)=X\in\Or_X$. Donc $f$ est continue
            \end{enumerate}
            \item Soient $(X,\Or_X),(Y,\Or_Y)$ et $(Z,\Or_Z)$ 3 espaces topologiques et soient $\funto{f}{X}{Y}$ et $\funto{g}{Y}{Z}$ 2 applications continues. Soit $O\in\Or_Z$. $g$ est continue, donc $g^{-1}(O)\in\Or_Y$ et donc $f^{-1}(g^{-1}(O))\in\Or_X$, donc $g\circ f$ est bien continue
        \end{enumerate}
    \end{proof}

    \begin{re}{}{}
        Si $A\subset X$ muni de la topologie induite par $\Or_X$ et $\funto{f}{X}{Y}$ est continue, alors $\funto{f_{|A}}{A}{Y}$ est continue.\\
        De façon analogue, $f(X)\subset B\subset Y$, notons $\fundfto{g}{X}{B}{x}{f(x)}$ alors $g$ est continue $\iff f$ est continue
    \end{re}

    \begin{df}{}{}
        Soient $(X,\Or_X),(Y,\Or_Y)$ deux espaces topologiques. Une application $\funto{f}{X}{Y}$ est un homéomorphisme si $f$ est continue, bijective et $f^{-1}$ est continue.\\
        Deux espaces sont homéomorphes s'il existe un homéomorphisme entre eux
    \end{df}

    \begin{re}{}{}
        Dire que $\funto{f}{X}{Y}$ est un homéomorphisme c'est dire que $f$ est bijective et que cette bijection induit, au niveau des sous-ensembles, une bijection de $\Or_X$ vers $\Or_Y$
    \end{re}

    \begin{ex}{}{}
        \begin{enumerate}
            \item Pour $d\in\N^*$, la boule unité $B(0,1)=\{x\in\R^d,\norme{x}<1\}$ est homéomorphe à $\R^d$ par l'application $\mfun{x}{\frac{x}{1-\norme{x}}}$
            \item On regarde les sous espaces de $\R$: $X=[0,1]$ et $[0,1[\cup \{2\}$.\\
            Soit $\begin{array}{ccccc}
                f : & & Y & \to & X \\
                & & x\in[0,1[ & \mapsto & x \\
                & &  2  & \mapsto & 1
\end{array}$ $f$ continue, bijective mais $f^{-1}$ n'est pas continue car $f(\{2\})=(f^{-1})^{-1}(\{2\})=\{1\}$ alors que $\{2\}$ est un ouvert de $Y$ et $\{1\}$ n'est pas un ouvert de $X$
        \end{enumerate}
    \end{ex}

    \subsection{Adhérence, intérieur et frontière}

    Dans toute la section, $(X,\Or_X)$ est un espace topologique et $A$ est une partie de $X$.

    \begin{df}{}{}
        \begin{enumerate}
            \item La fermeture (ou adhérence) de $A$, noté $\overline{A}$ est le plus petit fermé contenant $A$:
            \[\overline{A}=\bigcap_{\substack{F\text{ fermé}\\ A\subset F}}F=\{x\in X,\forall O\in\Or,x\in O\implies O\cap A\neq \emptyset\}\]
            \item L'intérieur de $A$, noté $\mathring{A}$ est le plus grand ouvert de $X$ contenu dans $A$:
            \[\mathring{A}=\bigcup_{\substack{O\subset A\\ O\text{ ouvert}}}O=\{x\in X,\exists O\in \Or,x\in O\subset A\}\]
            \item La frontière de $A$, noté $\partial A$, est donnée par:
            \[\partial A=\overline{A}\wt \mathring{A}\qquad \text{Elle est fermée}\]
            \item On dit que $A$ est dense dans $X$ si $\overline{A}=X$
        \end{enumerate}
    \end{df}

    \begin{proof}{Démonstration de l'égalité}{}
        Pour la 1, $F$ est un fermé de $X$ signifie qu'il existe $O\in\Or$ tel que $F=X\wt O=O^C$.\\
        Dire de plus que $F$ contient $A$ revient à dire que $A\cap O=\emptyset$ donc
        \[\bigcap_{\substack{F\text{ fermé}\\ A\subset F}}F=\bigcap_{\substack{O\in\Or\\ A\cap O=\emptyset}}O^C=X\wt \bigcup_{\substack{O\in\Or\\ A\cap O=\emptyset}}O=\{x\in X,\forall O\in \Or,x\in O\implies O\cap A\neq \emptyset\}\]
    \end{proof}

    \begin{re}{}{}
        En particulier:
        \[A\text{ est ouvert}\iff \mathring{A}=A\qquad A\text{ est fermé}\iff \overline{A}=A\]
    \end{re}

    \begin{prop}{}{}
        \begin{enumerate}
            \item L'adhérence du complémentaire de $A$ est le complémentaire de l'intérieur de $A$:
            \[\overline{X\wt A}=\overline{A^C}=X\wt\mathring{A}=(\mathring{A})^C\]
            \item L'intérieur du complémentaire de $A$ est le complémentaire de l'adhérence:
            \[\mathring{A^C}=(\overline{A})^C\]
        \end{enumerate}
    \end{prop}

    \begin{proof}{}{}
        Soit $A \subset X$,
        \begin{align*}
            \overline{X\wt A}&=\bigcap_{\substack{X\wt A\subset F\\ F\text{ fermé}}}F\\
            &=\left(\bigcup_{\substack{X\wt A\subset F\\ F\text{ fermé}}}X\wt F\right)^C\\
            &=\left(\bigcup_{\substack{A\subset X\wt F\\ F\text{ fermé}}}X\wt F\right)^C\\
            &=\left(\bigcup_{\substack{A\subset O\\ O\text{ ouvert}}}O\right)^C\\
            &=\left(\mathring{A}\right)^C
        \end{align*}
        Même manipulation pour la deuxième égalité
    \end{proof}

    \begin{cor}{}{}
        Pour que $A$ soit dense dans $X$, il faut et il suffit que son intersection avec tout ouvert non-vide de $X$ soit non-vide.
    \end{cor}

    \begin{proof}{}{}
        Découle de l'égalité de la définition de l'adhérence.\\
        Autre démonstration: Posons $B=X\wt A=A^C$\\
        \[\overline{A}=X\iff\mathring{B}=\emptyset\iff\Or\cap\Par(B)=\{\emptyset\}\iff\forall O\in\Or\wt\{\emptyset\},O\not\subset B=X\wt A\iff\forall O\in\Or\wt\{\emptyset\},O\cap A\neq \emptyset\]
    \end{proof}

    \begin{ex}{}{}
        On munit $\R$ de la topologie usuelle, $\Q$ est dense dans $\R$
    \end{ex}

    \subsection{Voisinages, convergence de suites et continuité de fonctions}

    Dans toute la section, $(X,\Or_X)$ est un espace topologique

    \subsubsection*{Voisinages}

    \begin{df}{}{}
        Soit $x\in X$. Une partie $X$ de $X$ est un voisinage de $x$, s'il existe un ouvert $O$ tel que $x\in O\subset V$. On note $\mathcal{V}(x)$ l'ensemble des voisinages de $x$. On parle de voisinage de $x$ pour désigner un ouvert qui contient $x$.
    \end{df}

    \begin{re}{}{}
        Une partie de $X$ est ouverte \ssi elle est voisinage de tous ses points
    \end{re}

    \begin{proof}{}{}
        Si $O$ est ouverte, $\forall x\in O,x\in O\subset O$, donc $O$ est un voisinage.\\
        ment, si $V$ est un voisinage de chacun de ses points, alors $\forall x\in V,\exists O_x\in\Or$ tel que $x\in O_x\subset V$. Donc $V=\bigcup_{x\in V}O_x$ est un ouvert.
    \end{proof}

    \begin{prop}{}{}
        Soit $x\in X$:
        \begin{enumerate}
            \item $X$ est un voisinage de $x$
            \item $\forall V\in\Vr(x),x\in V$
            \item Si $V\in\Vr(x)$, toute partie de $X$ qui contient $V$ est un voisinage de $x$
            \item L'intersection de deux voisinages de $x$ est un voisinage de $x$
        \end{enumerate}
    \end{prop}

    \begin{proof}{}{}
        Exercice
    \end{proof}

    \begin{prop}{}{}
        $x\in X$, $A\subset X$
        \begin{enumerate}
            \item $A\in \Vr(x)\iff x\in \mathring{A}$
            \item $x\in \overline{A}\iff \forall V\in\Vr(x),V\cap A\neq \emptyset$
        \end{enumerate}
    \end{prop}

    \begin{proof}{}{}
        On a les équivalences:
        \begin{enumerate}
            \item \[x\in\mathring{A}\iff\exists O\in\Or,x\in O\text{ et }O\subset A\iff A\in \Vr(x)\]
            \item \begin{align*}
                x\in \overline{A}&\iff \forall O\in\Or,(x\in O\implies O\cap A\neq \emptyset)\\
                &\forall O\in\Or,\forall V\in \Par(X),(x\in O\subset V\implies V\cap A\neq \emptyset) \\
                &\forall V\in\Vr(x),V\cap A\neq \emptyset \\
            \end{align*}
        \end{enumerate}
    \end{proof}

    \begin{ex}{}{}
        On considère $\R$ munit de la topologie usuelle et $A\subset \R,A\neq\emptyset$ et $A$ majorée. La borne supérieure de $A$ est adhérente à $A$. (Exercice)
    \end{ex}

    \subsubsection*{Convergence de suite}

    \begin{df}{}{}
        Une suite $(x_n)_{n\in\N}$ d'éléments de $X$ converge vers $x\in X$(pour la topologie $\Or$), lorsque:$\forall V\in\Vr(x),\exists n_V\in\R,\forall n\ge n_V,x_n\in V$.\\
        On note $x_n\longrightarrow x$
    \end{df}

    \begin{prop}{}{}
        Soit $A\subset X$ et $(x_n)$ une suite d'éléments de $A$. Si $(x_n)$ converge x, alors $x\in\overline{A}$
    \end{prop}

    \begin{proof}{}{}
        Si $x_n$ converge vers $x$. $\forall V\in \Vr(x),\exists n_V\in\N:\forall n\ge n_V,x_n\in V$. En particulier, $x_{n_V}\in V\cap A$ et $V\in A\neq \emptyset$. Donc $x\in\overline{A}$
    \end{proof}

    \begin{ex}{}{}
        \begin{enumerate}
            \item Pour $X\neq\emptyset$ muni de la topologie grossière, alors pour $x\in X,\Vr(x)=\{X\}$. Donc toute suite $(x_n)$ converge vers tout point de $X$.\\
            Pour $X\neq \emptyset$ muni de la topologie discrète, alors pour $x\in X,\{x\}\in\Vr(x)$. Donc une suite $(x_n)$ converge vers $x$ \ssi elle est stationnaire, $\exists N\in\N,\forall n\ge N,x_n=x$
            \item Si $X$ est munit de la topologie cofinie (Voir TD1-Exo7), si une suite ne prend qu'une et une seule fois chaque point de $X$ pour valeur, alors elle converge vers tout point de $X$.(Exercice)
        \end{enumerate}
    \end{ex}

    \begin{tm}{Convergence dans les espaces métriques}{}
        Soit $(X,d)$ un espace métrique.
        \begin{enumerate}
            \item Une suite $(x_n)$ d'éléments de $X$ converge vers $x\in X$ \ssi $d(x_n,x)\to 0$. Il y a unicité de la limite.
            \item Si $A\subset X$ et $x\in \overline{A}$, alors il existe une suite d'éléments de $A$ convergent vers $x$.
            \item Une partie $A$ de $X$ est fermée \ssi toute suite d'éléments de $A$ convergente a sa limite dans $A$.
        \end{enumerate}
    \end{tm}

    \begin{proof}{}{}
        Soit $\epsilon>0$, notons que $\cho{y\in B(x,\epsilon)\iff d(x,y)<\epsilon\\ B(x,\epsilon)\text{ est un voisinage de $x$}\\ \forall V\in\Vr(x),\exists\epsilon_{x,V},B(x,\epsilon_{x,V})\subset V}$
        \begin{enumerate}
            \item Exercice
            \item Soit $x\in \overline{A}$, alors $\forall n\in \N^*,B(x,\frac{1}{n})\cap A\neq \emptyset$. Choisissons pour chaque $n\in\N^*$, un point $x_n\in B(x,\frac{1}{n})\cap A$. La suite $(x_n)$ converge vers $x$.
        \end{enumerate}
    \end{proof}

    \subsubsection*{Continuité en un point, séquentielle et globale}

    \begin{df}{}{}
        Soient $X,Y$ deux espaces topologiques et $\funto{f}{X}{Y}$ une application.
        \begin{enumerate}
            \item La fonction $f$ est continue en $x\in X$ lorsque $\forall V\in\Vr(f(x)),\exists W\in \Vr(x),f(W)\subset V$
            \item La fonction est séquentiellement continue en $x$ lorsque pour toute suite $(x_n)$ de $X$ telle que $x_n\to x$, on a $f(x_n)\to f(x)$
        \end{enumerate}
    \end{df}

    \begin{re}{}{}
        La continuité (séquentielle) en un point est une propriété locale: Si $A\in\Vr(x)$ et $B\in\Vr(f(x))$ tel que $f(A)\subset B,\funto{f_{|A,B}}{A}{B}$ est (séquentiellement) continue en $x$ \ssi $f$ est (séquentiellement) continue en $x$.
    \end{re}

    \begin{tm}{}{}
        Si $\funto{f}{X}{Y}$ est continue en $x\in X$ alors $f$ est séquentiellement contine en $x$.
    \end{tm}

    \begin{proof}{}{}
        Soit $(x_n)$ une suite de $X$ telle que $x_n\to x$:$\forall U\in \Vr(x),\exists n_V\in\N,\forall n\ge n_V,x_n\in U$. Soit $V\in\Vr (f(x))$.\\
        Puisque $f$  est continue en $x$, il existe $W\in\Vr(x)$ telle que $f(W)\subset V$. Soit $n_W\in\N$:$\forall n\ge n_W,x_n\in W$. En particulier $f(x_n)\in f(W)\subset V,\forall n\ge n_W$. Donc $f(x_n)\to f(x)$.
    \end{proof}

    \begin{re}{}{}
        La réciproque n'est pas vraie en général
    \end{re}

    \begin{ex}{}{}
        On munit $X=\R$ de la topologie co-dénombrable dont les fermés sont $\R$ et les parties (au plus) dénombrables de $\R$.
        \begin{enumerate}
            \item On montre que les suites convergentes sont stationnaires
            \item On en déduit que toute application $\fun{X}{X}$ est séquentiellement continue en tout $x\in X$.
            \item On vérifie que $\funto{f}{X}{X}$ n'est nulle part continue: $f(x)=\cho{1\text{ si }x\ge 0\\ -1\text{ si }x<0}$
        \end{enumerate}
    \end{ex}

    \begin{tm}{Continuité et continuité séquentielle dans un espace métrique}{}
        Si la topologie de $X$ provient d'une métrique alors $\funto{f}{X}{Y}$ est continue en $x$ \ssi elle est séquentiellement continue en $x$.
    \end{tm}

    \begin{proof}{}{}
        Montrons par contraposée que séquentiellement continue $\implies$ continue.\\
        Supposons que $f$ n'est pas continue en $x$:
        \[\exists V\in\Vr(f(x)):\forall W\in\Vr(x),f(W)\not\subset V\]
        En particulier, $f^{-1}(V)$ n'est pas un voisinage de $x$, $x\not\in f^{-1}(V)$, on a alors $x\in\overline{A}=X\wt f^{-1}(V)$. Donc d'après le théorème sur la convergence dans les espaces métriques, il existe $(y_n)$ une suite de $A$ telle que $y_n\to x$. Mais $\forall n,f(y_n)\not\subset V$, donc la suite $(f(y_n))$ ne peut pas admettre $f(x)$ pour limite.
    \end{proof}

    \begin{re}{}{}
        Le point essentiel ici est qu'un ensemble $F$ de $(X,d)$ métrique est fermé \ssi $F=\{x,\exists(x_n)\in F^{\N},x_n\to x\}$. On parle alors d'espaces séquentiels.
    \end{re}

    \begin{tm}{Continuité en tout point et continuité}{}
        La fonction $\funto{f}{X}{Y}$ est continue \ssi elle est continue en tout point de $X$.
    \end{tm}

    \begin{proof}{}{}
        Supposons $f$ continue et soit $V$ un voisinage ouvert de $f(x)$. Alors $W=f^{-1}(V)$ est un ouvert contenant $x$, c'est donc un voisinage de $x$ qui vérifie $f(W)\subset V$.\\
        ment, si $f$ est continue en tout point de $X$. Soit $O_Y$ un ouvert de $Y$ et $x\in f^{-1}(O_Y)$. Puisque $f$ est continue en $x$ et que $O_Y$ est un voisinage de $f(x)$, il existe un voisinage $W$ de $x$ tel que $f(W)\subset O_Y$.\\
        Donc $W\subset f^{-1}(O_Y)$. On en déduit que $f^{-1}(O_Y)$ est un voisinage de chacun de ses points, c'est donc un ouvert, donc $f$ est continue.
    \end{proof}

    \begin{prop}{}{}
        Soit $X$ et $Y$ deux espaces topologiques et $\funto{f}{X}{Y}$ une application.
        \begin{enumerate}
            \item Soit $A\subset X$ et $x\in\overline{A}$. Si $f$ est continue en $x$, $f(x)\in f(\overline{A})$
            \item Si $f$ est continue pour tout $A\subset X$, on a $f(\overline{A})\subset \overline{f(A)}$
        \end{enumerate}
    \end{prop}

    \begin{proof}{}{}
        \begin{enumerate}
            \item Soit $V$ un voisinage de $f(x)$, alors $f^{-1}(V)$ est un voisinage de $x$ et puisque $x\in\overline{A}$, $f^{-1}\cap A\neq \emptyset$. Donc $V\cap f(A)\neq \emptyset$. Donc $f(x)\in\overline{f(A)}$
            \item Découle de 1.
        \end{enumerate}
    \end{proof}

    \subsubsection*{Limites d'applications}

    \begin{dfprop}{}{}
        Soient $X$ et $Y$ deux espaces topologiques, $A\subset X$ et $\funto{f}{A}{Y}$ une application. Soit $x\in\overline{A}$ et $l\in Y$. On dit que $l$ est limite de $f$ en $x$ si, pour tout voisinage $V$ de $l$ dans $Y$, il existe un voisinage $W$ de $x$ dans $X$ tel que $f(W\cap A)\subset V$.\\
        Dans ce cas, $l\in \overline{f(A)}$
    \end{dfprop}

    \begin{proof}{}{}
        Analogue au 1. précédent, laisser en exercice
    \end{proof}

    \begin{re}{}{}
        Dire que $l$ est limite de $f$ en $x$, c'est dire:
        \begin{itemize}
            \item Pour $x\in A$ et $f(x)=l$, que $f$ est continue en $x$
            \item Pour $x\notin A$, la fonction $\fundfto{g}{A\cup\{x\}}{Y}{y}{\cho{l\text{ si }y=x\\ f(y)\text{ si }y\neq x}}$ est continue en $x$
        \end{itemize}
    \end{re}

    \begin{re}{}{}
        Il n'y a pas unicité de la limite en général. Par exemple, si $Y$ est muni de la topologie grossière, $\forall y\in Y,\forall X,\forall A\subset X,A\neq \emptyset,\forall \funto{f}{A}{Y}$, $\forall x\in \overline{A}$, $y$ est limite de $f$ en $x$.
    \end{re}

    \begin{prop}{}{}
        Soient $X,Y,Z$ des espaces topologiques, $A\subset X,\funto{f}{A}{Y}$ une application, $B\subset Y$ contenant $f(A)$ et $\funto{g}{B}{Z},x\in \overline{A},y\in Y$ et $l\in Z$.\\
        Notons $\funto{f=g\circ f}{A}{Z}$. Si $y$ est limite de $f$ en $x$ et $l$ est limite de $g$ en $y$, alors $l$ est limite de $h$ en $x$.
    \end{prop}

    \begin{re}{}{}
        \begin{itemize}
            \item On a $\overline{f(A)}\subset \overline{B}$ donc on peut parler de limite de $g$ en $y$.
            \item En particulier, si $\funto{g}{Y}{Z}$ est continue et $y$ est limite de $f$ en $x$, alors $g(y)$ est limite de $h$ en $x$. Notamment, si $C\subset A\subset X,x\in \overline{C}\subset \overline{A}$, alors si $f$ admet $l$ pour limite en $x$, la restriction de $f$ à $C$ aussi.
        \end{itemize}
    \end{re}

    \begin{proof}{}{}
        Exercice
    \end{proof}

    \section{Espaces topologiques séparés}

    \begin{df}{}{}
        L'espace topologique $(X,\Or)$ est séparé ou de Hausdorff si pour toute paire de points $x,y$ de $X$ distincts, il existe un voisinage $V$ de $x$ et un voisinage $W$ de $y$ tels que $V\cap W=\emptyset$
    \end{df}

    \begin{ex}{}{}
        $\R$ muni de la topologie usuelle est séparé. Pour $x,y\in\R,x\neq y$, le plus petit des deux est dans $]-\infty,\frac{x+y}{2}[$ et le plus grand dans $]\frac{x+y}{2},+\infty[$
    \end{ex}

    \begin{prop}{}{}
        Dans un espace séparé, les singletons sont fermés. Plus généralement, toute partie finie est fermée.
    \end{prop}

    \begin{proof}{}{}
        On suppose $X$ séparé, $x\in X$, montrons que $\{x\}$ est fermé. Soit $y\in X\wt\{x\}$, il existe un voisinage $V$ de $y$ tel que $X\wt V$ contient un voisinage de $x$. En particulier, $x\in X\wt V$, c'est à dire $V\subset X\wt \{x\}$ et $X\wt \{x\}$ est un voisinage de $y$. Ainsi $X\wt \{x\}$ est un voisinage de chacun de ses points, c'est donc un ouvert, donc $\{x\}$ est fermé.
    \end{proof}

    \begin{prop}{}{}
        Soient $X$ et $Y$ deux espaces topologiques, $A\subset X$, $x\in A$. Si $Y$ est séparé, alors toute application $\funto{f}{A}{Y}$ admet au plus une limite en $x$. Si $f$ admet une limite en $x$, on note $\underset{z\to x}{\lim} f(z)=l$
    \end{prop}

    \begin{re}{}{}
        La proposition ne dit rien sur l'existence d'une limite
    \end{re}

    \begin{proof}{}{}
        Supposons que $\funto{f}{A}{Y}$ admet une limite $l$ en $x$ et soit $y\in Y,y\neq l$. Montrons que $y$ n'est pas limite de $f$ en $x$. Puisque $Y$ est séparé, il existe $V$ voisinage de $l$ et $W$ voisinage de $y$ dans $Y$ tels que $V\cap W=\emptyset$. Puisque $l$ est limite de $f$ en $x$, il existe un voisinage $U$ de $x$ dans $X$ tel que $f(U\cap A)\subset V$. Pour tout voisinage $U'$ de $x$, $U\cap U'$ est un voisinage de $x$ et $x\in A$ donc $U\cap U'\cap A\neq \emptyset$.\\
        En particulier, $\emptyset\neq f(U\cap U'\cap A)\subset f(U'\cap A)\cap V$. Puisque $V\cap W=\emptyset, f(U'\cap A)\not\subset W$.\\
        Donc $y$ n'est pas limite de $f$ en $x$.
    \end{proof}

    \begin{re}{}{}
        Sous les mêmes hypothèses, si $x\in A$ et $f$ admet $l$ pour limite en $x$ alors nécessairement $f(x)=l$: la preuve est analogue
    \end{re}

    \begin{prop}{}{}
        Tout espace métrique est Hausdorff.
    \end{prop}

    \begin{proof}{}{}
        Soit $(X,d)$ un espace métrique et $x,y\in X,x\neq y$. Pour tout $z\in X,d(x,y)\le d(x,z)+d(y,z)$ donc au moins un des deux nombres $d(x,z)$ ou $d(y,z)$ est supérieur ou égal à $\frac{1}{2}d(x,y)$.\\
        Donc $B(x,\frac{1}{2}d(x,y))\cap B(y,\frac{1}{2}d(x,y))=\emptyset$
    \end{proof}

    \begin{re}{}{}
        Ces espaces métriques satisfont une propriété de séparation plus forte dite axiome de séparation des espaces parfaitement normaux de Hausdorff.
    \end{re}

    \begin{re}{}{}
        Dans un espace métrique, les points sont fermés et pour tout fermés $F_0$ et $F_1$ tels que $F_0\cap F_1=\emptyset$, $\funto{f}{X}{[0,1]}$ telle que $f^{-1}(\{O\})$
    \end{re}

    \section{Topologie engendrée, topologie produit, topologie quotient}

    \subsubsection*{Topologie engendrée}

    \begin{dfprop}{}{}
        Soit $X$ un ensemble et $\Sigma\subset\Par(X)$:
        \begin{enumerate}
            \item Il existe une topologie $\Or_\Sigma$ sur $X$ pour laquelle tous les éléments de $\Sigma$ sont des ouverts et qui est moins fine que toute topologie ayant cette propriété.
            \item Notons $\Br=\{\bigcap_{i\in I}X_i,\text{$I$ fini et }X_i\in\Sigma\}\cup\{\emptyset,X\}$. Un sous-ensemble $O$ de $X$ est un ouvert pour la topologie $\Or_\Sigma$ \ssi il est réunion (éventuellement infinie) d'éléments de $\Br$
        \end{enumerate}
    \end{dfprop}

    \begin{proof}{}{}
        Notons $\tau$ l'ensemble des réunions d'éléments de $\Br$. Par construction, $\tau$ contient $\emptyset$ et $X$ et est stable par réunion quelconque. Notons que l'intersection de deux éléments de $\Br$ est encore un élément de $\Br$. Pour deux familles (éventuellement infinies) d'éléments de $\Br$, $(U_j)_{j\in J},(V_k)_{k\in K}$, on a
        \[\left(\bigcup_{j\in J} U_j\right)\cap\left(\bigcup_{k\in K} V_k\right)=\bigcup_{j,k}\left(U_j\cap V_k\right)\in \tau\]
        Donc $\tau$ est stable par intersection finie. Finalement $\tau$ est une topologie sur $X$. Soit $\Or$ une topologie sur $X$ pour laquelle tous les éléments de $\Sigma$ sont des ouverts. Par définition d'une topologie on a $\Br\subset \Or$ puis $\tau\subset \Or$. Donc $\tau$ est moins fine que $\Or$ et $\tau=\Or_{\Sigma}$ 
    \end{proof}

    \begin{re}{}{}
        L'ensemble $\Br$ constitue une base de la topologie $\Or_\Sigma$(à comprendre comme "Famille génératrice")
    \end{re}

    \begin{df}{}{}
        Soit $(X,\Or)$ un espace topologique et $\Br\subset \Or$ une famille d'ouverts. On dit que $\Br$ est une base de la topologie $\Or$ lorsque:
        \[\forall O\in \Or,\exists (B_i)_{i\in I}\subset \Br,O=\bigcup_{i\in I}B_i\]
    \end{df}

    \begin{ex}{}{}
        Pour $(X,d)$ espace métrique, $(\Br(x,\frac{1}{n}))_{\substack{x\in X\\ n\in\N^*}}$ est une base de la topologie métrique sur $X$, pour laquelle tout point possède une base dénombrable de voisinage. On dit alors que $X$ est à base dénombrable de voisinage
    \end{ex}

    \subsubsection*{Topologie engendrée par une famille finie d'applications}

    \begin{df}{Topologie initiale}{}
        Soit $X$ un ensemble, $(Y_i,\Or_i)_{i\in I}$ une famille d'espaces topologiques et $(\funto{f_i}{X}{Y_i})_{i\in I}$ une famille d'applications. La topologie engendrée par la famille $(f_i)_{i\in I}$ est la gendrée par $\Sigma=\{f_i^{-1}(O),i\in I,O\in \Or_i\}$. On l'appelle la topologie initiale de $X$.
    \end{df}

    \begin{ex}{}{}
        Si $(X,\Or)$ est un espace topologique et $A\subset X$ la topologie induite par $\Or$ sur $A$ est la topologie engendrée par l'injection $\funtoinj{i}{A}{X}$
    \end{ex}

    \begin{prop}{}{}
        \begin{enumerate}
            \item La topologie $\Or$ engendrée par la famille des fonctions $\funtoinj{f_i}{X}{Y_i}$ est la topologie la moins fine rendant les $f_i$ continues.
            \item Considérons $X$ muni de la topologie $\Or$: Une application $\funto{g}{Z}{X}$ est continue \ssi $\funto{f_i\circ g}{Z}{Y_i}$ est continue pour tout $i\in I$
        \end{enumerate}
    \end{prop}

    \begin{proof}{}{}
        \begin{enumerate}
            \item Evident
            \item La composée de deux applications continues est continue donc si $\funto{g}{Z}{X}$ est continue alors $\funto{f_i\circ g}{Z}{Y_i}$ l'est aussi pour tout $i\in I$.\\
            Réciproquement, supposons que $\funto{f\circ g}{Z}{Y_i}$ continues pour tout $i\in I$. Soit $U$ un ouvert de $X$, il est de la forme 
            \[U=\bigcup_{k\in K}\bigcap_{j\in J}{A_k}_j,{A_k}_j\in\Sigma, J\text{ finie}\]
            \[g^{-1}(U)=\bigcup_{k\in K}\bigcap_{j\in J}g^{-1}({A_k}_j)\]
            \begin{align*}
                g^{-1}(U)&=\bigcup_{k\in K}\bigcap_{j\in J} g^{-1}({A_k}_j)\\
                &=\bigcup_{k\in K}\bigcap_{j\in J}g^{-1}(f_{kj}^{-1}(O))\text{ car $\funinj{I}{J\times K}$ car $K$ queconque}
            \end{align*}
            Or $\forall (k,j)\in K\times J,({f_k}_j\circ g)^{-1}(O)\in \Or_\Sigma$ par hypothèse donc $g^{-1}(U)\in \Or_\Sigma$ comme intersection finie et union d'éléments de $\Or_\Sigma$
        \end{enumerate}
    \end{proof}

    \subsubsection*{Applications: Topologie produit}

    \begin{df}{}{}
        Soit $(X_i,\Or_i)_{i\in I}$ une famille d'espaces topologies et $X=\prod_{i\in I} X_i$ et $\funto{\pi_i}{X}{X_i}$ la projection canonique. La topologie produit est la topologie engendrée par les $\pi_i$
    \end{df}

    \begin{ex}{}{}
        La topologie usuelle sur $\R^n$ est la topologie produit des topologies usuelles sur chaque facteur $\R$
    \end{ex}

    \begin{prop}{}{}
        Soit $(X_i,\Or_i)_{i\in I}$ une famille d'espaces topologiques, $X=\prod_{i\in I} X_i$ muni de la topologie produit.
        \begin{enumerate}
            \item Une base de la topologie produit est donnée par les $\prod_{i\in I}O_i$ où il existe un sous-ensemble fini $J$ de $I$ tels que $O_i\in\Or_i$ pour $i\in J$ et $O_i=X_i$ pour $i\notin J$
            \item Les projections $\funto{\pi_i}{X}{X_i}$ sont ouvertes, c'est à dire que l'image d'un ouvert de $X$ est ouvert
            \item Si $a=(a_i)_{i\in I}\in \prod_{i\in I}X_i=X$ pour tout $n\in \N$, $a_n=(a_{n,i})\in \prod_{i\in I} X_i=X$, alors $a_n\to a$ \ssi pour tout $i\in I, a_{n,i}\to a_i$ dans $X_i$
            \item $X$ est séparé \ssi tous les $X_i$ sont séparés
        \end{enumerate}
    \end{prop}

    \begin{proof}{}{}
        \begin{enumerate}
            \item Pour $O_j\in\Or_j,U_j=\pi_j^{-1}(\Or_j)=O_j\times\prod_{\substack{i\in I\\ i\neq j}}X_i$. Pour $J\subset I$ fini, $\bigcap_{j\in J}U_j=\prod_{i\in I}O_i$ où $O_i\in \Or_i$ pour $i\in J$ et $O_i=X_i$ pour $i\neq j$
            \item On a $\pi_i(\bigcap_{j\in J}U_j)=O\in\Or_i$. Pour toute famille $(A_k)$ de parties de $X$ et $\funto{f_i}{X}{Y}$ une application. $f(\bigcup A_k)=\bigcup f(A_k)$. Tout ouvert $O$ de $X$ est réunion d'ouverts de la forme $\bigcap_{j\in J} U_j$ ($J$ finie) donc $\pi_i(O)$ est ouvert comme réunion d'ouverts.
            \item Puisque $\pi_i$ est continue, elle est séquentiellement continue donc si $a_n\to a,\pi_i(a_i)=a_{n,i}\to\pi_i(a)=a_i$. ment, supposons que pour tout $i, a_{n,i}\to a_i$ dans $X_i$.\\
            Soit $\Vr_j=\prod O_i$ un élément de la base de topologie produit associé à un $J\subset I$ fini, contenant $a$: $a_i\in O_i\in \Or_i$ pour $i\notin J,O_i=X_i$. On a $\forall j\in J, n_J,n\ge n_J,{a_n}_j\in O_j$.\\
            Si $n_{\Vr_j}=\max{n_j,j\in J}$, alors $a_n\in \Vr_j$ pour tout $n\ge n_{\Vr_j}$ donc $a_n\to a$
            \item Soit $x=(x_i)_{i\in I},y=(y_i)_{i\in I}$ dans $X$, si $x\neq y$, il existe $i\in I$ tel que $x_i\neq y_i$\\
            $\impliedby$ Si $X_i$ est séparé, $\exists U,O\in\Or_i$ tels que $O\cap U=\emptyset$, $x_i\in O$ et $y_i\in U$ alors $\prod_{i\neq j}X_j\times U$ et $\prod_{j\neq i}X_j\times O$ sont des voisinages ouverts respectivement de $x$ et $y$ qui sont disjoints.\\
            $\implies$ Supposons $X$ séparé. Soit $i\in I, a,b\in X_i$ tels que $a\neq b$. Soit $x=(x_i)_{i\in I},y=(y_i)_{i\in I}$ dans $X$ tel que $x_j=y_j$ pour $j\neq i$ et $x_i=a,y_i=b$. Alors $x\neq y$, $X$ est séparé donc il existe deux ouverts $V$ et $W$ de $X$ tel que $V\cap W=\emptyset,x\in V$ et $y\in W$.\\
            $\pi_i(V)$ est un ouvert de $X_i$ tel que $a\in\pi_i(V)$ et $\pi_i(W)$ est un ouvert de $X_i$ tel que $b\in \pi(W)$ car $\pi_i$ est ouverte. Si $\pi_i(V)\cap\pi_i(W)\neq \emptyset$, Soit $c\in\pi(V)\cap \pi_i(W)$, alors $z=(z_k)_{k\in I}$ tel que $z_k=x_k=y_k$ pour $k\neq i$ et $z_i=c$. $z$ est dans $V\cap W$, ce qui contredit le fait que $V\cap W=\emptyset$ donc $X_i$ est séparé
        \end{enumerate}
    \end{proof}

    \begin{prop}{}{}
        Un espace topologique est séparé \ssi sa diagonale $\Delta_X=\{(x,x),x\in X\}$ est fermée dans $X\times X$
    \end{prop}

    \begin{proof}
        On considère $\Delta_X^C=(X\times X)\wt \Delta_X$.
        \begin{align*}
            X\text{ est séparé }&\iff\forall (x,y)\in \Delta_X^C,\exists U,V \text{ deux ouverts de $X$ tel que }U\cap V=\emptyset,(x,y)\in U,V\\
            &\iff\forall(x,y)\in\Delta_x,\exists U,V\text{ ouverts de $X$ tel que }(x,y)\in U\times V\text{ et }U\times V\subset \Delta_X^C\\
            &\iff\Delta_X^C\text{ est ouvert}\\
            &\iff \Delta_X\text{ est fermé}
        \end{align*}
    \end{proof}

    \begin{cor}{}{}
        Soient $f,g$ deux applications continues de $X$ dans $Y$ où $Y$ est séparé:
        \begin{enumerate}
            \item L'ensemble $\{x\in X,f(x)=g(x)\}$ est fermé dans $X$
            \item Le graphe de $f, G_G=\{(x,f(x)),x\in X\}$ est fermé dans $X\times Y$
        \end{enumerate}
    \end{cor}

    \begin{proof}
        \begin{enumerate}
            \item On pose $A=\{x\in X,f(x)=g(x)\}$. $f=g$ \ssi $A=X$.\\
            Si $f\neq g$, $A\neq X$. Soit $x\in A^C,f(x)\neq g(x)$\\
            $\exists O_f\in \Or_Y,\exists O_g\in \Or_Y,f(x)\in O_f,g(x)\in O_g,O_f\cap O_g=0$ ($Y$ séparé).\\
            \[x\in f^{-1}(O_f)\Or_X\text{ et }x\in g^{-1}(O_g)\in \Or_X\text{ et }O_f\cap O_g=\emptyset\text{ ($f$ et $f$ continue)}\]
            donc $x\in f^{-1}(O_f)\cap g^{-1}(O_g)\in \Or_X$. Il faut montrer que $f^{-1}(O_f)\cap g^{-1}(O_g)\cap A=\emptyset$, or
            \begin{align*}
                f^{-1}(O_f)\cap g^{-1}(O_g)\cap A&=\{x\in X,f(x)\in O_f,g(x)\in O_g,f(x)=g(x)\}\\
                &=\{x\in X,f(x)=g(x)\in O_f\cap O_g\}\\
                &=\{x\in X,f(x)=g(x)=\emptyset\}=\emptyset
            \end{align*} 
            Donc $\forall x\in A^C$, il existe un ouvert $O_x$ tel que $x\in O_x$ et $O_x\cap A=\emptyset$ donc $A^C=\bigcup_{x\in X}O_x\in \Or_X$ donc $A$ est fermé.
        \end{enumerate}
    \end{proof}

    \begin{prop}{Produit dénombrable d'espaces métriques}{}
        Soit $(X_i,d_i)_{i\in I}$ une famille dénombrable d'espaces métriques dans la topologie produit sur $X=\prod_{i\in}X_i$ provient d'une métrique ($X$ est métrisable)
    \end{prop}

    \begin{proof}
        Pour $I$ fini, posons $d(x,y)=\max\{d(x_i,y_i),i\in I\}$. C'est une métrique sur $X$ dont la topologie associée coincide avec la topologie produit (Exercice).\\
        Pour $I=\N$, posons $d(x,y)=\underset{i\in I}{\max}(\min d_i(x_i,y_i),\frac{1}{2^i})$. On vérifie que $d$ est une métrique sur $X$ dont la topologie coincide avec la topologie produit (Exercice).\\
        Remarque: Notons $V_n=\min(d_n(x_n,y_n),\frac{1}{2^n})$, la suite $V_n$ de $\R_+^*$ converge vers $O$. Il existe donc $N\in\N$ tel que $d(x,y)=\underset{n\in\llbracket 0,N\rrbracket}{\max}(V_n)$.\\
        Donc pour mesurer $d(x,y)$, on ne considère que ce qu'il se passe sur le produit finie $\prod_{n\in\llbracket 0,N\rrbracket}X_n$
    \end{proof}

    \subsubsection*{Topologie finale}

    \begin{df}{}{}
        Soit $X$ un ensemble, $(Y_i,\Or_i)_{i\in I}$ une famille d'espaces topologiques et $(\funto{f_i}{Y_i}{X})_{i\in I}$ une famille d'applications. L'ensemble des $O\subset X$ tels que $f_i^{-1}(O)$ est ouvert dans $Y_i$, pour tout $i\in I$, constitue une topologie sur $X$ appelée topologie finale. On note $\Or_f$ cette topologie
    \end{df}

    \begin{prop}{}{}
        \begin{enumerate}
            \item $\Or_f$ est la topologie la plus fine qui rend les $f_i$ continues.
            \item Si $X$ est muni de la topologie $\Or_f$, une application $\funto{g}{X}{Z}$ est continue \ssi $\funto{g\circ f_i}{Y_i}{Z}$ est continue $\forall i\in I$ 
        \end{enumerate}
    \end{prop}

    \begin{proof}
        Exercice
    \end{proof}

    \subsubsection*{Topologie quotient}

    \begin{df}{}{}
        Soit $(X,\Or)$ un espace topologique et $\Rr$ une relation d'équivalence sur $X$. La topologie quotient sur $X/\Rr$ est la topologie finale associée à la projection canonique $\fun{X}{X/\Rr}$
    \end{df}

    \begin{ex}{}{}
        $X=\R\times\{1,2\}$ muni de la topologie usuelle dont une base est donnée par $\{]a,b[\times \{1\},(a,b)\in\R,a\le b\}\cup\{]a,b[\times \{2\},(a,b)\in\R,a\le b\}$. Soit $\Rr$ la relation d'équivalence sur $X$ donnée par:
        \[\forall x,y\in \R,(i,j)\in\{1,2\}^2,(x,i)\Rr (u,j)\iff x=y<0\text{ ou }(x=y\ge 0\text{ et }i=j)\]
        Une base de la topologie quotient est donnée par:
        \begin{itemize}
            \item $]a,b[,a\le b<0$
            \item $]a,b[\times \{i\}$ avec $i\in\{1,2\}$ et $0\le a\le b$
            \item $]-c,0[\cup [0,d[\times\{1,2\}$ avec $c>0$ et $d>0$
        \end{itemize}
        Tout voisinage de $(0,1)$ rencontre tout voisinage de $(0,2)$ donc $X/\Rr$ n'est pas séparé
    \end{ex}

    \begin{re}{Séparé d'un espace topologique}{}
        Soit $X$ espace topologique et $\Rr$ la relation d'équivalence sur $X$ donnée par $\forall(x,y)\in X,x\Rr y\iff \forall Y$ séparé et toute application continue $\funto{f}{X}{Y}$, on a $f(x)=f(y)$.\\
        Alors $X/\Rr$ est séparé et pour tout espace topologique séparé $Y$ et toute application continue $\funto{f}{X}{Y}$, il existe une application continue $\funto{g}{X/\Rr}{Y}$ tel que (Voir Diagramme p.24)
    \end{re}

    \chapter{Compacité}

    Dans toute la suite $(X,\Or)$ est un espace topologique

    \section{Espace Quasi-compact, espaces compacts}

    \begin{df}{}{}
        Une famille $A=(A_i)_{i\in I}$ de parties de $X$ est un recouvrement de $X$ si $X=\bigcup_{i\in I}A_i$.\\
        Si $J\subset I$ et $\bigcup_{j\in J} A_j=X$, on dit que $(A_j)_{j\in J}$ est un sous-recouvrement de $A$, si de plus $J$ est fini on parle de sous-recouvrement fini.
    \end{df}

    \begin{df}{}{}
        On dit que $X$ est quasi-compact si de tout recouvrement d'ouverts, on peut extraire un sous-recouvrement fini (Propriété de Borel-Lebesgue).\\
        Une partie $K$ de $X$ est quasi-compact si elle est quasi-compacte pour la topologie induite, pour toute famille $(O_i)_{i\in I}$ d'ouverts de $X$ tel que $K\subset \bigcup_{i\in I}$, il existe $J\subset I$ fini tel que $K\subset \bigcup_{j\in J}O_j$
    \end{df}

    \begin{ex}{}{}
        \begin{enumerate}
            \item Tout ensemble fini est quasi-compact. Un ensemble discret est quasi-compact \ssi il est fini
            \item Soit $(x_n)_{n\in N}$ une suite de $X$ admettant une limite $l$, alors $K=\{x_n,n\in \N\}\cup \{l\}$ est quasi compact. En effet, si $(O_i)_{i\in I}$ est une famille d'ouverts tel que $K\subset \bigcup_{i\in I}O_i$. Soit $i_0\in I$ tel que $l\in O_{i_0}$, puisque $l$ est une limite de $(x_n)_n$, il existe $N\in \N$ tel que si $n\ge N$, $x_n\in O_{i_0}$. Il ne reste plus qu'à recouvrir $\{x_0,\ldots,x_n\}$
        \end{enumerate}
    \end{ex}

    \begin{prop}{}{}
        \begin{enumerate}
            \item L'espace $X$ est quasi-compact \ssi pour toute famille $\Fr$ de fermés de $X$, si $\bigcap_{F\in \Fr}F=\emptyset$, il existe une sous-famille finie $f$ de $\Fr$ telle que $\bigcap_{F\in f}F=\emptyset$
            \item Si $K_1$ et $K_2$ sont deux sous-espaces quasi-compacts de $X$, alors $K_1\cup K_2$ aussi.
            \item Si $X$ est quasi-compact et $(F_n)$ est une suite décroissante (pour l'inclusion) de fermés non-vides de $X$, alors $\bigcap_{n\in \N}F_n \neq\emptyset$
        \end{enumerate}
    \end{prop}

    \begin{proof}
        \begin{enumerate}
            \item On passe au complémentaire
            \item Immédiat
            \item Soit $J\subset \N$ finie, $J\neq \emptyset$ et $N$ le plus grand élément de $J$. Alors puisque la suite $(F_n)$ est décroissante $\bigcap_{j\in J}F_J=F_N\neq \emptyset$. Donc d'après (1), $\bigcap_{n\in J}F_n\neq\emptyset$
        \end{enumerate}
    \end{proof}

    \begin{df}{}{}
        On dit que $X$ est compact lorsqu'il est quasi-compact et séparé
    \end{df}

    \begin{ex}{}{}
        $[0,1]$ est un sous-espace compact de $\R$. (Exercice)
    \end{ex}

    \begin{re}{}{}
        Partout dans le monde sauf en France, quasi-compact=compact. On ne demande pas sur les espaces compacts d'être séparés\\
        Compact et quasi-compact sont synonymes dans les espaces séparés
    \end{re}

    \section{Propriétés fondamentales}

    \begin{prop}{}{}
        \begin{enumerate}
            \item Un sous-espace fermé d'un espace quasi-compact est quasi-compact.
            \item Une partie quasi-compacte d'un espace topologique séparée et fermée.
            \item Dans un espace topologique compact, les parties compactes sont les parties fermées
        \end{enumerate}
    \end{prop}

    \begin{proof}
        \begin{enumerate}
            \item Supposons que $X$ est quasi-compact. Soit $F\subset X$ fermé et $(O_i)_{i\in I}$ une famille d'ouverts de $X$ telle que $F\subset \bigcup_{i\in I}O_i$, alors $X=(X\wt F)\cup(\bigcup_{i\in I}O_i)$ est un sous-recouvrement ouvert de $X$ qui est quasi-compact. Donc il existe $J\subset I$ fini tel que $X=(X\wt F)\cup (\bigcup_{j\in J}O_j)$. Par suite $F\subset\bigcup_{j\in J}O_j$
            \item Supposons $X$ séparé et soit $K\subset X$ quasi-compact. Montrons que $X\wt K$ est ouvert. Soit $x\in X\wt K$, puisque $X$ est séparé.
            \[\forall y\in X,y\neq x,\exists O_y\in \Or,y\in O_y,\exists O_x^y,x\in O_x^y\text{ et } O_j\cap O_x^y=\emptyset\]
            Puisque $K$ est quasi-compact, si on peut extraire du recouvrement $(O_y)_{y\in K}$ de $K$ un sous-recouvrement fini $(O_{y_i})_{i\in\llbracket 1,n\rrbracket}$. Alors $\bigcap_{i\in\llbracket1,n\rrbracket}O_x^{y_i}$ est un voisinage ouvert de $x$ contenu dans $X\wt K$. Donc $X\wt K$ est ouvert.
            \item Est un corollaire de (1) et (2)
        \end{enumerate}
    \end{proof}

    \begin{cor}{}{}
        Soit $X$ un espace topologique séparé, $(K_i)_{i\in I}$ une famille de parties compactes de $X$. L'intersection $\bigcap_{i\in I}K_i$ est compacte
    \end{cor}

    \begin{proof}
        On suppose que $I\neq \emptyset$. Soit $i_0\in I$, $X$ étant séparé, chaque $K_i$ est fermé donc $\bigcap_{i\in I}K_i$ est fermée et $\bigcap_{i\in I}K_i\subset K_{i_0}$ qui est compact. Donc $\bigcap_{i\in I}K_i$ est compact
    \end{proof}

    \begin{cor}{}{}
        Une partie $K\subset \R^n$ est compacte \ssi elle est fermée et bornée.
    \end{cor}

    \begin{proof}
        On recouvre $\R^n$ par les boules $B(0,m)$, avec $m\in \N^*$. Si $K\subset \R^n$ est comapct, on peut trouver un entier $N$ tel que 
        \[K\subset \bigcup_{n\in \llbracket 1,N\rrbracket}=B(0,N)\]
        Donc $K$ est borné. Puisque $\R^n$  est séparé, si $K$ est compact, il est fermé d'après le (2) de la proposition.\\
        ment, si $K$ est borné dans $\R^n$, il existe $N>0$ tel que $K\subset {\left[-N,N\right]}^n$. Or $\left[-N,N\right]$ est compact dans $\R$ et (Voir plus loin) un produit fini d'espaces compacts est compact. Donc si $K$ est de plus fermé, il est compact comme fermé dans le compact ${\left[-N,N\right]}^n$
    \end{proof}

    \begin{tm}{}{}
        On suppose que $X$ et $Y$ sont des espaces topologiques et $\funto{f}{X}{Y}$ une application continue
        \begin{enumerate}
            \item Si $X$ est quasi-compact alors $f(x)$ est quasi compact dans $Y$
            \item Si $f$ est bijective, $X$ est quasi-compact et $Y$ séparé, alors $f$ est un homéomorphisme.
            \item Si $Y=\R$ ($\funto{f}{X}{\R}$ continue) et $X\neq \emptyset$ est quasi-compact, alors $f$ est bornée et atteint ses bornes,
            \[\exists x_0,x_1\in X,\forall x\in X, f(x_0)\subset f(x)\subset f(x_1)\]
        \end{enumerate}
    \end{tm}

    \begin{re}{}{}
        On utilise souvent:\\
        1bis. L'image d'un compact, par une application continue à valeurs dans un espace séparé est un compact
    \end{re}

    \begin{proof}
        \begin{enumerate}
            \item Soit $(U_i)_{i\in I}$, une famille  d'ouverts de $Y$ telle que $f(x)\subset \bigcup_{i\in I}U_i$. Puisque $f$ est continue, pour tout $i\in I$, $f^{-1}(U_i)$ est un ovuert de $X$. Donc $(f^{-1}(U_i))_{i\in I}$ est un recouvrement ouvert de $X$ qui est quasi-compact donc il existe $J\subset I$ fini tel que $X=\bigcup_{j\in J}f^{-1}(U_j)$. Donc $f(X)\subset \bigcup_{j\in J} U_j$
            \item Il suffit de montrer que $f^{-1}$ est continue ou que pour tout $F$ fermé de $X$, $f(F)$ est fermé dans $Y$. Soit $F\subset X$ un fermé, puisque $X$ est quasi-compact, $F$ est quasi-compact et donc $f(F)$ est quasi-compact dans $Y$ qui est séparé donc fermé.
            \item Puisque $f(X)$ est un quasi-compact de $\R$ qui est séparé, $f(X)$ est compact dans $\R$ et donc fermé et borné. De plus $f(X)\neq \emptyset$ puisque $X\neq \emptyset$. Donc $f(X)=\overline{f(X)}$, il contient sa borne sup et sa borne inf
        \end{enumerate}
    \end{proof}

    \begin{tm}{Bolzano-Weierstrass}{}
        Dans un espace topologique compact $X$, toute partie infinie $A$ admet des points d'accumulation.
        \[\exists x\in X,x\in \overline{A\wt\{x\}}\]
    \end{tm}

    \begin{proof}
        Supposons que $A$ ne possède pas de point d'accumulation et montrons que $A$ est finie. Soit $x\in A^C=X\wt A$, alors $x$ n'est pas un point d'accumulation de $A$, donc il existe $V\in\Vr_x,V\cap(A\wt\{x\})=V\cap A=\emptyset$. Donc $V\subset X\wt A$. Donc $A^C$ est ouvert et par suite $A$ est fermé. Donc $A$ est fermée dans $X$ compact, donc $A$ est compact. De plus, $A$ est discret. Si $x\in A,x\notin A\wt \{x\}$, donc il existe $V\in \Vr_x,V\cap A=\{x\}$. Or un ensemble discret est quasi-compact \ssi il est fini. Donc $A$ est fini
    \end{proof}

    \begin{tm}{}{}
        Soit $(X_i)_{i\in \llbracket 1,n\rrbracket}$ une famille finie d'espaces topologiques (quasi)-compacts, alors leur produit $\displaystyle{\prod_{i\in\llbracket 1,n\rrbracket}X_i}$ est (quasi)-compact
    \end{tm}

    \begin{proof}
        Rappellons qu'un produit d'espaces séparés est séparé. Il suffit de montrer le théorème pour le produit $X\times Y$ de deux espaces $X$ et $Y$ quasi-compacts.
    \end{proof}

    \begin{tm}{}{}
        $(X_i)_{i\in I}$ est une famille dénombrable d'espace topologiques (quasi)-compact alors $X=\prod_{i\in I} X_i$ est (quasi)-compact
    \end{tm}

    \begin{proof}
        Soient $X,Y$ deux espaces quasi compacts. Soit $(O_i)_{i\in I}$ un recouvrement ouvert $X\times Y$. Pour tout $(x,y)\in X\times Y$. Il existe $k(x,y)\in I$ tel que $(x,y)\in O_{k(x,y)}$ et des ouverts $U_{x,y}$ de $X$ et $V_{x,y}$ de $Y$ tels que $(x,y)\in U_{x,y}\times V_{x,y}\subset O_{k(x,y)}$. Pour $x$ finie, on obtient un recouvrement $(X_{x,y})_{y\in Y}$ de $Y$. $Y$ étant compact, on extrait un sous-recouvrement fini $(V_{x,y|s})_{s\in\llbracket 1,N(x)\rrbracket}$ de $(V_{x,y})_{y\in Y}$. Alors $\bigcup x=\bigcap_{s\in \llbracket 1,N\rrbracket} U_{x,y|s}$ est un ouvert de $X$ qui contient x et qui vérifie $U_x\times V_{x,y|s}\subset O_{k(x,y)}$ pour tout $s\in\llbracket 1,N\rrbracket$ (Voir diagramme Antoine Pioche) et $\{x\}\times Y\subset U_x\times (\bigcup_{s\in \llbracket 1,N(x)\rrbracket} V_{x,y|s})$\\
        La famille $(U_x)_{x\in X}$ fournit un recouvrement ouvert de $X$ qui quasi-compact est quasi-compact donc on peut en extraire un sous-recouvrement fini $(U_{x,s})_{s\in \llbracket 1,N\rrbracket}$. On obtient alors un sous-recouvrement finis
    \end{proof}

    \begin{re}{}{}
        Il existe un théorème beaucoup plus difficile, le théorème de Tikhonov, qui implique que le produit d'une famille quelconque d'espaces (quasi)-compact est (quasi)-compact.
    \end{re}

    \section{Espaces localement compacts}

    \begin{df}{}{}
        Un espace topologique est dit localement compact s'il est séparé et si chacun de ses points admet un voisinage compact.
    \end{df}

    \begin{ex}{}{}
        \begin{enumerate}
            \item Un espace topologique compact est localement compact 
            \item $\R,\C,\R^n,\C^n$ sont localement compacts car les boules fermées sont compacts
            \item Un espace topologique discret est localement compact
            \item $\Q$ sous-espace topologique de $\R$ n'est pas localement compact (exercice)
        \end{enumerate}
    \end{ex}

    \begin{prop}{}{}
        \begin{enumerate}
            \item Une partie fermée (??? ouverte) d'un espace localement compact est localement compact.
            \item Un produit fini d'espaces localement compacts est localement compact
            \item L'intersection de deux sous-espaces localement compacts d'un espace topologique séparé est localement compact
        \end{enumerate}
    \end{prop}

    \begin{proof}
        (Exercice)
    \end{proof}

    \begin{re}{}{}
        La réunion de deux parties localement compactes n'est pas toujours localement compacte. Idem pour une intersection quelconque ou un produit quelconque.
    \end{re}

    \begin{ex}{}{}
        Dans $\R$, $A=\R\wt\left(\{0\}\cup\{\frac{1}{n},n\in \N^*\}\right)$ et $B=\{0\}$
        $A$ et $B$ sont localement compacts et $A\cup B=\R\wt\left( \{\frac{1}{n},n\in \N^*\}\right)$ ne l'est pas car $0$ n'a pas de voisinage compact dans $A\cup B$
    \end{ex}

    \subsection*{Compactification}

    \begin{tm}{Compactification d'Alexandrov}{}
        Soit $X$ un espace localement compact, il existe un espace topologique compact $\tilde{X}$ possédant un point noté $\infty$ tel que $X$ est homéomorphe à $\tilde{X}\wt\{\infty\}$
    \end{tm}

    \begin{proof}
        Idée de preuve: Notons $\Or_X$ la topologie sur $X$. On considère un élément noté $\infty$ et on pose $\tilde{X}=X\cup\{\infty\}$. Soit $\Or=\Or_X\cup \{\tilde{X}\wt K,K\text{ compact de }X\}\subset \Pr(X)$. On montre alors:
        \begin{enumerate}
            \item $\Or$ est une topologie sur $\tilde{X}$
            \item La topologie induite par $\Or$ sur $X$ est la topologie $\Or_X$
            \item $\tilde{X}$ est séparé: Si $x\in X$, on se donne un voisinage compact $K$ de $x$, $\mathring{K}$ et $\tilde{X}\wt K$ sont deux ouverts de $\tilde{X}$ disjoints, $\mathring{K}$ est un voisinage de $x$ et $\tilde{X}\wt K$ un voisinage de $\infty$
            \item $\tilde{X}$ est quasi-compact. Si $(O_i)_{i\in I}$ est un recouvrement ouvert de $\tilde{X}$. Soit $i_0\in I$ tel que $\infty\in O_{i_0}$ et $K$ compact de $X$ tel que $O_{i_0}=\tilde{X\wt K}$. Alors $(O_i\cap K)_{i\in I}$ donne un recouvrement ouvert de $K$ qui est compact, il existe $\delta\subset I$ fini tel que $K\subset \bigcup_{j\in J}O_j$. Finalement, $\tilde{X}=O_{i_0}\cup(\bigcup_{j\in J}O_j)$
        \end{enumerate}
    \end{proof}

    \chapter{Connexité}

    \begin{dfprop}{}{}
        Un espace topologique $X$ est connexe s'il vérifie l'une des conditions suivantes:
        \begin{enumerate}
            \item Les seules parties de $X$ à la fois ouvertes et fermées sont $X$ et $\emptyset$
            \item Il n'existe par de partition de $X$ en deux ouverts non-vides
            \item Il n'existe pas de partition de $X$ en deux fermés non-vides
            \item Toute application continue de $X$ à valeurs dans $\{0,1\}$ est constante
        \end{enumerate}
    \end{dfprop}

    \begin{proof}
        On a les équivalences suivantes:
        \begin{align*}
            A\in\Pr(X)\wt\{X,\emptyset\}\text{ est ouvert et fermé}&\iff\text{$A$ et $A^C$ sont ouvert et fermés non-vide}\\
            &\iff X=A\sqcup A^C\text{ est une partition de $X$ en deux ouverts non-vides}\\
            &\iff X=A\sqcup A^C\text{ est une partition de $X$ en deux fermés non-vides}\\
            &\iff \funto{f}{X}{\{0,1\}}\text{ telle que }f^{-1}(\{0\})=A\text{ et }f^{-1}(\{1\})=A^C\text{ est continue et non-constante}
        \end{align*}
    \end{proof}

    \begin{ex}{}{}
        \begin{enumerate}
            \item Tout espaces grossier est connexe
            \item Un espace topologique discret est connexe \ssi il est réduit à un point
            \item Tout espace homéomorphe à un espace connexe est connexe
            \item Les parties connexes de $\R$ sont les intervalles (exercice).
        \end{enumerate}
    \end{ex}

    \begin{prop}{}{}
        \begin{enumerate}
            \item Si une partie $A$ d'un espace topologique est connexe (pour la topologie induite), alors $\overline{A}$ aussi
            \item Si $\funto{f}{X}{Y}$ est continue et $X$ est connexe alors $f(X)$ aussi
            \item Si $(A_i)_{i\in I}$ est une famille de parties connexes de $X$ telle que $\bigcap_{i\in I}A_i\neq \emptyset$, alors $\bigcup_{i\in I} A_i$ est connexe
            \item Si $(X_i)_{i\in I}$ est une famille d'espaces topologiques connexes, alors $\prod_{i\in I}X_i$ est connexe (pour la topologie produit)
        \end{enumerate}
    \end{prop}

    \begin{proof}
        \begin{enumerate}
            \item Soit $A$ une partie connexe de $X$ et soit $B$ et $C$ deux parties fermées disjointes de $\overline{A}$ telles que $\overline{A}=B\sqcup C$. Alors $A=(A\cap B)\sqcup(A\cap C)$ est une partition de $A$ en deux fermés et puisque $A$ est connexe, $A\cap B=A$ ou $A\cap C=A$, c'est à dire $A\subset B$ ou $A\subset C$ et donc $\overline{A}\subset \overline{B}=B$ ou $\overline{A}\subset \overline{C}=C$ donc $\overline{A}=B$ ou $\overline{A}=C$
            \item Soit $A$ une partie ouverte et fermée de $f(X)\subset Y$. Puisque $f$ est continue, $f^{-1}(A)$ est une partie ouverte et fermée de $X$ qui est connexe. Donc soit $f^{-1}(A)=\emptyset$, soit $f^{-1}(A)=X$, c'est à dire $A=\emptyset$ ou $A=f(X)$
            \item Soit $A=\bigcup_{i\in I}A_i$, on suppose que $A=(O\cup O')\cap A$ où $O\cap A$ et $O'\cap A$ sont des ouverts disjoints de $A$. Soit $x\in \bigcap_{i\in I}A_i \subset A$, et supposons que $x\in O\cap A$, puisque $A_i$ est connexe et que $x\in A_i$ pour un $i\in I$, on a $A_i\subset O\cap A_i$, d'où $A\subset O\cap A$. On a alors $A=O\cap A$ et $O'\cap A=\emptyset$
            \item Puisque $\funto{\pi_i}{X=\prod_{i\in I}}{X_i}$ la projection canonique est continue si $X$ est connexe, $\pi(X)=X_i$ aussi d'après 2..\\
            $\implies$ Le produit de deux espaces topologiques non-vides connexes est connexe: Soit $Y$ et $Z$ connexe, soit $\funto{f}{Y\times Z}{\{0,1\}}$ continue, soit $(y_0,z_0)$ et $(y_1,z_1)$ dans $Y\times Z$. (Voir schema sur photo)\\
            Pour tout $y\in Y$, l'application $\fundfto{g_y}{Z}{Y\times Z}{z}{(y,z)}$ est continue, donc $\funto{f\circ g_y}{Z}{\{0,1\}}$ est continue et puisque $Z$ est connexe, elle est constante. Pour tout $z\in Z$, $f\circ g_{y_0}=f(y_0,z)=f(y_0,z_0)=f(y_0,z_1)$. De même pour tout $y\in Y$, $f(y,z_1)=f(y_0,z_1)=f(y_1,z_1)$.\\
            Finalement $f$ est nécessairement constante. Donc $Y\times Z$ est connexe. Donc lorsque $I$ est fini, si les $X_i$ sont connexes, $\prod_{i\in I}X_i$ aussi. Soit $(X_i)_{i\in I}$ une famille (pas nécessairement fini) d'espaces connexes et soit $X=\prod_{i\in I}X_i$ et $U,V$ deux ouverts tels que $X=U\cup V$\\
            Il existe $J\subset I$ fini tel que pour tout $i\in I\wt J$, $\pi_i(U)=\pi_i(V)=X_i$. On peut donc trouver $x\in U$ et $y\in V$ telle que $\pi_i(x)=\pi_i(y)=a_i$ pour tout $i\in I\wt J$. Soit alors $\fundfto{h}{\prod_{j\in J}X_j}{X=\prod_{i\in I}X_i}{z}{\cho{a_i \text{ si $i\in I\wt J$}\\ z_j\text{ si }j\in J}}$ $z=(z_j)_{j\in J}$ et en particulier $h(\pi_j(x))=x$ et $h(\pi_y(y))=y$. L'application $h$ est continue, donc $\prod_{j\in J} X_j=h^{-1}(X)=h^{-1}(U)\cup h^{-1}(V)$ est un recouvrement en deux ouverts non-vides du produit fini $\prod_{j\in J} X_j$ qui est connexe. Ceci n'est possible que si $h^{-1}(U)\cap h^{-1}(V)\neq \emptyset$ ou encore $U\cap V\neq \emptyset$. Donc $X$ est connexe. 
        \end{enumerate}
    \end{proof}

    \begin{df}{}{}
        La composante connexe $C_x$ de $x\in X$ est la plus grande partie connexe de $X$ qui contient $x$: $C_x$ est la réunion de toutes les parties connexes de $X$ qui contiennent $x$.
    \end{df}

    \begin{re}{}{}
        Conséquences de la propriété
        \begin{enumerate}
            \item Les composantes connexes sont fermées d'après 1.
            \item Deux composantes connexes sont disjoints ou égales et forment donc une partition de $X$ par des fermés. La relation $y Rx$ lorsque $y\in C_x$ est donc une relation d'équivalence.
            \item Dire que $X$ est connexe, c'est à dire qu'il a une seule composante connexe
        \end{enumerate}
    \end{re}

    \begin{df}{Connexité par arc}{}
        Un espace topologique est dit connexe par arc lorsque pour tout $x_0,x_1$ de $X$, il existe une application continue $\funto{c}{[0,1]}{X}$ telle que $c(0)=x_0$ et $c(1)=x_1$. On dit alors que $c$ est un chemin ou un arc entre $x_0$ et $x_1$
    \end{df}

    \begin{prop}{}{}
        \begin{enumerate}
            \item Si $\funto{f}{X}{Y}$ est continue et $X$ est connexe par arc alors $f(X)$ aussi.
            \item Si $(A_i)_{i\in I}$ est une famille de parties connxes par arc de $X$ telle $\bigcap_{i\in I}A_i\neq \emptyset$ alors $\bigcup_{i\in I}$ est connexe par arc.
            \item Si $(X_i)_{i\in I}$ est une famille d'espaces topologiques connexes par arc, alors $\prod_{i\in I}X_i$ aussi.
        \end{enumerate}
    \end{prop}

    \begin{proof}
        Exercice
    \end{proof}

    \begin{prop}{}{}
        Si $X$ est connexe par arc alors $X$ est connexe.
    \end{prop}

    \begin{proof}
        Soit $x,y\in X$ et $\funto{c}{[0,1]}{X}$ continue tels que $c(0)=x$ et $c(1)=y$. Alors $c([0,1])$ est une partie connexe de $X$ qui contient $x$ et $y$ donc $C_x=C_y$. Ceci étant vrai pour tout $x,y$ dans $X$, $X$ n'a qu'une composante connexe. Donc $X$ est connexe
    \end{proof}

    \begin{re}{}{}
        La réciproque est fausse
    \end{re}

    \begin{ex}{}{}
        On considère dans $\R^2$, $A=\{(x,\sin(\frac{1}{x})),x\in \R_+^*\}$\\
        Puisque $\fundfto{f}{]0,+\infty[}{A}{x}{(x,\sin(\frac{1}{x}))}$ est un homéomorphisme, $A$ est connexe et connexe par arc. On a donc $\overline{A}=A\cup\{0\}\times[-1,1]$ qui est connexe mais $\overline{A}$ n'est pas connexe par arc.
    \end{ex}

    \chapter{Propriétés topologiques des espaces métriques}

    Dans toute la suite $(X,d)$ est un espace métrique, et $X$ est muni de la topologie métrique
    
    \section{Avant-propos: questions de métriques}

    \begin{dfprop}{}{}
        Soit $d_1,d_2$ deux métriques sur $X$ et $\fundfto{\Id_{1,2}}{(X,d_1)}{(X,d_2)}{x}{x}$, $\fundfto{\Id_{2,1}}{(X,d_2)}{(X,d_1)}{x}{x}$
        \begin{enumerate}
            \item $d_1$ et $d_2$ sont topologiquement équivalentes si elles définissent la même topologie, ce qui équivaut à dire:
            \[\forall \epsilon>0,\forall x\in X,\exists\alpha>0,\forall y\in X(d_1(x,y)<\alpha\implies d_2(x,y)<\epsilon)\text{ et }(d_2(x,y)<\alpha\implies d_1(x,y)<\epsilon)\]
            Ce qui équivaut à dire que $\Id_{1,2}$ est un homéomorphisme
            \item $d_1$ et $d_2$ sont métriquement équivalents Si
            \[\forall \epsilon>0,\exists\alpha > 0,\forall x,y\in X(d_1(x,y)<\alpha\implies d_2(x,y)<\epsilon)\text{ et }(d_2(x,y)<\alpha\implies d_1(x,y)<\epsilon)\]
            Ceci équivaut à $\Id_{1,2}$ et $\Id_{2,1}$ sont uniformément continues.
            \item $d_1$ et $d_2$ sont équivalentes s'il existe des constantes $m$ et $M$ dans $\R_+^*$ telles que
            \[\forall x,y\in X,md_1(x,y)\le d_2(x,y)\le Md_1(x,y)\]
            Ceci équivaut à dire que $\Id_{1,2}$ et $\Id_{2,1}$ sont lipschitiennes.
        \end{enumerate}
    \end{dfprop}

    \begin{proof}
        Exercice
    \end{proof}

    \begin{re}{}{}
        On a équivalentes $\implies$ métriquelement équivalentes $\implies$ topologiquement équivalentes
    \end{re}

    \section{Suites de Cauchy et complétude}

    \subsection*{Suites de Cauchy}

    \begin{df}{}{}
        La suite $(x_n)_{n\in \N}$ de $(X,d)$ est de Cauchy lorsque
        \[\forall \epsilon >0,\exists N\in \N,\forall n\ge N,m\ge N, d(x_n,x_m)\le \epsilon\]
    \end{df}

    \begin{prop}{}{}
        \begin{enumerate}
            \item Une suite convergente est de Cauchy
            \item Une suite extraite d'une suite de Cauchy est de Cauchy
            \item Une suite de Cauchy est bornée
            \item Une suite de Cauchy a au plus une valeur d'adhérence: Si deux suites extraites d'une suite de Cauchy convergent, elles convergent vers la même limite.
            \item Une suite de Cauchy converge \ssi elle admet une valeur d'adhérence
            \item L'image d'une suite de Cauchy par une application: uniformément continue et de Cauchy
        \end{enumerate}
    \end{prop}

    \begin{proof}
        Exercice
    \end{proof}

    \begin{re}{}{}
        En particulier, les suites de Cauchy restent de Cauchy si on remplace la métrique $d$ par une métrique $\delta$ métriquement équivalente à $d$.
    \end{re}

    \subsection*{Espaces métriques complets}

    \begin{df}{}{}
        On dit qu'un espace métrique est complet lorsque toute suite de Cauchy converge.
    \end{df}

    \begin{re}{}{}
        Ce n'est pas une notion topologique: un espace peut être complet pour une métrique et non complet pour une métrique topologiquement équivalente.
    \end{re}

    \begin{prop}{}{}
        L'espace métrique $(X,d)$ est complet \ssi pour toute suite décroissante $(F_n)_{\nN}$ de fermés non-vides de $X$ telle que la suite des diamètres $(\delta(F_n))_\nN$ converge vers $0$, on a $\bigcap_\nN F_n=\{x\}$ où $x\in X$
    \end{prop}

    \begin{proof}
        Supposons $(X,d)$ complet, soit $(F_n)_\nN$ une suite décroissante de fermés non-vides telle que $\delta(F_n)\underset{n\to\infty}{\longrightarrow} 0$.\\
        Pour chaque $\nN$, ssoit $x_n\in F_n$. Pour $m\ge n$, $F_m\subset F_n$. Donc $x_n,x_m\in F_n$ et $d(x_n,x_m)\le \delta (F_n)$. Puisque $\delta(F_n)\underset{n\to\infty}{\longrightarrow} 0$, $(x_n)_\nN$ est de Cauchy et elle converge vers $x$. Puisque $(x_k)_{k\ge n}$ est une suite de $F_n$ qui est fermée donc $x\in F_n$ et ceci pour tout $n$. Donc $x\in \bigcap_{\nN}F_n$ donc $\bigcap_\nN F_n\neq \emptyset$. Si $x$ et $y\in \bigcap_{\nN} F_n$ alors, $d(x,y)\le \delta(F_n)$ pour tout $\nN$ donc $d(x,y)=0$ et $x=y$. ment, supposons que l'intersection de toute suite décroissante de fermés non-vides de $X$ est non-vide. Soit $(x_n)_\nN$ une suite de Cauchy de $X$ et pour $\nN$, $F_n=\overline{\{x_k,k\ge n\}}$. Alors $(F_n)_\nN$ est une suite décroissante de fermés. Si $x\in \bigcap_{\nN}$, pour tout $\nN,d(x,x_n)\le\delta(F_n)$. $\delta(F_n)\Longrightarrow 0$ car $(x_n)_n$ est de Cauchy, donc $x_n\longrightarrow x$
    \end{proof}

    \begin{ex}{}{}
        \begin{enumerate}
            \item $\fundfto{h}{]0,+\infty[}{]]0,+\infty[}{x}{\frac{1}{x}}$\\
            Pour $d_1(x,y)=|x-y|$ et $d_2(x,y)=|\frac{1}{x}-\frac{1}{y}|$ sont topologiquement équivalentes mais non-métriquement équivalentes. Car $h$ n'est pas uniformément continue
            \item Sur $]0,+\infty[$, soit $d_3(x,y)=\min(1,|x-y|)$. $d_3$ est topologiquement équivalente à $d_1$ et même métriquement équivalente à $d_1$, mais $d_1$ et $d_2$ ne sont pas équivalentes
        \end{enumerate}
    \end{ex}

    \begin{tm}{}{}
        Tout espace métrique compact est complet
    \end{tm}

    \begin{proof}
        On psait déjà que si $X$ est quasi-compact et $(F_n)$ une suite décroissantede fermés non-vides de $X$ alors $\bigcap_{n\in \N}F_n\neq \emptyset$. La proposition permet de conclure.
    \end{proof}

    \begin{tm}{Théorème de Baire}{}
        Si $X$ est complet alors il est de Baire: Toute intersection dénombrable d'ouverts denses est dense, ou de façon équivalente toute réunion dénombrable de fermés d'intérieur vide est d'interieur vide
    \end{tm}

    \begin{proof}
        Soit $(O_n)_{n\in \N}$ une famille d'ouverts denses de $X$, on doit montrer que si $O$ est un ouvert de $X$, alors $O\cap(\bigcap_{\nN})\neq \emptyset$. Puisque $O_0$ est dense, il existe $x\in O\cap O_0$ qui est ouvert, donc il existe $r_0$ tel que $0<r_0<1$ et $\overline{B}(x_0,r_0)\subset O\cap O_0$. On construit par réccurence une suite $(x_n)_{\nN}$ de $X$. On construit $(r_n)_\nN$ de $]0,1[$ telles que:
        \begin{itemize}
            \item $0<r_n<2^{-n}$
            \item $\overline{B}(x_n,r_n)\subset B(x_{n-1},r_{n-1})\cap O_n$
        \end{itemize}
        On obtient une suite décroissante de fermés non-vides $(\overline{B}(x_n,r_n))_\nN$ dont le diamètre tend vers $0$ de $X$ qui est complet.\\
        Donc $\emptyset\neq \bigcap \overline{B}(x_n,r_n)\subset O\cap(\bigcap_{n\in\N} O_n)$
    \end{proof}

    \subsection*{Complété d'un espace métrique}

    Soit $(X,d)$ un espace métrique, on construit son complété: espace métrique $(\tilde{X},\tilde{d})$ complet qui satisfait:
    \begin{itemize}
        \item $X$ est un sous-espace dense de $\tilde{X}$.
        \item $\forall \funto{f}{X}{Y}$ uniformément continue avec $Y$ complet, il existe une unique $\funto{\tilde{f}}{\tilde{X}}{Y}$ continue telle que $\tilde{f}_{|X}=f$
    \end{itemize}
    De plus $\tilde{f}$ est uniformément continue. L'ensemble $\tilde{X}$ est formé des "limites formelles" de suites de Cauchy:
    \begin{itemize}
        \item Soit $X^{\N,C}\subset X^\N$ l'ensemble des suites de Cauchy de $X$
        \item Soit $\Rr$ la relation d'équivalence sur $X^{\N,C}$ donnée par:
        \[(x_n)\Rr (y_n)\iff d(x_n,y_n)\underset{n\to+\infty}{\longrightarrow}0\]
        \item Soit $\tilde{X}=X^{\N,C}/\Rr$\\
        Pour $(x_n),(y_n)$ dans $X^{\N,C}$, la suite réelle $(d(x_n,y_n))_\nN$ est de Cauchy dans $\R^+$, donc elle converge et on pose $\tilde{d}([(x_n)],[(y_n)])=\underset{n\to+\infty}{\lim}d(x_n,y_n)$. On vérifie que $\tilde{d}$ est une métrique sur $\tilde{X}$.
        \item L'injection qui à $x\in X$ associe la suite constante égale à $x$ induit une isométrie $\funto{\iota}{X}{\tilde{X}}$ qui permet d'identifier $X$ à un sous espace de $\tilde{X}$. De plus pour tout $x=[(x_n)]$ de $\tilde{X}$, où $(x_n)\in X^{\N,C}$ la suite $(\zeta (x_n))_n$ converge dans $\tilde{X}$ vers $a$. Donc $X$ est dense dans $\tilde{X}$.
        \item $(\tilde{X},\tilde{d})$ est complet. Si $(a_n)$ est une suite de Cauchy de $\tilde{X}$, puisque $X$ est dense dans $\tilde{X}$\\
        $\forall \nN,\exists x_n\in X$ tel que $\tilde{d}(\iota(x_n),a_n)<2^{-n}$. Pour $n,m\in \N$,
        \[d(x_n,y_n)=\tilde{d}(\iota(x_n),\iota(x_m))\le \tilde{d}(\iota(x_n),a_n)+\tilde{d}(a_n,a_m)+\tilde{d}(a_m,\iota(x_m))\le 2^{-n}+2^{-n}+\tilde{d}(a_n,a_m)\]
        Donc on déduit que la suite $(x_n)$ de $X$ est de Cauchy. De plus la suite $(\iota(x_n))$ converge vers $a=[(x_n)]\in \tilde{X}$ et $\tilde{d}(a,a_n)\le \tilde{d}(a,\iota(x_n))+\tilde{d}(\iota(x_n),a_n)$. Donc la suite $(a_n)_\nN$ converge vers $a$. On termine en montrant la propriété sur le prolongement des application continues.
    \end{itemize}

    \section{Compacité dans les espaces métriques}

    \subsection*{Caractérisation de la compacité dans les espaces métriques}

    \begin{tm}{}{}
        Soit $(X,d)$ un espace métrique. Les propositions suivantes sont équivalentes:
        \begin{enumerate}
            \item $X$ est compact
            \item Propriété de Bolzano-Weierstrass: Toute suite de $X$ admet une sous-suite convergente.\\
            Lorsqu'elles ont lieu, elles entraînent la propriété suivante
            \item Lemme de Lebesgue: Pour tout recouvrement $(O_i)_{i\in I}$ de $X$ par des ouverts, il existe $r>0$ tel que pour tout $x\in X$, il existe $i_x\in I$, $B(x,r)\subset O_{i_x}$        \end{enumerate}
    \end{tm}

    \begin{proof}
        $(1)\implies (2)$ Soit $X$ compact et $(x_n)_\nN$ une suite de $X$. Si $(x_n)_\nN$ n'admet pas de sous-suite convergente alors:
        \begin{itemize}
            \item $(x_n)_\nN$ prend un nombre infini de valeurs
            \item $\forall x\in X,\exists r_x>0$ tel que $|B(x,r_n)\cap{x_n,\nN}|$ est fini
        \end{itemize}
        On obtient alors un recouvrement ouvert de $X=\bigcup_{x\in X} B(x,r_x)$ dont on peut extraire un sous-recouvrement fini, chacun des éléments de ce recouvrement ne contenant qu'un nombre fini de termes de la suite $(x_n)_\nN$. C'est en contradiction avec $(1)$\\
        $(2)\implies (3)$ Si le lemme de Lebesgue est en défaut:
        \[\forall n\in \N,\exists x_n\in X,\forall i\in I,B(x_n,2^{-n})\not\subset O_i\qquad (*)\]
        On peut extraire une sous-suite convergente $(x_{n_k})_{k\in \N}$ de la suite $(x_n)_\nN$ qui converge vers $x\in X$.\\
        Si $x\in O_{i_x}$, il existe $n,N\in \N$ tel que
        \begin{itemize}
            \item $b(x,2^{-n})\subset O_{i_x}$
            \item $\forall k\ge N,d(x,x_{n_k})< 2^{-n-1}$
        \end{itemize}
        Donc $B(x_N,2^{-n-1})\subset O_{i_x}$ ce qui contredit $(*)$\\
        $(2)\implies (1)$ Soit $(O_i)_{i\in I}$ un recouvrement ouvert de $X$ et $r>0$ tel que
        \[\forall x\in X\exists i_x\in I,X(x,r)\subset O_{i_x}\]
        Supposons qu'il n'existe pas de sous-ensemble fini $J$ de $I$ tel que $\bigcup_{j\in J}O_j=X$. Alors pour toute partie finie $A$ de $X$, comme $\bigcup_{x\in A}B(x,r)$ est contenue dans une réunion finie de $O_i$, il existe $x\in X$ tel que $d(x,a)>r$ pour tout $a\in A$.\\
        On construit alors une suite de $X$ de la façon suivante:
        \begin{itemize}
            \item On choisi $x_0\in X$
            \item Pour tout $n\in\N$, on choisit $x_{n+1}\notin \bigcup_{k\in\llbracket 0,n}B(x_n,r)$
        \end{itemize}
        De sorte que pour tout $p,q\in \N,d(x_q,x_p)>r$. La suite $(x_n)_\nN$ n'admet pas de sous-suite convergente.
    \end{proof}

    \subsection*{Produits dénombrables d'espaces compactes}

    \begin{tm}{Théorème de Tikhinov métrique dénombrable}{}
        Soit $(X_n,d_n)_{n\in \N}$ une famille dénombrable d'espaces métriques compacts alors leur produit $X=\prod_{x\in \N}X_n$ est métrisable et compact.
    \end{tm}

    \begin{proof}
        Nous savons déjà que $d(x,y)=\max_{\nN}\{\min(d_n(x_n,y_n),2^{-n})\}$ est une métrique sur $X$. Il reste à montrer que $X$ est compact.\\
        La preuve repose sur le procédé diagonal de Cantor. Notons pour tout $\nN$, $\funto{\pi_n}{X}{X_n}$ la projection canonique (continue)\\
        Soit $\Xr=(\Xr^n)$ une suite de $X$.\\
        \begin{center}
            \begin{tabular}{c|cccc}
                $n$&$X_0$&$X_1$&$X_2$&$\cdots$\\ \hline
                $0$&$\Xr_0$&$\Xr_1$&$\Xr_2$&$\cdots$\\
                $1$&$\Xr_0^1$&$\Xr_1^1$&$\Xr_2^1$&$\cdots$\\
                $\vdots$&$\vdots$&$\vdots$&$\vdots$&$\ddots$
            \end{tabular}
        \end{center}
        On veut extraire une sous-suite convergente de $\Xr$. Puisque $X_0$ est compact, on peut extraire une sous-suite $\Xr^{(0)}=\left(\Xr^{\varphi_0(c)}\right)_\nN$ de $\Xr$ telle que $\pi_0(X^{(0)})$ converge vers $\Xr_0$.\\
        Puisque $X_1$ est compact, on peut extraire une suite $\Xr^{(1)}=\left( \Xr^{\varphi_0\circ\varphi_1(n)}\right)_\nN$ de $\Xr^{(0)}$ telle que $\pi_1(\Xr^{(1)})$ converge vers $X_1$ et $\pi_0(\Xr^{(1)})$ converge toujours vers $x_0$ comme suite extraite d'une suite convergente.\\
        On procède ainsi par réccurence et on obtient une suite $(\Xr^{(k)})_{k\in \N}$ de suites extraites de $\Xr$ où $\Xr^{(k)}=\left( \Xr^{\varphi_0\circ\varphi_1\circ\cdots\circ \varphi_k(n)}\right)_\nN$ est telle que $\pi_j(\Xr^{(k)})$ converge vers $x_j$ pour $j\in\llbracket 0,k\rrbracket$.\\
        Considérons alors la suite extraite de $\Xr$ suivante $\Xr^{\varphi}=(\Xr^{\varphi(n)})_\nN$ où $\varphi(n)=\varphi_0\circ\varphi_1\circ\cdots\circ \varphi_n(n)$
        \begin{center}
            $\Xr^{\varphi}$
            \begin{tabular}{c|cccc}
                $n$&$X_0$&$X_1$&$X_2$&$\cdots$\\ \hline
                $0$&$\Xr_0^{\varphi_0(0)}$&$\Xr_1^{\varphi_0(0)}$&$\Xr_2^{\varphi_0(0)}$&$\cdots$\\
                $1$&$\Xr_0^{\varphi_0\circ\varphi_1(1)}$&$\Xr_1^{\varphi_0\circ\varphi_1(1)}$&$\Xr_2^{\varphi_0\circ\varphi_1(1)}$&$\cdots$\\
                $2$&$\Xr_0^{\varphi_0\circ\varphi_1\circ \varphi_2(2)}$&$\Xr_1^{\varphi_0\circ\varphi_1\circ \varphi_2(2)}$&$\Xr_2^{\varphi_0\circ\varphi_1\circ \varphi_2(2)}$&$\cdots$\\
                $\vdots$&$\vdots$&$\vdots$&$\vdots$&$\ddots$
            \end{tabular}
            \qquad
            $\Xr^{(1)}$
            \begin{tabular}{c|cccc}
                $n$&$X_0$&$X_1$&$X_2$&$\cdots$\\ \hline
                $0$&$\Xr_0^{\varphi_0\circ\varphi_1(0)}$&$\Xr_1^{\varphi_0\circ\varphi_1(0)}$&$\Xr_2^{\varphi_0\circ\varphi_1(0)}$&$\cdots$\\
                $1$&$\Xr_0^{\varphi_0\circ\varphi_1(0)}$&$\Xr_1^{\varphi_0\circ\varphi_1(0)}$&$\Xr_2^{\varphi_0\circ\varphi_1(0)}$&$\cdots$\\
                $\vdots$&$\vdots$&$\vdots$&$\vdots$&$\ddots$
            \end{tabular}
        \end{center}
        Pour chaque colonne $n$, on obtient une suite extraite à partir du rang $n$ de la suite $\pi_n(X^{(n)})$ et qui converge donc vers $x_n$
    \end{proof}

    \subsection*{Compacité et continuité uniforme}

    \begin{tm}{Théorème de Heine}{}
        Toute application continue d'un espace métrique compact à valeur dans un espace métrique est uniformément continue
    \end{tm}

    \begin{proof}
        Soit $(X,d_X)$ un espace métrique compact, $(Y,d_Y)$ un espace métrique et $\funto{f}{X}{Y}$ continue.\\
        Soit $\epsilon>0$, pour tout $x\in X$,
        \[O_x=\{z\in X,d_Y(f(x),f(z))<\epsilon\}\]
        Pour tout $y_0\in Y, \funto{d_Y(y_0,\cdot)}{Y}{\R^+}$ continue, $f$ est continue donc pour $x\in X$, la fonction $\fun{z\in X}{d_Y(f(x),f(z))}$ est continue.\\
        Donc $O_x$ est un ouvert de $X$ qui contient $x$. La famille $(O_x)_{x\in X}$ est un recouvrement ouvert de $X$ qui est compact, choisissons $r$ un nombre de Lebesgue de ce recouvrement.\\
        Pour $x,y\in X$ tel que $d(x,y)<r$, il existe $z\in X$ tel que $y\subset B(z,r)\subset O_z$ et donc
        \begin{align*}
            d_Y(f(x),f(y))&\le d_Y(f(x),f(z))+d_Y(f(z),f(y))\\
            &\le 2\epsilon
        \end{align*}
    \end{proof}

    \subsection*{Séparabilité des espaces métriques compacts}

    \begin{prop}{}{}
        Un espace métrique compact est séparable.
    \end{prop}

    \begin{proof}
        Soit $(X,d)$ un espace métrique compact. Pour tout $k>0$, il existe une famille finie $(x_s^k)_{s\in\llbracket1,N_k\rrbracket}$ de points de $X$ telle que
        \[X=\bigcup_{s\in\llbracket 1,N_k\rrbracket} B(x_s^k,\frac{1}{k})\]
        Alors la famille dénombrable $\{x_0^k,k\in \N^*,s\in \llbracket 1,N_k\rrbracket\}$ est dense dans $X$.
    \end{proof}

    \chapter{Espaces de fonctions continues sur un espace compacts}

    \section{Espace $C_b(X,Y)$ et $C(X,Y)$}

    Soit $(X,d_X)$ et $(Y,d_Y)$ deux espaces métriques

    \subsubsection*{Avant-propos}

    On désigne par:
    \begin{itemize}
        \item $C(X,Y)$ l'espace des fonctions continues de $X$ dans $Y$
        \item $B(X,Y)$ l'espace des fonctions bornées de $X$ dans $Y$
        \item $C_b(X,Y)$ l'espace des fonctions continues et bornées de $X$ dans $Y$
    \end{itemize}

    \begin{re}{}{}
        \begin{itemize}
            \item Si $X$ est compact $C_b(X,Y)=C(X,Y)$.
            \item Ces espaces $B(X,Y)$ et $C_b(X,Y)$ sont des espaces métriques pour 
            \[d_\infty(f,g)=\sup_{x\in X}d_Y(f(x),g(x))\]
            \item Si $d_\infty(f_n,f)\underset{n\to +\infty}{\longrightarrow}0$, on dit que la suite $(f_n)_\nN$ converge uniformément vers $f$.
        \end{itemize}
    \end{re}

    \begin{prop}{}{}
        \begin{enumerate}
            \item La limite uniforme d'une suite de fonctions continues est continue:
            \[C_b(X,Y)\text{ est fermé dans $B(X,Y)$}\]
            \item Si $Y$ est complet alors $B(X,Y)$ est complet
            \item Si $Y$ est complet alors $C_b(X,Y)$ est complet.
        \end{enumerate}
    \end{prop}

    \begin{proof}
        Si $(f_n)_\nN$ une suite de $B(X,Y)$ et $\funto{f}{X}{Y}$ une application telle que $d_\infty(f_n,f)\to0$.\\
        Soit $x,x_0\in X$ et $\nN$,
        \begin{align*}
            d_Y(f(x),f(x_n))&\le d_Y(f(x),f_n(x))+d_Y((f_n(x),f_n(x_0)))+d_Y(f_n(x_0),f(x_0))\\
            &\le 2 d_\infty (f_n,f)+\sup_{z,y\in X}d_Y(f_n(z),f_n(y))
        \end{align*}
        On en déduit que
        \begin{enumerate}
            \item $f$ est bornée puisque $f_n$ l'est
            \item Si $f_n$ est continue en $x_0$ et à partir d'un certain rang $f$ aussi
        \end{enumerate}
        On a donc $(1)$.\\
        Montrons $(2)$. Soit $(f_n)_{\nN}$ une suite de Cauchy (pour $d_\infty$) de $B(X,Y)$.\\
        Soit $\epsilon>0$ et $N_\epsilon \in \N$. $\forall n,m\ge N_{\epsilon},d_\infty(f_n,f_m)< \epsilon$.
        Pour tout $x\in X, d_Y(f_n(x),f_m(x))\le d_\infty (f_n,f_m)<\epsilon$. Donc la suite $(f_n(x))_n$ est de Cauchy dans $Y$ qui est complet donc elle converge et on note $f(x)=\lim_{n\to \infty}f_n(x)$\\
        En faisant tendre $n$ vers $+\infty$ dans $(*)$, on obtient
        \[\forall x\in X, d_Y(f_n(x),f(x))\le \epsilon\]
        Donc $d_\infty(f_n,f)\underset{n\to \infty}{\longrightarrow}$ d'après $(1)$ ci-dessus, $f$ est bornée puisque $f_n$ est bornée pour tout $\nN$.
        Montrons $(3)$, $C_b(X,Y)$ est complet car fermé dans $B(X,Y)$ qui est complet. 
    \end{proof}

    On en déduit de la preuve précédente le corollaire suivant

    \begin{cor}{}{}
        Soit $X$ un ensemble non-vide, $(Y,d_Y)$ un espace métrique et $(g_n)_\nN$ une suite de Cauchy de fonctions de $X$ dans $Y$ (pour $d_\infty$). On suppose que pour tout $x\in X$, la suite $(g_n(x))_\nN$ converge vers une limite que l'on note $g(x)$. Alors la suite $(g_n)_\nN$ converge uniformément vers $g$.\\
    \end{cor}

    \section{Parties relativement compactes de $C(X,Y)$ et théorème d'Ascoli}

    \begin{df}{}{}
        Soit $A\subset C_b(X,Y)$ et $x_0\in X$. On dit que $A$ est équicontinue en $x_0$ si:
        \[\forall \epsilon>0,\exists \delta>0,\forall f\in A,\forall x\in X, d_X(x_0,x)<\delta\implies d_Y(f(x_0),f(x))<\epsilon\]
        On dit que $A$ est équicontinue si $A$ est équicontinue en tout point de $X$.
    \end{df}
    
    \begin{ex}{}{}
        Pour $M>0$, l'ensemble $A_M=\{\funto{f}{X}{Y},d_Y(f(x),f(y))\le M d_X(x,y)\}$ des fonctions $M$-lipschitziennes est équicontinue
    \end{ex}

    \begin{prop}{}{}
        Lorsque $X$ est un espace métrique compact, une partie $A$ de $C(X,Y)$ est équicontinue \ssi elle est uniformément équicontinue.
        \[\forall \epsilon>0,\exists \delta>0,\forall f\in A,\forall x,y\in X, d_X(x,y)<\delta\implies d_Y(f(x),f(y))<\epsilon\]
    \end{prop}

    \begin{proof}
        Résulte d'une adaptation du théorème de Heine
    \end{proof}

    \begin{tm}{Théorème de Arzelà-Ascoli}{}
        Soit $(X,d)$ un espace métrique compact non-vide, $(Y,d)$ un espace métrique et $A$ une partie de $C(X,Y)$. Les assertions suivantes sont équivalentes:
        \begin{enumerate}
            \item $A$ est relativement compacte dans $(C(X,Y),d_\infty)$: $\overline{A}$ est compact $C(X,Y)$
            \item $A$ est (uniformément) équicontinue et pour tout $x\in X,\{ f(x),f\in A\}$ est relativement compact dans $Y$ (ie $\overline{\{f(x),f\in A\}}$ est compact)
        \end{enumerate}
    \end{tm}

    \begin{proof}
        $(1)\implies (2)$ Pour tout $x\in X, \fundfto{\varphi_x}{C(X,Y)}{Y}{f}{f(x)}$ est $1$-lipschitzienne donc continue.\\
        Donc si $\overline{A}$ est compact, $\varphi_x(\overline{A})=\{f(x),f\in \overline{A}\}$ est compact dans $Y$. De plus $\{f(x),f\in A\}\subset \{f(x),f\in \overline{A}\}$ donc $\overline{\{f(x),f\in A\}}$ est un fermé dans un compact, donc est compact.\\
        Montrons que $A$ est uniformément équicontinue. Soit $\epsilon>0,x_0\in X$ et $r>0$, \\
        on définie $O_{x_0,r}^\epsilon=\{f\in C(X,Y),\sup_{x\in B(x_0,r)} d_Y(f(x),f(x_0))<\epsilon\}$\\
        Objectif: Montrer que $A$ est continue dans un des $O_{x_0,r}^\epsilon$ ce qui donne l'équicontinuité en $x_0$. Pour cela, on montre que les $O_{x_0,r}^\epsilon$ forment un recouvrement ouvert de $C(X,Y)$ puis on utilise $\overline{A}$ compact.\\
        Si $f\in O_{x_0,r}^\epsilon$ alors soit $\rho:=\sup_{x\in B(x_0,r)}\epsilon -d_Y(f(x),f(x_0))>0$.\\
        Si $g\in B_{C(X,Y)}(f,\frac{1}{2}\rho)$ alors pour tout $x\in B_X(x_0,r)$
        \[d_Y(g(x),g(x_0))\le \underset{\le \frac{1}{2}\rho}{d_Y(g(x),f(x))}+d_Y(f(x),f(x_0))+\underset{\le \frac{1}{2}\rho}{d_Y(f(x_0),g(x_0))}\le \epsilon\]
        Donc $B_{C(X,Y)}\subset O_{x_0,r}^\epsilon$ et $O_{x_0,r}^\epsilon$ est ouvert dans $C(X,Y)$.\\
        On a évidemment que $\bigcup_{r>0}O_{x_0,r}^\epsilon=C(X,Y)$ et comme $\overline{A}$ est compact, il existe $r_{x_0}>0$ (par le lemme de Lebesgue) tel que $A\subset O_{x_0,r_{x_0}}^\epsilon$ (car les ouverts sont emboîtés lorsque $r$ vraie). Donc $A$ est équicontinue en $x_0$ et ceci pour tout $x_0$.\\
        $(2)\implies (1)$ Soit $(f_n)_\nN$ une suite de $\overline{A}$, montrons qu'elle admet une sous-suite convergente dans $C(X,Y)$
        \begin{itemize}
            \item On peut considérer les $f_n$ sont dans $A$. En effet par définition de $\overline{A}$, il existe $(g_n)_\nN$ une suite de $A$ telle que $d_\infty(f_n,g_n)\underset{n\to\infty}{\longrightarrow}0$. Donc $(f_n)_\nN$ admet une sous-suite $(f_{\varphi(n)})_\nN$ convergente vers $f$ \ssi $(g_{\varphi(n)})_\nN$ converge vers $f$. On suppose que $(f_n)_\nN\in A^\N$
            \item Puisque $X$ est compact, il est réparable. Soit $D=\{x_k,k\in \N\}\subset X$ telle que $\overline{D}=X$. On va construire une suite extraite de $(f_n)_\nN$ ponctuellement sur $D$.
        \end{itemize}
        Par hypothèse, $\overline{\{f(x_k),f\in A\}}$ est compact dans $Y$ pour tout $k\in \N$. La suite $(f_n(x_k))_\nN$ de $\overline{\{f(x_k),f\in A\}}$ admet donc une sous-suite convergente.\\
        On construit alors par réccurence, pour tout $j\in \N$ une suite d'extractices $\funto{\varphi_j}{\N}{\N}$ telle que la suite $(f_{\varphi_0\circ\varphi_1\circ\ldots\circ\varphi_k(n)}(x_k))_n$ converge pour tout $k\in \N$ (comme dans la preuve de Tikhinov métrique dénombrable)\\
        On obtient une sous-suite $(f_{\varphi(n)})_n$ de $(f_n)_\nN$ où $\varphi(n)=\varphi_0\circ \ldots\circ \varphi_n(n)$ (Procédé diagonal de Cantor) qui est telle que pour tout $k\in \N$, $f_{\varphi(n)}(x_k)$ converge lorsque $n$ tend vers $\infty$ vers une limite que l'on note $f(x_k)$.\\
        Il reste à montrer que $A$ est relativement compacte dans $C(X,Y)$
    \end{proof}

    \begin{re}{}{}
        \begin{enumerate}
            \item Si $X=\{1,\ldots,n\}$ avec la métrique discrète $d(x,y)=1$ pour $x\neq y$, alors $\Cr(X,\R)=\R^n\qquad (*)$.\\
            La condition $\overline{\{f(x),f\in A\}}$ est compact signifie la bornitude et l'uniforme équi-continuité est toujours vérifiée.
        \end{enumerate}
    \end{re}

    \section{Théorème de Stone-Weierstrass}

    Soit $X$ un espace compact et $\K=\R$ ou $\C$. On considère l'algèbre $\Cr(X,\K)$

    \subsection*{Structure de $\Cr(X,\K)$}

    $\Cr(X,\K)$ est une algèbre sur $\K$ de Banach unitaire. C'est à dire que:
    \begin{itemize}
        \item $\Cr(X,\K)$ est un espace métrique pour $d_\infty$
        \item $\Cr(X,\K)$ est un $\K$-\ev et la norme $d_\infty$ est associée à la norme $\norme{f}_\infty=\sup \{|f(x)|,x\in X\}$
        \item $\Cr(X,\K)$ est muni d'un produit $\fun{(f,g)}{f\times g}$ où $f\times g(x)=f(x)\times g(x),\forall x\in X$.\\
        Compatible avec la structure d'espace vectoriel : $f(g+\lambda h)=f g+\lambda f h$
        \item $\norme{fg}_\infty\le \norme{f}_\infty\norme{g}_\infty$ ce qui implique que le produit est continue,
        \item $\Cr(X,\K)$ est complet donc de Banach
        \item La fonction $\fundfto{\1}{X}{\R}{x}{1}$ est dans $\Cr(X,\K)$ et $\1$ est l'unité pour $\Cr(X,\K)$
    \end{itemize}

    \begin{df}{}{}
        Une sous-algèbre $A$ de $\Cr(X,\K)$ est un sous-espace vectoriel stable par produit. Elle est dite unitaire si $\1\in A$.
    \end{df}

    \begin{re}{}{}
        Si $A$ est une sous-algèbre de $\Cr(X,\K)$, par continuité des opérations $\overline{A}$ aussi.
    \end{re}

    \begin{df}{}{}
        Une partie $A$ de $\Cr(X,\K)$ est séparante si
        \[\forall x,y\in X,\exists f\in A, x\neq y\implies f(x)\neq f(y)\]
    \end{df}

    \begin{tm}{Théorème de Stone-Weierstrass(Cas réel)}{}
        Soit $X$ un espace topologique compact. Toute sous-algèbre unitaire séparante $A$ de $\Cr(X,\R)$ est dense pour la topologie de la convergence uniforme $(d_\infty)$
    \end{tm}

    \begin{lm}{}{}
        Si $A\subset C(X,\R)$ est une sous-algèbre unitaire qui sépare les points:
        \[\forall x,y\in X,\forall a,b\in \R,\exists f\in A, f(x)=a\text{ et }f(y)=b\]
    \end{lm}

    \begin{proof}
        Soit $g\in A$ tel que $g(x)\neq g(y)$, alors puisque $A$ est une sous-algèbre unitaire:
        \[f=a+\frac{b-a}{g(y)-g(x)}(g-g(x))\text{ convient}\]
        La condition $A$ unitaire peut être remplacée par la condition:
        \[\forall x\in X,\exists f\in A,f(x)\neq 0\]
        On remplace alors "$a$" par "$a\frac{f}{f(x)}$" dans cette preuve
    \end{proof}

    \begin{lm}{}{}
        La fonction sur $[0,1]$ $t \mapsto \sqrt{t}$ est limite uniforme de polynômes réels en $t$.
    \end{lm}

    \begin{proof}
        On construit par réccurence une suite de polynômes $(P_n)_\nN$ en posant:
        $\cho{P_0(t)=0\\P_{n+1}(t)=P_n(t)=\frac{1}{2}(t-(P_n(t))^2)\quad (*)}$\\
        On montre par réccurence que (Laisser en exercice)
        \begin{enumerate}
            \item $\forall t\in [0,1],(P_n(t))_\nN$ est croissante
            \item $\forall t\in [0,1],\forall \nN, 0<\sqrt{t}-P_n(t)\le (1-t)^n(\sqrt{t}-t)$\\
        \end{enumerate}
        Pour $t\in [0,1]$, la suite croissante majorée $(P_n(t))_\nN$ dans $[0,\sqrt{t}]$ converge vers une limite notée $f(t)$. Par passage à la limite dans $(*)$, $f(t)=f(t)+\frac{1}{2}(t-(f(t))^2)$ d'où $f(t)=\sqrt{t}$\\
        L'inégalité $\sqrt{t}-P_n(t)\le (1-t)^n(\sqrt{t}-t)$ permet d'obtenir que la convergence est uniforme.
    \end{proof}

    \begin{lm}{}{}
        Soit $A$ une sous-algèbre unitaire fermée de $C(X,\R)$ où $X$ est compact.
        \begin{enumerate}
            \item Si $f\in A$, $|f|\in A$
            \item Pour $f,g$ dans $A$, les applications
            \[\fundfto{\min(f,g)}{X}{\R}{x}{\min(f(x),g(x))}\text{ et }\fundfto{\max (f,g)}{X}{\R}{x}{\max(f(x),g(x))}\text{ sont dans $A$}\]
        \end{enumerate}
    \end{lm}

    \begin{re}{}{}
        \begin{itemize}
            \item Si $A$ est une sous-algèbre unitaire de $C(X,\R)$, alors $\overline{A}$ est une sous-algèbre unitaire fermée de $C(X,\R)$ et on pourra lui appliquer le lemme
            \item Si $A$ est une $\K$-algèbre unitaire, pour tout $P\in \K[X]$ et tout $f\in A$, $P(f)\in A$
        \end{itemize}
    \end{re}

    \begin{proof}
        \begin{enumerate}
            \item Si $f\in A$ est à valeurs dans $\R_+$, soit $a\in \R_+^*$ tel que $a\norme{f}_\infty\le 1$ et $\fundfto{h}{X}{[0,1]}{x}{af(x)}$. On note $\fundfto{g}{[0,1]}{\R}{t}{\sqrt{t}}$\\
            D'après le lemme 2, $g\circ f$ est limite uniforme de la suite $(P_n\circ h)_\nN$, donc pour tout $x\in \N, P_n\circ f\in A$ et sa limite $g\circ h\in \overline{A}=A$ où $g\circ h(x)=\sqrt{af(x)}$. Comme $\sqrt{a}\neq 0,\sqrt f\in A$.\\
            Enfin pour $f\in A$ (quelconque), $f^{2}\in A$ est à valeurs dans $\R_+$ et $|f|=\sqrt{f^2}\in A$
            \item On a $\min(f,g)=\frac{f+g-|f-g|}{2}$ et $\max(f,g)=\frac{f+g+|f-g|}{2}$ 
        \end{enumerate}
    \end{proof}

    \begin{re}{}{}
        Une réccurence immédiate montre que sous les hypothèses du lemme 3, si $J=\{f_1,\ldots,f_n\}$ est une partie finie de $A$, $\min(J)$ et $\max(J)$ sont dans $A$
    \end{re}

    \begin{lm}{}{}
        Soit $A$ une sous-algèbre unitaire fermée de $C(X,\R)$ qui sépare les points, où $X$ est compact. On a
        \[\forall f\in C(X,\R),\forall \epsilon>0,\exists g\in A \norme{g-f}_\infty <\epsilon\]
    \end{lm}

    \begin{proof}
        Soit $f\in C(X,\R),\epsilon>0$ et $I=\{h\in A,\forall y\in X,h(y)>f(y)-\epsilon\}$. Pour $h\in I$, on pose $V_n=\{x\in X,h(x)<f(x)+\epsilon\}$ de sorte que pour $h\in I$ et $x\in V_n,|h(x)|<\epsilon$. Le sous-ensemble $V_n$ est un ouvert de $X$ car $h$ et $f$ sont continues.\\
        Montrons que les $V_n$ forment un recouvrement ouvert de $X$. Soit $x\in X$, on va montrer qu'il existe $g_x\in I$ tel que $x\in V_{g_x}$. Notons $I_x=\{g\in A,g(x)=f(x)\}$ et pour $g\in I_x$, $U_g=\{y\in X,g(y)>f(y)-\epsilon\}$ qui est ouvert dans $X$. D'après le lemme 1
        \[\forall y\in X,\exists g\in A:f(x)=g(x)\text{ et }f(y)=g(y)\]
        Alors $g\in I_x,y\in U_g$ (donc $I_x$ et $U_g$ sont non-vides).\\
        Donc $(U_g)_{g\in I_x}$ forme un recouvrement ouvert de $X$ dont on peut extraire un sous-recouvrement fini $(U_g)_{g\in J_x}$ où $J_x\subset I_x$ est fini. Posons $g_x=\max(J_x)$, alors $g_x\in A$ d'après le lemme 3. On a évidemment $g_x(x)=f(x)$.\\
        Pour tout $z\in X, g_x(z)=\max_{h\in J_x} h(z)$ donc:\\
        $\forall y\in X,\exists g\in J_x$, tel que $y\in U_g$ et $f(y)-\epsilon<g(y)\le g_x(y)$. Autrement dit: $g_x\in I, g_x(x)=f(x)$ et donc $x\in V_{g_x}$. Donc $(V_n)_{n\in I}$ forme un recouvrement ouvert de $X$ qui est compact et dont on extrait un sous-recouvrement fini: $(V_n)_{n\in J}$ où $J\subset I$ est fini.\\
        Soit $g=\min (J)$ qui est dans $A$ par le lemme 3. Soit $x\in X$, on a:
        \begin{itemize}
            \item $\forall h\in J: h(x)>f(x)-\epsilon$ donc $g(x)>f(x)-\epsilon$
            \item $\exists h\in J, x\in V_n$ c'est à dire $h(x)<f(x)+\epsilon$ et donc $g(x)\le h(x)<f(x)+\epsilon$. Donc $|f(x)-g(x)|<\epsilon$. Enfin, puisque $X$ est compact, $\sup_{x\in X}\{|f(x)-g(x)|\}$ est atteint d'où $\norme{f-g}_\infty<\epsilon$
        \end{itemize}
    \end{proof}

    \begin{proof}{Preuve du théorème de Stone-Weierstrass (cas réel)}{}
        Soit $A$ une sous-algèbre unitaire qui répare les points de $C(X,\R)$ où $X$ est compact et $f\in C(X,\R)$, alors d'après le lemme 4,
        \[\forall n\in N,\exists f_n\in \overline{A}:\norme{f-f_n}_\infty<\frac{1}{n+1}\]
        Donc la suite $(f_n)_\nN$ converge vers $f$ et comme $\overline{A}$ est fermée $f\in \overline{A}$ donc $\overline{A}=C(X,\R)$
    \end{proof}

    \begin{re}{}{}
        Le théorème de Stone-Weierstrass ci-dessus ne peut pas s'étendre tel quel à $C(X,\C)$. En effet, si $A=\C[Z]$ et $X=\{z\in \C,|z|\le 1\}$, pour tout $P\in C(X,\C)$
        \[\int_{|z|=1}P(z)\, dz=i\int_{0}^{2\pi} P(e^{i\theta})e^{i\theta}\, d\theta=0\]
        Si $\C[Z]$ était dense dans $C(X,\C)$, par continuité de l'intégrale, on aurait $\int_{|z|=1}f(z)\, dz=0$ pour tout $f\in C(X,\C)$. Ce qui est faux, par exemple: $\int_{|z|=1}\overline{z}\, dz=2\pi i$
    \end{re}

    \begin{tm}{Théorème de Stone-Weierstrass (Cas complexe)}{}
        Soit $X$ un espace topologique compact. Toute sous-algèbre unitaire $A$ de $C(X,\C)$ qui sépare les points et qui est stable par conjugaison complexe est dense pour la topologie de la convergence uniforme.
    \end{tm}

    \begin{proof}
        Soit $f\in C(X,\C)$, si $f\in A$, puisque $A$ est stable par conjugaison complexe, alors $\Re(f)=\frac{1}{2}(f+\overline{f})$ et $\Im(f)=\Re(-if)$ sont dans $A$. Soit $x,y$ dans $X$, $x\neq y$ et $f\in A$ tel que $f(x)\neq f(y)$ alors nécessairement $\Re(f(x))\neq \Re (f(y))$ ou $\Im(f(x))\neq \Im (f(y))$. Ainsi $A_\R=\{\Re(f),f\in A\}=A\cap C(X,\R)$ est une sous-algèbre unitaire séparante de $C(X,\R)$, elle est donc dense dans $C(X,\R)$ par Stone-Weierstrass réel.\\
        Si $f\in C(X,\C),f=\Re(f)+ i \Im(f)$, si $\epsilon >0$, $h$ et $g$ dans $A_\R$ tel que $d_\infty(\Re(f),h)<\epsilon$ et $d_\infty (\Im(f),g)<\epsilon$ alors $h+ig\in A$ et $d_\infty(f,h+ig)<\sqrt{2}\epsilon$
    \end{proof}

    \begin{cor}{}{}
        Soit $X$ compacte de $\R$, l'ensemble des fonctions polynomiales est dense dans $C(X,\R)$
    \end{cor}

    \begin{proof}
        L'ensemble des fonctions polynomiales de $X\subset \R$ dans $\R$ est une algèbre unitaire qui sépare les points.
    \end{proof}

    \begin{re}{}{}
        Autrement dit, $\forall \funto{f}{X}{\R}$ continue, $\exists$ suite de polynôme qui converge vers $f$
    \end{re}

    \begin{cor}{Théorème de Weierstrass trigonométrique}{}
        Toute fonction continue $2\pi$-périodique de $\R$ dans $\C$ est limité uniforme sur $\R$ d'une suite de polynômes trigonométriques.
    \end{cor}

    \begin{re}{}{}
        Un polynôme trigonométrique est une fonction de $\R$ dans $\C$ de la forme:
        \[T(\theta)=\sum_{k=-n}^{n}a_n e^{ik\theta},n\in \N\]
    \end{re}

    \begin{proof}
        On identifie les fonctions $2\pi$-périodique continue de $\R$ dans $\C$ aux fonctions continues de $S_1$ dans $\C$ via $F(e^{i\theta})=f(\theta)$ puis on applique Stone-Weierstrass complexe
    \end{proof}

    \chapter{Espaces vectoriels normés}

    Dans tout ce chapitre, $\K=\R$ ou $\C$ et $E$ et $F$ sont des $\K$-\evs

    \section{Propriétés des normes}

    \begin{ra}{}{}
        \begin{itemize}
            \item Une norme est une application $\funto{N}{E}{\R_+}$ telle que
            \begin{itemize}
                \item $N(x)=0\iff x=0$
                \item $\forall \lambda\in \R,\forall x\in E, N(\lambda x)=|\lambda| N(x)$
                \item $\forall x,y\in E, N(x,y)\le N(x)+N(y)$
            \end{itemize}
            \item Un \evn est un couple $(E,N)$ où $E$ \ev et $N$ est une norme sur $E$.
            \item Lorsque $(E,N)$ \evn, $\fundfto{d_N}{E\times E}{\R_+}{(x,y)}{N(x-y)}$ est la métrique associée à $N$
        \end{itemize}
    \end{ra}

    \begin{ex}{}{}
        \begin{enumerate}
            \item $(\K^p,\norme{\cdot}_p)$ pour $p\in[1,+\infty]$ où pour $x=(x_1,\ldots x_n)\in \K^n$
            \begin{itemize}
                \item Si $p\in [1,+\infty[, \norme{x}_p=(|x_1|^p+\ldots+|x_n|^p)^{\frac{1}{p}}$
                \item Si $p=+\infty, \norme{x}_\infty=\max_{i\in \llbracket 1,n \rrbracket}|x_i|$
            \end{itemize}
            \item Ces normes s'étendent aux suites à valeurs dans $\K$ pour $x=(x_n)_\nN\in \K^\N$
            \begin{itemize}
                \item $p\in [1,+\infty[\implies \norme{x}_p=(\sum_{i\in \N}|x_i|^p)^{\frac{1}{p}}$
                \item $p=+\infty\implies \norme{x}_\infty=\sup_{i\in \N}|x_i|$.
            \end{itemize}
            On désigne par $l^p(\N,\K)$ l'espace des suites $(x_n)_\nN$ telles que $\norme{x}_p< +\infty$
            \item Ces normes s'étendent à l'espace des fonctions continues $\funto{f}{[0,1]}{\K}$
            $\norme{f}_p=(\int_{0}^{1}|f(t)|^p\, dt)^{\frac{1}{p}}$ et $\norme{f}_\infty=\sup_{t\in [0,1]}|f(t)|$ puis avec espaces mesurables...
            \item L'ensemble $C^k([0,1],\K)$ des fonctions $k$-fois dérivables et dont la k-ième dérivée est continue avec:
            \[\norme{f}=\max_{i\in \llbracket 0,k\rrbracket}\norme{f^{(i)}}_\infty\]
        \end{enumerate}
    \end{ex}

    \begin{re}{}{}
        La difficulté pincipale principale est de montrer que l'on a bien des normes dans les exemples 1 à 3 est de vérifier l'inégalité triangulaire. Cela passe par l'inégalité de Hölder (qui généralise Cauchy-Schwartz):\\
        Soient $p,q\in [0,\infty]$ tels que $\frac{1}{p}+\frac{1}{q}=1$, pour $x=(x_1,\ldots,x_n)$ et $y=(y_1,\ldots,y_n)$ dans $\K^n$, on a:
        \[\left|\sum_{i=1}^{n} x_i y_i\right|< \norme{x}_p\norme{y}_q\text{ (Inégalité de Hölder)}\]
        Dont on déduit (avec un peu de travail) $\norme{x}_p=\sup\left\{|\sum_{i=1}^{n}x_i y_i|,y\in \K^n\text{ et }\norme{y}_q\le 1\right\}$. Puis $\norme{x+y}_p\le \norme{x}_p+\norme{y}_p$ (Inégalité de Minkowski)
    \end{re}

    \begin{prop}{}{}
        Soit $(E,N)$ un espace vectoriel normé.
        \begin{enumerate}
            \item L'application $\funto{N}{(E,d_N)}{\R_+}$ est $1$-lipschitienne (donc continue)
            \item La topologie induite par $d_N$ sur $E$ est compatible avec les opérations d'espaces vectoriels.
            \[\fundf{E\times E}{E}{(x,y)}{x+y}\text{ et }\fundf{\K\times E}{E}{(\lambda,x)}{\lambda x}\text{ sont continues}\]
            \item Contrairement à ce qu'il se passe dans un espace métrique en général,
            \begin{enumerate}
                \item Si $x\in E$ et $r\in \R_+^*, \overline{B(x,r)}=\overline{B}(x,r)$
                \item Si $x\in E$ et $r\in \R_+^*, \ring{\overline{B}(x,r)}=B(x,r)$
            \end{enumerate}
            \item L'adhérence d'un sous-espace vectoriel est un sous-espace vectoriel
        \end{enumerate}
    \end{prop}

    \begin{proof}
        Laisser en exercice
    \end{proof}

    \begin{df}{}{}
        Un espace vectoriel complet est un espace de Banach
    \end{df}

    \begin{re}{}{}
        La notion de métriques équivalentes induit la notion de normes équivalentes. Deux normes $N_1$ et $N_2$ sur $E$ sont équivalentes, si il existe $a,b\in \R_+^*$ tels que
        \[\forall x\in E, aN_1(x)\le N_2(x)\le bN_1(x)\]
        Donc pour $N_1$ équivalente à $N_2$, on a que $(E,N_1)$ est un Banach si et seulement si $(E,N_2)$ est un Banach.
    \end{re}

    \subsection*{Convergence des séries dans un Banach}

    \begin{df}{}{}
        Soit $(E,\norme{\cdot}_E)$ un \evn et $(x_n)_\nN \in E^\N$
        \begin{enumerate}
            \item La série $\sum_{\nN} x_n$ converge lorsque la suite $(S_n)_\nN$ de $E^\N$ converge où $S_n=\sum_{i=0}^{n} x_i$
            \item La série $\sum_{\nN} x_n$ est normalement convergente lorsque la suite $(T_n)_\nN\in \R^\N$ converge où $T_n=\sum_{i=0}^{n}\norme{x_i}_E$
        \end{enumerate}
    \end{df}

    \begin{re}{}{}
        La suite $(T_n)_\nN$ est croissante donc $\sum_{n\in \N}x_n$ est normalement convergente $\iff \exists M>0$ tel que $\forall n\in \N, \sum_{i=0}^{n}\norme{x_i}_E<M$
    \end{re}

    \begin{prop}{}{}
        Dans un espace de Banach, une série normalement convergente est convergente
    \end{prop}

    \begin{proof}
        On suppose que la suite $(T_n)_\nN$ converge où $T_n=\sum_{i=0}^{n}\norme{x_i}_E$. Soit $S_n=\sum_{i=0}^{n}x_i$, alors pour $p,q\in \N$, $p\le q$:
        \[\norme{S_p-S_q}_E=\norme{\sum_{i=p}^{q-1}}_E\le \sum_{i=p}^{q-1}\norme{x_i}_E=T_q-T_p\]
        Puisque $(T_n)_\nN$ converge, c'est une suite de Cauchy donc $(S_n)_\nN$ est de Cauchy dans un Banach, donc elle converge.
    \end{proof}

    \begin{re}{}{}
        Dans un \evn non-complet, il existe des séries normalement convergentes, mais non-convergentes. Il suffit de considérer une suite de Cauchy $(x_n)_\nN$ qui ne converge pas. Quitte à considérer une suite extraite, on peut supposer que $\norme{x_{n+1}-x_n}\le 2^{-n}$. Alors la série $\sum (x_{n+1}-x_n)$ est normalement convergente mais non-convergente.
    \end{re}

    \section{Applications linéaires continues}

    \begin{tm}{}{}
        Soit $(E,\norme{\cdot}_E)$ et $(F,\norme{\cdot}_F)$ deux $\K$-espaces vectoriels normés et $\funto{f}{E}{F}$ une application linéaire, sont équivalentes:
        \begin{enumerate}
            \item $f$ est continue en $0_E$
            \item $f$ est continue
            \item $f$ est uniformément continue
            \item $f$ est lipschitzienne
            \item $\exists k\in \R_+$ tel que $\forall x\in E, \norme{f(x)}_F\le k\norme{x}_E$
        \end{enumerate}
    \end{tm}

    \begin{proof}
        En notant que pour $x,y\in E$, la linéarité de $f$ donne:
        \[f(x-y)=f(x)-f(y)\text{, on a évidemment }5\implies 4\implies 3\implies 2\implies 1\]
        Il suffit de montrer que $1\implies 5$. Si $f$ est continue en $0_E$, il existe $\eta>0$ tel que:
        \[\forall x\in E, \norme{x}_E<\eta,\text{ alors, }\norme{f(x)}_F\le 1\]
        Donc pour $x\in E\wt\{0_E\},\norme{\frac{\eta x}{\norme{x}_E}}_E=\eta$ et
        \[\norme{f(\frac{\eta x}{\norme{x}_E})}_F=\frac{\eta}{\norme{x}_E}\norme{f(x)}_E\le 1\]
        D'où $\norme{f(x)}_F\le \frac{1}{\eta}\norme{x}_E$. Pour $x=0_E$, l'inégalité est encore vérifiée.
    \end{proof}

    \begin{cor}{}{}
        Deux normes sur $E$ sont équivalentes si et seulement si elles définissent la même topologie
    \end{cor}

    \begin{proof}
        Les normes $N_1$ et $N_2$ sur $E$ sont équivalentes s'il existe $a,b\in \R_+^*$ tels que
        \[\forall x\in E, aN_1(x)\underset{(*)}{\le} N_2(x)\underset{(**)}{\le} bN_1(x)\]
        $(*)$ signifie que $\funto{Id}{(E,d_{N_2})}{(E,d_{N_1})}$ est continue et $(**)$ signifie que $\funto{Id}{(E,d_{N_1})}{(E,d_{N_2})}$ est continue.
    \end{proof}

    \begin{df}{}{}
        Soit $(E,\norme{\cdot}_E)$ et $(F,\norme{\cdot}_F)$ deux $\K$-\evn. On note $L(E,F)$ l'espace vectoriel des applications linéaires et continues de $E$ dans $F$.\\
        Pour $f\in L(E,F)$, la norme (d'opérateurs) de $f$, notée $\norme{f}_{L(E,F)}$ est
        \[\norme{f}_{L(E,F)}=\inf\{k\in \R_+^*,\forall x\in E, \norme{f(x)}_F\le k\norme{x}_E\}\]
        Ainsi $(L(E,F),\norme{\cdot}_{L(E,F)})$ est un \evn.
    \end{df}

    \begin{prop}{}{}
        Soit $(E,\norme{\cdot}_E), (F,\norme{\cdot}_F)$ deux $\K$-\evn et $f\in L(E,F)$. On a:
        \begin{enumerate}
            \item $\forall x\in E, \norme{f(x)}_F\le \norme{f}_{L(E,F)}\norme{x}_E$
            \item \begin{align*}\norme{f}_{L(E,F)}&=\sup_{\norme{x}_E<1}\\
                &=\sup_{\norme{x}_E\le 1} \norme{f(x)}_F\\
                &=\sup_{x\neq 0_E}\frac{\norme{f(x)}_F}{\norme{x}_E}\\
                &=\sup_{\norme{x}_E=1}\norme{f(x)}_F
                \end{align*}
            \item Si $(G,\norme{\cdot}_G)$ est un $\K$-\evn et $g\in L(F,G)$
            \[\norme{g\circ f}_{L(E,F)}\le \norme{g}_{L(E,F)}\norme{f}_{L(E,F)}\]
        \end{enumerate}
    \end{prop}

    \begin{proof}
        \begin{enumerate}
            \item Immédiat dès que l'on note
            \[A=\{\forall k\in \R_+,\norme{f(x)}_F\le k\norme{x}_E\}=\bigcap_{x\in E} \phi_x^{-1}([0,+\infty[)\]
            où $\fundfto{\phi_x}{\R_+}{\R}{k}{k\norme{x}_E -\norme{f(x)}_F}$\\
            $A$ est fermé, non-vide dans $\R$, minoré par $0$ donc $A$ contient une borne inférieure.
            \item Puisque $f(0_E)=0_F$, la condition $\norme{f(x)}_F\le k\norme{x}_E$ est toujours vraie pour $x=0_E$. Donc
            \begin{align*}
                \norme{f}_{L(E,F)}&=\inf \{k\in \R_+, \forall x\in E\wt \{0_E\},\norme{f(x)}_F\le k\norme{x}_E\}\\
                &=\inf \{k\in \R_+,\forall x\in E\wt \{0_E\},\frac{\norme{f(x)}_F}{\norme{x}_E}\le k\}\\
                &=\sup_{x\neq 0_E}\frac{\norme{f(x)}_F}{\norme{x}_E}
            \end{align*}
            On remarque que pour $x_0\neq 0$ fixé et $\lambda\in [0,1]$,
            \[\norme{f(\lambda x)}_F=|\lambda|\norme{f(x_0)}_F\le \norme{f(x_0)}_F\]
            Donc $\sup_{\norme{x}_E\le 1}\norme{f(x)}_F=\sup_{\norme{x}_E=1}\norme{f(x_0)}_F$.\\
            De plus, $\frac{\norme{f(x_0)}_F}{\norme{x_0}_E}=\norme{\frac{f(x_0)}{\norme{x_0}_E}}_F=\norme{f\left( \frac{x_0}{\norme{x_0}_E}\right)}_F$ et $\norme{\frac{x_0}{\norme{x_0}_E}}_E=1$ donc
            \[\sup_{\norme{x}_E=1}\norme{f(x_0)}_F=\sup_{x\neq 0_E}\frac{\norme{f(x)}_F}{\norme{x}_E}\]
            On a de plus, pour $\funto{g}{X}{\R}$, continue et $X$ métrique, si $A\subset X$ est tel que $g(A)$ est borné, alors
            \[\sup_{x\in A}g(x)=\sup_{x\in \overline{A}}g(x)\text{ (Laisser en exercice)}\]
            On applique ceci à $\fundfto{g}{E}{\R}{x}{\norme{g(x)}_F}$ et $A=B(0_E,A),\overline{A}=\overline{B}(0_E,1)$
            \item Laisser en exercice
        \end{enumerate}
    \end{proof}

    \begin{tm}{}{}
        Si $F$ est un espace de Banach, $L(E,F)$ aussi
    \end{tm}

    \begin{proof}
        Soit $(f_n)_\nN$ une suite de Cauchy de $L(E,F)$. Pour $p,q\in \N$ et $x\in E$
        \[\norme{f_p(x)-f_q(x)}_F\le \norme{f_p-f_q}_{L(E,F)}\norme{x}_E\]
        Donc si $(f_n)_\nN$ est de Cauchy dans $L(E,F)$, $(f_n(x))_\nN$ est de Cauchy dans $F$.\\
        Puisque $F$ est complet, la suite $(f_n(x))_\nN$ converge, on note $f(x)$ sa limite. Il reste à vérifier que $f\in L(E,F)$ et que $(f_n)_\nN$ converge vers $f$ pour la norme $\norme{\cdot}_{L(E,F)}$.\\
        $f$ linéaire: Soit $x,y\in E,\lambda \in \K$ et $\nN$,
        \begin{align*}
            f_n(x+y)&=f_n(x)+f_n(y)\\
            f_n(\lambda x)=\lambda f_n(x)
        \end{align*}
        Par continuité de la somme et du produit extérieur, on trouve à la limite:
        \[f(x+y)=f(x)+f(y)\text{ et }f(\lambda x)=\lambda f(x)\]
        Donc $f$ est linéaire\\
        $f$ est continue et limite de $(f_n)_\nN$ dans $L(E,F)$.\\
        Soit $\epsilon>0$ et $N\in \N$ tels que si $p,q\ge \N, \norme{f_p-f_q}_{L(E,F)}\le \epsilon$\\
        Pour $x\in E, \norme{f_p(x)-f_q(x)}_F\le \epsilon \norme{x}_E$. On fixe $q$ et on fait tendre $p$ vers $+\infty$:
        \[\norme{f(x)-f_q(x)}_F\le \epsilon \norme{x}_E\]
        Donc $f-f_q$ est continue et $\norme{f-f_q}_{L(E,F)}\le \epsilon\qquad (*)$\\
        Puisque $f_q$ est continue, $f=(f-f_q)+f_q$ est continue.\\
        D'après $(*)$, $\lim_{n\to +\infty}\norme{f-f_q}=0$ donc la suite $(f_n)_\nN$ converge vers $f$ dans $L(E,F)$
    \end{proof}
    
    \begin{re}{}{}
        En particulier, $L(E):=L(E,E)$ est une algèbre (Le produit est la composition), unitaire, généralement non commutative qui est une algèbre de Banach lorsque $E$ est de Banach.
    \end{re}

    \begin{ex}{}{}
        $M_n(\R)=L(\R^n,\R^n)$
    \end{ex}

    \begin{re}{}{}
        En dimension infinie, les applications linéaires ne sont pas toujours continues.\\
        Par exemple, $\fundfto{\varphi}{C([0,1],\R)}{\R}{f}{f(0)}$ est linéaire mais pas continue lorsque $C([0,1],\R)$ est muni de la topologie de la norme $\norme{\cdot}_p$ pour $p\in [1,+\infty[$
    \end{re}

    \begin{tm}{Prolongement des applications linéaires continues}{}
        Soit $F$ un Banach, $D$ un \sev de $E$ un \evn tel que $\overline{D}=E$.\\
        Soit $f\in L(D,F)$, il existe un unique prolongement $\tilde{f}\in L(E,F)$ de $f$. De plus, $\norme{\tilde{f}}_{L(E,F)}=\norme{f}_{L(D,F)}$
    \end{tm}

    \begin{proof}
        Puisque $f\in L(D,F)$, $f$ est uniformément continue. Par le théorème de prolongement des applications uniformément continues (CF Partiel Exo 3), on sait que $f$ se prolonge en une unique application uniformément continue $\funto{\tilde{f}}{E}{F}$.\\
        Soit $\lambda \in \K$, l'application $\fundf{E\times E}{F}{(x,y)}{\tilde{f}(x+\lambda y)-\tilde{f}(x)-\lambda \tilde{f}(y)}$ est continue et nulle sur $D\times D$ qui est dense dans $E\times E$. Donc elle est nulle et $\tilde{f}$ est linéaire.\\
        Soit $(x_n)_\nN$ une suite d'éléments de $D$ qui converge vers $x\in E$.\\
        Pour $\nN$,
        \[\norme{\tilde{f}(x_n)}_F=\norme{f(x_n)}_F\le \norme{f}_{L(D,F)}\norme{x_n}_E\]
        En passant à la limite quand $n\to +\infty$
        \[\norme{\tilde{f}(x)}_F\le \norme{f}_{L(D,F)}\norme{x}_E\]
        Donc $\norme{\tilde{f}}_{L(E,F)}\le \norme{f}_{L(D,F)}$, puisque $D\subset E$,
        \[\{k\in \R_+,\forall x\in E,\norme{f(x)}_F\le k\norme{x}_E\}\subset \{k\in \R_+,\forall x\in D, \norme{f(x)}_F\le k\norme{x}_E\}\]
        Donc $\norme{f}_{L(D,F)}\le \norme{f}_{L(E,F)}$
    \end{proof}

    \subsection*{Conséquences du théorème de Baire}

    \begin{ra}{}{}
        $(X,d)$ un espace métrique complet est un espace de Baire, toute réunion dénombrable de fermés d'intérieur vide est d'intérieur vide.
    \end{ra}

    \begin{cor}{}{}
        Si un espace métrique non-vide complet est réunion dénombrable de fermés $(F_n)_\nN$, alors
        \[\exists n\in \N, \mathring{F_n}\neq \emptyset\]
    \end{cor}

    \begin{tm}{Théorème de l'application ouverte}{}
        Soit $E,F$ deux espaces de Banach et $T\in L(E,F)$, si $T$ est surjective alors elle est ouverte
    \end{tm}

    \begin{proof}
        \begin{enumerate}
            \item Une bijection $\funto{f}{X}{Y}$ entre deux espaces métriques est un homéomorphisme \ssi:
            \[\forall A\subset X,f(\overline{A})=\overline{f(A)}\text{ et }f(\mathring{A})=\mathring{f(A)}\]
            Pour un \evn $E$, $x\in E$ et $\lambda\in \K$, les applications:
            $\fundfto{\tau_x}{E}{E}{y}{x+y}$ et $\fundfto{m_\lambda}{E}{E}{y}{\lambda y}$. sont des homéomorphismes et on a pour $r\in \R_+^*$,
            \[B_E(x,r)=\tau_x\circ m_r(B(0,1))\]
            De plus, si $F$ est un \evn et $\funto{T}{E}{F}$ une application linéaire,
            \[ T\circ \tau_x\circ m_r=\tau_x\circ m_r\text{ car }T(x+ry)=T(x)+rT(y),\forall y\in E\]
            Donc $T(B_E(x,r))=\tau_{T(x)}\circ m_r\circ T(B_E(0,1))$
            \item On a les équivalences:
            \begin{enumerate}
                \item $T$ est ouverte
                \item $\forall U\subset E$ ouvert, $T(U)$ est ouvert.
                \item $\forall x\in E,\forall r\in \R_+^*,\exists \varphi>0, B_F(T(x),\varphi)\subset T(B_E(x,r))$
                \item $\exists R>0$, $B_F(0,1)\subset T(B_E(0,R))$
            \end{enumerate}
            \item Puisque $T$ est surjective, $F=\bigcup_\nN F_n$ où
            \[F_n:= \overline{T(B_E(0,n))}\]
            Puisque $F$ est complet, il est de Baire. Donc il existe $n\in \N$ tel que $F_n\neq \emptyset$. Donc il existe $y\in F_N$ et $\eta>0$ tel que
            \[B(y,\eta)\subset T_N=\overline{T(B_E(0,N))}\]
            Pour $n$ assez grand, si on note $M=mN$, on a:
            \[B_F(0,1)\subset B_F(y,m\eta)\subset \overline{T(B_E(0,M))}=\overline{T(B_E(0,r))}\qquad (*)\]
            D'après $(*)$,
            \[\forall n\in\N, B_F(0,\left(\frac{1}{2}\right)^n)\subset \overline{T(B_E(0,\left(\frac{1}{2}\right)^nM))}\]
            Soit $y_0\in B_F(0,1)\subset \overline{T(B_E(0,M))}$, $B_F(y_0,\frac{1}{2})\cap T(B_E(0,M))\neq \emptyset$ donc il existe $x_0\in B_E(0,M)$ tel que $\norme{y_0-x_0}\le \frac{1}{2}$. On note $y_1=y_0-T(x_0)$.\\
            On construit par réccurence une suite $(x_n)_{n\in\N}$ de $E$ et une suite $(y_n)_{n\in \N}$ de $F$ telles que:
            \begin{itemize}
                \item $x_n\in B_E(0,\left(\frac{1}{2}\right)^nM)$
                \item $y_{n+1}=y_n-T(x_n)=y_0-T(x_0+\ldots+x_n)\in B_F(0_F,\left(\frac{1}{2}\right)^{n+1})$
            \end{itemize}
            La série $\sum x_n$ est normalement convergente et $E$ est un Banach donc elle converge vers une limite que l'on note $X$. De plus,
            \[\norme{X}_E=\norme{\sum_{n=0}^{\infty}x_n}_E\le \sum_{n=0}^{\infty}\norme{x_n}_E\le \sum_{n=0}^{\infty}2^{-n}M=2M\]
            Par passage à la limite, $y_0=T(X)\in T(B_E(0,3M))$ donc $B_F(0,1)\subset T(B_E(0,3M))$ et $T$ est ouverte
        \end{enumerate}
    \end{proof}

    \begin{cor}{}{}
        Soient $E$ et $F$ deux Banach. Toute application linéaire continue et bijective est un homéomorphisme.
    \end{cor}

    \begin{proof}
        Laisser en exercice
    \end{proof}

    \begin{tm}{Théorème du graphe fermé}{}
        Soit $E$ et $F$ deux Banach et $f$ une application linéaire de $E$ dans $F$. Alors $f$ est continue \ssi son graphe est fermé.
    \end{tm}

    \begin{proof}
        On sait déjà que si $\funto{T}{E}{F}$ est continue, alors le graphe de $T$ est fermé dans $E\times F$ muni de la norme $\norme{(x,y)}_\infty=\max(\norme{x}_E,\norme{y}_F)$.\\
        Inversement, puisque $E$ et $F$ sont des Banach, $(E\times F,\norme{\cdot}_\infty)$ est un Banach. Si $\funto{T}{E}{F}$ est linéaire, son graphe $G_T=\{(x,T(x)),x\in E\}$ est un \sev de $E\times F$ que l'on suppose fermé. Donc $G_T$ est un Banach. La projection $\funto{p_E}{E\times F}{E}$ est continue et sa restriction $p_{E|G_T}$ est bijective. C'est donc un homéomorphisme et sa réciproque est continue. Alors $T=p_F\circ\left(p_{E|G_T}\right)^{-1}$ est continue.
    \end{proof}

    \begin{tm}{Théorème de Banach-Steinhaus}{}
        Soit $E$ un Banach et $F$ un \evn et $A\subset L(E,F)$. Si pour tout $x\in E$, $\{f(x),f\in A\}$ est borné dans $F$, alors $A$ est bornée dans $L(E,F)$.\\
    \end{tm}

    \begin{proof}
        On suppose $E\neq \{0_E\}$. Soit $n\in \N$ et 
        \begin{align*}
            C_n&=\{x\in E,\forall f\in A,\norme{f(x)}_F\le n\}\\
            &=\bigcap_{f\in A}\{x\in E,\norme{f(x)}\le n\}\text{ qui est fermé dans $E$ comme intersection de fermés}
        \end{align*}
        Puisque $A\subset L(E,F)$, pour tout $x\in \N$,
        \[C_n=nC_1=\{nx,x\in C_1\}=m_n(C_1)\]
        Par hypothèse, $E$ est réunion des $C_n$, par le corollaire du théorème de Baire:
        \[N\in \N,\mathring{C_n}\neq\emptyset\]
        Où $m_n$ est un homéomorphisme donc $\mathring{C_1}\neq \emptyset$. Soit $x_0\in \mathring{C_1}$:\\
        \[\exists r>0, B(x_0,r)\subset C_1=\{x\in E,\forall f\in A,\norme{f(x)}\le 1\}\]
        Soit $x\in B(0_E,1)$, alors $x_0\pm rx\in B(x_0,r)$ et :
        \[\forall f\in A, f(x)=\frac{1}{2r}(f(x_0+rx)-f(x_0-rx))\text{ d'où}\]
        \[\norme{f(x)}_F\le \frac{1}{2r}(\norme{f(x_0+rx)}+\norme{f(x_0-rx)})\le \frac{1}{r}\]
        Donc $\forall f\in A$, $\norme{f}_{L(E,F)}\le \frac{1}{r}$.
    \end{proof}

    \begin{cor}{}{}
        Soit $E$ un Banach, $F$ un \evn. Soit $(T_n)_\nN$ une suite de $L(E,F)$ qui converge simplement vers $\funto{T}{E}{F}$. Alors $T\in L(E,F)$
    \end{cor}

    \begin{proof}
        Laisser en exercice
    \end{proof}

    \section{Propriétés liées à la dimension des \evn}

    \subsection*{Le cas des \evn de dimension finie}

    \begin{prop}{}{}
        Munissons $\K^n$ de la norme $\norme{\cdot}_\infty$:
        \begin{enumerate}
            \item Toute application linéaire de $\K^n$ dans un \evn est continue.
            \item Toute application linéaire bijective de $\K^n$ dans un \evn est un homéomorphisme.
        \end{enumerate}
    \end{prop}

    \begin{proof}
        Soient $(e_1,\ldots,e_n)$ la base canonique de $\K^n$, $(E,\norme{\cdot}_E)$ un \evn et $\funto{\varphi}{\K^n}{E}$ une application linéaire
        \begin{enumerate}
            \item $\forall x=(x_1,\ldots,x_n)\in \K^n$
            \[\norme{\varphi(x)}_E=\norme{x_1\varphi(e_1)+\ldots+x_n\varphi(e_n)}_E\le \norme{x}_\infty\sum_{i=1}^{n}\norme{\varphi(e_i)}_E\]
            Donc $\varphi$ est continue et $\norme{\varphi}_{L(\K^n,E)}\le \sum_{i=1}^{n}\norme{\varphi(e_i)}_E$
            \item Supposons $\varphi$ bijective et notons $S=\{x\in \K^n,\norme{x}_\infty=1\}$. $S$ est un compact de $\K^n$. L'application $\fundfto{\norme{\varphi(\cdot)}_E}{\K^n}{\R_+}{x}{\norme{\varphi(x)}_E}$ est continue (d'après $(1)$, $\varphi$ est continue et $\norme{\cdot}_E$ est continue).\\
            Puisque $f$ est injective, $\forall x\in S$, $\norme{\varphi(x)}_E>0$ et puisque $S$ est compact, il existe $a\in \R_+^*$ qui minore $\{\norme{\varphi(x)}_E,x\in S\}$\\
            Soit $x\in \K^n\wt \{0_{\K^n}\}$, alors $\frac{x}{\norme{x}_{\K^n}}\in S$ et $\norme{\varphi(\frac{x}{\norme{x}_\infty})}_E=\frac{\norme{\varphi(x)}_E}{\norme{x}_\infty}\ge a$ ou encore $\norme{\varphi(x)}_E\ge a\norme{x}_\infty\qquad (*)$.\\
            Cette inégalité est encore vérifiée en $x=0$. Finalement,
            \[\forall u\in E,u=\varphi(\varphi^{-1}(u))\text{ et }\norme{u}_E\ge a\norme{\varphi^{-1}(u)}_\infty\]
            Donc $\varphi^{-1}$ est continue et $\norme{\varphi^{-1}}_{L(E,\K^n)}\le a^{-1}$
        \end{enumerate}
    \end{proof}

    \begin{tm}{}{}
        Soit $E$ un \evn de dimension finie:
        \begin{enumerate}
            \item Toutes les normes sur $E$ sont équivalentes.
            \item On muni $E$ d'une norme. Toute application linéaire de $E$ dans un \evn est continue.
        \end{enumerate}
    \end{tm}

    \begin{proof}
        Soit $n=\dim E$ et $\funto{\varphi}{\K^n}{E}$ une application linéaire bijective.
        \begin{enumerate}
            \item Soient $N_1,N_2$ deux normes sur $E$. D'après la proposition précédente, les applications
            \[\varphi_1=\funto{\varphi}{(\K^n,\norme{\cdot}_\infty)}{(E,N_1)}\text{ et }\varphi_2=\funto{\varphi}{(\K^n,\norme{\cdot}_\infty)}{(E,N_2)}\]
            sont des homéomorphismes linéaires. Donc $\funto{\Id}{(E,N_1)}{(E,N_2)}$ qui vaut $\varphi_2\circ \varphi_1^{-1}$ est un homéomorphisme linéaire. Ce qui équivaut à dire que $N_1$ et $N_2$ sont équivalentes
            \item Soit $\funto{\Psi}{E}{F}$ une application linéaire où $F$ est un \evn, alors d'après la proposition précédente, $\Psi\circ\varphi$ et $\varphi^{-1}$ sont continues. Donc $\Psi=(\Psi\circ\varphi)\circ\varphi^{-1}$ aussi.
        \end{enumerate}
    \end{proof}

    \begin{prop}{}{}
        Soit $(E,\norme{\cdot}_E)$ un \evn de dimension finie.
        \begin{enumerate}
            \item L'espace $E$ est complet
            \item Une partie $A$ de $E$ est complète \ssi elle est fermée.
            \item Une partie $A$ de $E$ est compacte \ssi elle est fermée et bornée.
        \end{enumerate}
    \end{prop}

    \begin{proof}
        Soit $n=\dim E$ et $\funto{\varphi}{\K^n}{E}$ une application linéaire bijective. Puisque toutes les normes sur $E$ sont équivalentes, la norme $\norme{\cdot}_E$ est équivalente à $\norme{\varphi^{-1}(\cdot)}_\infty$.\\
        De plus, $\varphi$ est une isométrie, $\fun{(\K^n,\norme{\cdot}_\infty)}{(E,\norme{\varphi^{-1}(\cdot)}_\infty)}$.\\
        Il suffit donc de vérifier les propriétés pour $(\K^n,\norme{\cdot}_\infty)$
        \[(N(\varphi(x))=\norme{\varphi^{-1}(\varphi(x))}_\infty=\norme{x}_\infty)\]
    \end{proof}

    \begin{cor}{}{}
        Tout \sev de dimension finie d'un \evn est fermé
    \end{cor}

    \begin{proof}
        Laisser en exercice
    \end{proof}

    \subsection*{Le cas des \evn de dimension infinie}

    \begin{tm}{Théorème de Riesz}{}
        La boule unité fermée d'un \evn est compacte \ssi la dimension de $E$ est finie.
    \end{tm}

    \begin{lm}{}{}
        Soit $F$ un \sev fermé de $E$ et $y\in E\wt F$. Il existe $z\in \K y\oplus F$ tel que $\norme{z}=1$ et $d(z,F)>\frac{1}{2}$
    \end{lm}

    \begin{proof}
        Soit $\delta=d(y,F)=\inf \{\norme{z-y},z\in F\}$. Puisque $y\not \in F$ et $F$ est fermé, $\delta>0$. Soit $x\in F$ tel que $\norme{x-y}\le 2\delta$ et soit $z=\frac{y-x}{\norme{y-x}}$\\
        Alors pour $z\in S_E=\{e\in E,\norme{e}_E=1\}$ et pour $z'\in F$,
        \[\norme{z-z'}=\norme{\frac{1}{\norme{y-x}}(y-(x+\norme{y-x}z'))}\ge \frac{1}{\norme{y-x}}d(y,F)>\frac{1}{2}\]
    \end{proof}

    \begin{proof}
        Preuve du Théorème de Riesz: Si $\dim E$ est finie, la boule unité fermée étant fermée et bornée, elle est compacte.\\
        Réciproquement, si $\dim E=+\infty$, montrons que $\overline{B_E}(0,1)$ n'est pas compacte en construisant une suite qui n'admet que pas de sous-suite convergente.\\
        Soit $S_E=\delta B_E$ la sphère unité. Soit $x_0\in S_E$ et $F_1=\text{Vect}(x_0)$.\\
        Puisque $\dim F_1<+\infty$, $F_1$ est fermé dans $E$. Puisque $\dim E=+\infty$, il existe $y_1\in E\wt F$.\\
        D'après le lemme, $\exists x_1\in S_E\cap (\K y_1\oplus F_1)$ tel que $d(x_1,F_1)>\frac{1}{2}$ et donc $\norme{x_1-x_0}>\frac{1}{2}$.\\
        On construit aussi par réccurence une suite $(x_n)_n$ de $S_E$ une suite de \sev emboités $F_n=\text{Vect}(x_0,\ldots,x_{n-1}),\forall n\in \N$, une suite $(y_n)_{n\in\N^*}$ de parties de $E$ où $y_n\in E\wt F_n$ tels $x_n\in S_E\cap (\K y_n\oplus F_n)$ et $d(x_n,F_n)>\frac{1}{2}$. En particulier, $\norme{x_n-x_k}>\frac{1}{2}$ pour $k\le n-1$. Finalement, la suite $(x_n)_\nN$ de $S_E\subset \overline{B_E}$ est telle que
        \[p,q\in \N,p\neq q\implies \norme{x_p-x_q}\ge \frac{1}{2}\]
        Donc la suite $(x_n)_\nN$ n'admet pas de sous-suite convergente et $B_E$ n'est pas compact.
    \end{proof}

    \chapter{Formes linéaires continues}

    Dans toute la suite, $E$ est un $\K$-\evn, $\K=\R$ ou $\C$.

    \section{Dual d'un espace vectoriel normé}

    \subsection*{Dual topologique}

    Puisque $\R$ et $\C$ sont des espaces de Banach, $L(E,\K)$ est un espace de Banach

    \begin{df}{}{}
        Le dual topologique de $E$, que l'on note $E'$ est l'espace de Banach $L(E,\K)$.
    \end{df}

    \begin{ra}{}{}
        Pour $E$ un $\K$-\ev, une forme linéaire sur $E$ est une application linéaire de $E$ dans $\K$. On note $E^*$ l'ensemble des formes linéaires de $E$, c'est le dual algébrique.
    \end{ra}

    \begin{re}{}{}
        Si $\K=\C$, $E$ est aussi un $\R$-\evn, on peut alors considérer:
        \[E_\R'=L(E,\R)\text{ et }E_\C'=L(E,\C)\]
        Notons $\funto{\Re}{\C}{\R}$ qui à $z\in \C$ associe la partie réelle de $z$.
    \end{re}

    \begin{prop}{}{}
        L'application linéaire $\Re$ est une bijection isométrique
    \end{prop}

    \begin{proof}
        Notons que $\norme{\Re}=1$. Soit $l\in E_\C'$, on a:
        \[\norme{\Re-l}\le \norme{\Re}\norme{l}_{E_C'}=\norme{l}_{E_\C'}\]
        Soit $x\in \overline{B_E(0,1)}, \lambda\in \C$ telle que $\lambda l(x)=|l(x)|$ et $|\lambda|=1$. On a
        \begin{align*}
            \left|l(x)\right|&=\lambda l(x)\\
            &= l(\lambda x)\\
            &=\Re\circ l(\lambda x)\\
            &\le \norme{\Re\circ l}_{E'_\R}\norme{x}_E\\
            &\le \norme{\Re\circ l}_{E'_\R}
        \end{align*}
        Donc
        \begin{align*}
            \norme{l}_{E'_\C}&=\sup\{\left| l(x)\right|,x\in \overline{B_\epsilon}(0,1)\}\\
            &\le \norme{\Re\circ l}_{E'_\R}
        \end{align*}
        Finalement $\norme{\Re\circ l}_{E'_\R}=\norme{l}_{E'_\C}$. Donc $\Re$ est isométrique. En particulier, $\Re$ est injective. Soit $f \in E'_\R$ et $\funto{g_f}{x\in E}{f(x)-i f(ix)}$
    \end{proof}

    \subsection*{Hyperplans}

    \begin{prop}{}{}
        Soit $E$ un $\K$-\ev et $H$ un \sev de $E$. Sont équivalents:
        \begin{enumerate}
            \item $\exists e\in E,e\neq 0$ tel que $E=\K e\oplus H$
            \item $E\neq H$ et $\forall e\in E\wt H$, $E=\K e\oplus H$
            \item $\exists f\in E^*$ non nulle telle que $H=\ker f$.
        \end{enumerate}
        Un hyperplan de $E$ est un \sev de $E$ qui satisfait ces propositions
    \end{prop}

    \begin{prop}{}{}
        Soit $E$ un \evn. Un hyperplan $H$ de $E$ est fermé \ssi il existe $f\in E'=L(E,\K)$ non nulle telle que $H=\ker f$ 
    \end{prop}

    \begin{proof}
        Si $f\in E'$ alors $f$ est continue, $\{0\}$ est fermé dans $\K$ donc $\ker f\in f^{-1}(\{0\})$ est fermé dans $E$.\\
        Inversement, soit $\Hr$ un hyperplan fermé de $E, f\in E^*$ non-nulle telle que $\Hr=\ker f$ et $e\in E\wt \Hr$ tel que $f(e)=1$.\\
        Alors $E=\Hr\oplus \K e$ et pour $y\in \Hr$ et $\lambda \in \K$, $f(y+\lambda e)=\lambda$. Montrons que $f$ est continue. Soit $\delta=d(e,\Hr)=\inf\{\norme{e-y}_E,y\in \Hr\}$. Puisque $\Hr$ est fermée et $e\notin \Hr,\delta>0$. Soit $x\in E$:\\
        Si $x\notin \Hr$, soit $\lambda\in \K^*$ et $y\in \Hr,x=y+\lambda e$
        \[\norme{x}_E=\norme{y+\lambda e}_E =\left|\lambda\right|\norme{e-\left(\frac{-1}{\lambda}y\right)}_E\ge \left|\lambda\right|\delta=\left|f(x)\right|\delta\]
    \end{proof}

    \begin{dfprop}{}{}
        Soit $E$ un $\K$-\ev et $\Hr\subset E$. $\Hr$ est un hyperplan affine s'il satisfait une des condtions équivalentes suivantes:
        \begin{enumerate} 
            \item $\exists H$ un hyperplan vectoriel et $x\in E$,
            \[\Hr=x+H=\{x+y,y\in H\}\]
            \item $\exists f\in E^*$ et $\alpha\in \K$ telle que: $\Hr=\{z\in E, f(z)=\alpha\}=f^{-1}(\alpha)$
        \end{enumerate}
        Si $E$ est un \evn et $\Hr$ un hyperplan affine de $E$ alors:
        \begin{enumerate}
            \item Si $\Hr$ n'est pas fermé dans $E$, alors il est dense.
            \item $\Hr$ est fermé $\iff \exists f\in E',\alpha\in \K, \Hr=f^{-1}(\alpha)$
        \end{enumerate}
    \end{dfprop}

    \begin{proof}
        Laisser en exercice
    \end{proof}

    \section{Théorème de Hahn-Banach}

    \subsection*{Les formes géométriques}
    
    Dans cette partie $E$ est un espace vectoriel normé sur $\R$.

    \begin{lm}{Connexité et hyperplan}{}
        Soit $H$ un \sev fermé de $E$. La partie $E\wt H$ est connexe \ssi $H$ n'est pas un hyperplan et dans ce cas $E\wt H$ est connexe par arc.
    \end{lm}

    \begin{proof}
        Si $H$ est un hyperplan fermé, soit $f\in E'$ une forme linéaire continue de noyau $H$. Comme $f(E\wt H)=\R^*$ n'est pas connexe, il en résulte que $E\wt H$ n'est pas connexe.\\
        Inversement, supposons que $H$ n'est pas un hyperplan fermé et soient $x,y$ dans $E\wt H$.
        \begin{itemize}
            \item Si $y\not\in \R x\oplus H$, le segment $[x,y]=\{(1-t)x+ty,t\in [0,1]\}$ est un arc qui joint $x$ à $y$ dans $E\wt H$.
            \item Si $y\in \R x\oplus H$, puisque $y\not \in H$, il existe $\lambda \in \R^+$ et $h\in H$ tel que $y=\lambda x+ h$ et donc $\R x\oplus H= \R y\oplus H$. Comme $H$ n'est pas un hyperplan fermé, il existe $y\in E\wt \R x\oplus H$.\\
            la ligne brisée fermée des segments $[x,z]$ et $[z,y]$ dans $E\wt H$ et joint $x$ à $y$.
        \end{itemize}
    \end{proof}

    \begin{tm}{Théorème de Hahn-Banach géométrique}{}
        Soit $C\subset E$ un ouvert convexe non-vide avec $0_E\not\in C$. Il existe un hyperplan (vectoriel) fermé $H$ de $E$ tel que $C\cap H=\emptyset$.
    \end{tm}

    \begin{ra}{}{}
        $A\subset E$ est convexe $\iff \\forall x,y\in A,[x,y]=\{tx+(1-t)y,t\in [0,1]\}\subset A$
    \end{ra}

    \begin{re}{}{}
        Le théorème de $Hahn-Banach$ géométrique dit que l'on peut séparé $E$ en deux parties par un hyperplan fermé. $E=\R e \oplus H=E^+\Cup H\Cup E^-$ où $E^+=\{\lambda e+h,\lambda\in \R_+^*\text{ et }h\in H\}$ et $E^-=\{\lambda e+h,\lambda\in \R_-^*,h\in H\}$ et $C\subset E^+$ ou $C\subset E^-$.
    \end{re}

    \begin{lm}{Point clé du théorème de Hahn-Banach géométrique}{}
        Soit $C\subset E$ un ouvert convexe non-vide, avec $0_E\not\in C$ et $F$ un \sev de $E$ tel que $F\cap C=\emptyset$.\\
        Si $F$ n'est pas un hyperplan fermé, il existe un \sev $G$ de $E$ tel que $G\cap C=\emptyset$ et $F\subsetneq G$.
    \end{lm}

    \begin{re}{}{}
        Notons $\Fr$ l'ensemble des \sev $F$ de $E$ tel que $F\cap C=\emptyset$, et on a $\{0\}\in \Fr$. Montrer le théorème de Hahn-Banach revient à montrer que $\Fr$ contient un hyperplan.\\
        L'ensemble $\Fr$ est ordonné par l'inclusion. Le lemme dit précisemment que pour que $F$ soit un élément maximal de $\Fr$ (i.e. $F\in \Fr$ est maxiaml si $\forall G\in \Fr$, $F\subset G\implies F=G$). Il est nécessaire que $F$ soit un hyperplan fermé. Une fois le lemme montré, il n'y aura plus qu'à montrer que $\Fr$ admet (au moins) un élément maximal.
    \end{re}

    \begin{proof}
        Preuve du lemme: Puisque $C$ est ouvert, si $F\cap C=\emptyset$, alors $\overline{F}\cap C=\emptyset$.\\
        En effet, si $x\in \overline{F}\cap C$, alors $C$ est un voisinage ouvert de $x$ et puisque $x\in \overline{F}$, $C\cap F\neq \emptyset$.\\
        Donc si $F$ n'est pas fermé, il suffit de prendre $G=\overline{F}$ et le lemme est montré. Supposons $F$ fermé et que $F$ n'est pas un hyperplan. En particulier; $E\wt F$ est connexe. On cherche $G$ sous la forme $G=F\oplus \R^x$ où $x\not\in F$ et tel que $(F\oplus \R x)\cap C=\emptyset$. Notons $C+F=\{c+y,c\in C,y\in F\}$. On a:
        \[x\not\in F\qquad (F\oplus \R x)\cap C=\emptyset\iff x\not\in F\text{ et }\forall t\in \R, tx\not \in C+F(*)\]
        On va montrer que l'ensembles des $x\in E\wt F$ qui ne satisfont pas les propriétés équivalentes $(*)$, c'est à dire:
        \[A=\{x\in E\wt F,\exists t\in \R, tx\in C+F\}\neq E\wt F\]
        Puisque $C\cap F=\emptyset, 0 x=0_E\oplus C+F$ donc $A=\{x\in E\wt F,\exists t\in \R^*, tx\in C+F\}$. Soit
        \begin{align*}
            W_+&=\{x\not \in F, \exists t>0,tx\in C+F\}\\
            &=\{sc+y,s\in \R_+^*,c\in C, y\in F\}\\
            &=\bigcup_{s\in \R_+^*,y\in F} (sC+y)
        \end{align*}
        où $sC+y=\{sc+y,c\in C\}=\tau_y\cdot n_s(C)$ est un ouvert donc $W_+$ est un ouvert. De même,
        \begin{align*}
            W_-&= \{x\not\in F,\exists t<0,tx\in C+F\}\\
            &=\{sc+y,s\in \R_-^*,c\in C,y\in F\}=\bigcup_{s\in \R_-^*,y\in F}(sC+y)
        \end{align*}
        On a deux ouverts $W_+$ et $W_-$ contenus dans $E\wt F$ qui sont tels que $A=W_+\cup W_-$.\\
        De plus, $W_+\cup W_-=\emptyset$. Si $s,t\in \R_+^*,b,c\in C$ et $y,z\in F$ sont tels que $sb+y=-(tc+z)\in W_+\cup W_-$
        \[\frac{s}{s+t}b+\frac{t}{s+t}c=\frac{-1}{s+t}(y+z)\in C\cap F \text{(Exclu par hypothèse)}\]
        Puisque $E\wt F$ est connexe, il n'admet pas de partition par deux ouverts non vides, donc $W_+\cup W_-\subsetneq E\wt F$, il existe donc $x\in E\wt F$ et tel que $tx\notin C+F$ pour tout $t\in \R$
    \end{proof}

    \begin{lm}{Lemme de Zorn}{}
        Tout ensemble inductif admet nécessairement un élément maximal.
    \end{lm}

    \begin{proof}
        Preuve du théorème de Hahn-Banach géométrique:
        \begin{enumerate}
            \item Cas où $\dim E$ est finie. Soit $\Hr\in \Fr$ de dimension maximale. D'après le lemme précédent c'est un hyperplan fermé.
            \item Cas général. L'ensemble $\Fr$ est inductif pour l'inclusion
            \begin{itemize}
                \item $\{0_E\}\in \Fr$ et donc $\Fr$ est non vide.
                \item Si $T$ est un sous-ensemble non-vide totalement ordonné pour l'inclusion de \sevs de $E$ ne contenant pas $C$, alors $\bigcup_{F\in T}F$ est un majorant de $T$ dans $\Fr$.\\
                "Tout sous-ensemble totalement ordonné de $\Fr$ possède un majorant". Donc $\Fr$ admet un élément maximal qui d'après le lemme précédent est un hyperflan.
            \end{itemize}
        \end{enumerate}
    \end{proof}

    \begin{cor}{Théorème de Hahn-Banach géomégtrique avec les formes linéaires continues}{}
        Soit $C\subset E$ un ouvert convexe non-vide, $0_E\notin C$. Il existe $f\in E'$ telle que $\forall x\in C,f(x)>0$.
    \end{cor}

    \begin{proof}
        Par le théorème de Hahn-Banach, il existe $\Hr$ hyperplan fermé de $E$ tel que $\Hr\cap C=\emptyset$. Soit $f\in E'$ tel que $\Hr=\ker f$.\\
        L'image d'un convexe par une application affine est un convexe. Donc $f(X)$ est un convexe de $\R$ et les convexes de $\R$ sont les intervalles.\\
        Puisque $C\cap \Hr=\emptyset, 0\notin f(C)$. Finalement $f(C)\subset \R^*_+$ ou $f(C)\subset \R_-^*$. Quitte à changer $f$ en $-f$, on a le résultat.
    \end{proof}

    \begin{cor}{Théorème de séparation de Hahn-Banach version fermée}{}
        Soient $A$ et $B$ deux ensembles convexes disjoints non-vides de $E$, avec $A$ fermé et $B$ compact.\\
        Il existe une forme linéaire continue $f\in E'$ et $a,b$ dans $\R$ tels que:
        \[\forall x\in A,\forall y\in B,f(x)\le a<b\le f(y)\]
        Ou de façon équivalente, si $c\in \R$ est tel que $a<c<b$, l'hyperplan affine $\Hr=f^{-1}({c})$ sépare $A$ et $B$
    \end{cor}

    \begin{proof}
        Exercice
    \end{proof}

    \begin{cor}{Théorème de séparation de Hahn-Banach version ouverts}{}
        Soit $A$ et $B$ deux ensembles convexes disjoints non-vides de $E$ avec $A$ ouvert.\\
        Il existe une forme linéaire $f\in E'$ et $a\in \R$ tel que
        \[\forall x\in A,f(x)>a\text{ et }\forall y\in B, f(y)\le a\]
        De façon équivalente, l'hyperplan affine fermé $f^{-1}\left(\{a\}\right)$ sépare $A$ et $B$
    \end{cor}

    \begin{proof}
        Pour tout $y\in B$, notons $V_y=\{x-y,x\in A\}=\tau_y(A)$, c'est un ouvert.\\
        On considère
        \[C=A-B=\{x-y,x\in A,y\in B\}=\bigcup_{y\in B}V_y\]
        Alors
        \begin{itemize}
            \item $0_E\notin C$ car $A\cap B=\emptyset$
            \item $C$ est ouvert comme réunion d'ouverts
            \item $C$ est convexe: pour $x_1,x_2\in A$ et $y_1,y_2\in B$ et $t\in [0,1]$
            \[t(x_1-y_1)+(1-t)(x_2-y_2)=tx_1+(1-t)x_2 -(ty_1+(1-t)y_2)\in A-B=C\]
        \end{itemize}
        Par le corollaire précédent, il existe $f\in E'$ tel que $f(C)\subset \R_+^*$. Donc $\forall x\in A,\forall y\in B, f(y)<f(x)$.\\
        Ainsi $f(A)$ et $f(B)$ sont deux intervalles de $\R$ disjoints tels que $\forall \alpha \in f(A),\forall \beta \in f(B), \beta<\alpha$. En particulier $f(A)$ est minoré et $f(B)$ est majoré et $\sup(f(B))\le \inf (f(A))$.\\
        Soit $a=\sup(f(B))$, on a $\forall y\in B,f(y)\le a$. De plus, $A\subset \{x\in E,f(x)\ge a\}=D$ et puisque $A$ est ouvert
        \[\mathring{A}=A\subset \mathring{D}=\{x\in E,f(x)>a\}\] 
    \end{proof}

    \subsection*{Les formes analytiques}

    On suppose ici que $E$ est un $\K$-espace vectoriel normé, $\K=\R$ ou $\K=\C$

    \begin{tm}{Théorème de prolongement de Hahn-Banach}{}
        Soit $E$ un $\K$-\evn et $F$ un \sev de $E$.\\
        Pour tout $l\in F'$, il existe $\tilde{l}\in E'$ qui prolonge $l$ et telle que $\norme{\tilde{l}}_{E'}=\norme{l}_{F}$
    \end{tm}

    \begin{proof}
        Si $l$ est nulle, c'est trivial. De plus, on a vu que si $\K=\C$, les duaux réels (i.e. où $E$ est considéré comme un $\R$-\ev) et complexes sont isométriques. Il suffit donc de traiter le cas réel.\\
        On suppose donc que $\K=\R$ et que $l$ n'est pas réelle. Soit $A=B_E(0_E,1)$ la boule ouverte de $E$ et $B=\{x\in F, l(x)=\norme{l}_{F'}\}$. On a:
        \begin{itemize}
            \item $A$ est un convexe ouvert non-vide de $E$.
            \item Puisque $l\neq 0_{F'}$, elle est surjective, $B$ est non-vide et $0_E\notin B$. Pour $x_0\in B$, elle est surjective, $B$ est non-vide et $0_E\notin B$. Pour $x_0\in B,B=l^{-1}\left( \norme{l}_{F'}\right)=\{x_0+y,y\in \ker l\}$ est un hyperplan affine de $F$ et en particulier convexe. De plus, $F=\ker l\oplus x_0 \R=\text{Vect}(B)$
            \item Pour $x\in A\cap B$, on a $\left|l(x)\right|\le \norme{l}_{F'}\norme{x}_E<\norme{l}_{F'}$. Donc $A\cap B=\emptyset$.\\
        \end{itemize}
        On peut appliquer le théorème de séparation de Hahn-Banach version ouverte qui assure qu'il existe $f\in E'$ et $a\in \R$ tel que $f(A)\subset ]-\infty,a[$ et $f(B)\subset [a,+\infty[$.\\
        Puisque $0_E\in A$ et $f(0_E)=0,a>0$.\\
        Puisque $B$ est un sous-espace affine de $E$, son image par une application linéaire est un sous-espace affine et donc $f(B)$ est un sous-espace affine sur $\R$. Puisque ce n'est pas $\R$, c'est un point.\\
        Il existe donc $b\ge a>0$ tel que $f(B)=\{0\}$. Quitte à multiplier $f$ par $b^{-1}\norme{l}_\infty$, on peut supposer $b=\norme{l}_{F'}$. Alors $l$ et $f$ coïncident sur $B$, qui est un hyperplan affine ne contenant pas $0_E$. Mais $B$ engendre $F$, donc $l$ et $f$ coïncident sur $F$. En particulier $\norme{l}_{F'}\le \norme{f}_{E'}$.\\
        De plus
        \begin{align*}
            \norme{f}_{E'}&=\sup\{\left|f(x)\right|,x\in A\}\\
            &\underset{(*)}{=}\sup \{f(x),x\in A\}\\
            &\le a \le \norme{l}_{F'}
        \end{align*}
        D'où l'égalité ($(*)$ est vraie car si $f(x)\le 0$ alors $f(-x)=-f(x)\ge 0$ et puisque $A$ est une boule ouverte centrée en $0_E$, si $x\in A,-x\in A$)
    \end{proof}

    \begin{cor}{Corollaire du théorème de prolongement de Hahn Banach}{}
        Soit $E$ un $\K$-\evn:
        \begin{enumerate}
            \item Pour tout $x\in E,x\neq 0_E$, il existe $f\in E'$ telle que $\norme{f}_{E'}=1$ et $f(x)=\norme{x}_E$.
            \item Pour tout $x\in E,\norme{x}_E=\sup_{\norme{f}_{E'}=1}|f(x)|$
            \item Soit $G$ un \sev fermé de $E$, $F\neq E$. Il existe $f\in E'$ telle que $\norme{f}_{E'}=1$ et $f(x)=0$ pour tout $x\in G$.
            \item Soit $G$ un \sev de $E$ alors $\overline{G}=E$ \ssi pour tout $f\in E'$, si $f_{|G}=0$, alors $f=0$
        \end{enumerate}
    \end{cor}

    \begin{re}{}{}
        En particulier, (1) montre que si $E\neq \{0\}$, alors $E'\neq \{0\}$.
    \end{re}

    \begin{proof}
        Laisser en exercice\\
        Indications: 
        \begin{enumerate}
            \item Il suffit d'appliquer le théorème de prolongement de Hahn-Banach à $F=\K x$ et $ l(\lambda x)=\lambda \norme{x}_E$
            \item Découle de $(1)$
            \item Soit $x_0\notin G,F=G\oplus x_0 \R$ et $\funto{l}{F}{\K}$ donc par $l(y+\lambda x_0)=\lambda$. On vérifie que $l\in F'$ et $l$ est nulle sur $F$.
            \item Découle de $(3)$
        \end{enumerate}
    \end{proof}

    \subsection*{Bidual}

    Si $E$ est un $\K$-\evn, $E'=L(E,\K)$ son dual topologique est lui-même un $\K$-\evn et même un Banach.\\
    Le dual topologique de $E'$, que l'on note $E''$ est appelé le bidual topologique de $E$.

    \begin{cor}{Injection de $E$ dans un bidual}{}
        L'espace $E$ s'injecte isométriquement dans son bidual.\\
        L'application $\fundfto{J}{E}{E''}{x}{\fundfto{J(x)}{E'}{\K}{l}{l(x)}}$ est une isométrie.
    \end{cor}

    \begin{proof}
        Pour $x\in E, \fundfto{J(x)}{E'}{\K}{l}{l(x)}$ est bien linéaire et pour $l\in E'$
        \[\left|J(x)(l)\right|=\left|l(x)\right|\le \norme{l}_{E'}\norme{x}_E\]
        Donc $J(x)$ est continue. $J(x)\in L(E',\K)$ et $\norme{J(x)}_{E'}\le \norme{x}_{E}$.
        \begin{itemize}
            \item Si $x=0_E$, on a $J(x)=0_{E'}$ et donc $\norme{J(0_E)}_{E''}=\norme{0_E}_E=0$
            \item Si $x\neq 0_E$, par le théorème de prolongement de Hahn-Banach, il existe $f\in E'$ telle que $\norme{f}_{E'}=1$ et $f(x)=\norme{x}_E$. Alors:
            \[\left|J(x)(l)\right|=\left|f(x)\right|=\norme{x}_E\]
            Donc $\norme{J(x)}_{E''}=\norme{x}_E$. $J$ est évidemment linéaire, elle est continue de norme $1$ et c'est une isométrie (En particulier $J$ est injective).
        \end{itemize}
    \end{proof}

    \begin{df}{}{}
        Un \evn $E$ est dit reflexif si l'injection $\funto{J}{E}{E''}$ est surjective et donc bijective. Dans ce cas $E$ et $E''$ sont isométriques et $E$ est nécessairement un espace de Banach (puisque $E'$ est un Banach)
    \end{df}

    \begin{ex}{}{}
        \begin{itemize}
            \item Si $\dim E$ est finie, $E$ est reflexif
            \item Les espaces de Hilbert, les espaces $l^p$ et $L^p$ pour $1<p<+\infty$ sont reflexifs (Voir dans ce chapitre)
        \end{itemize}
    \end{ex}

    \section{Convergence faible et faible*}
    
    Dans cette section, $E$ est un $\K$-\evn où $\K=\R$ ou $\C$

    \subsection*{Convergence faible}

    \begin{df}{}{}
        Soit $(x_n)_\nN$ une suite de $E$. On dit que $(x_n)_\nN$ converge faiblement vers $x\in E$ Si
        \[\forall l\in E',l(x_n)\underset{n\to +\infty}{\longrightarrow}l(x)\]
        On note $x_n \underset{n\to \infty}{\rightharpoonup} x$
    \end{df}

    \begin{re}
        Il y a unicité de la limite faible quand elle existe. En effet, si $l(x-y)=0,\forall l\in E'$, par le $(2)$ du corollaire du théorème de prolongement de Hahn Banach, $\norme{x-y}_E=0$
    \end{re}

    \begin{re}{}{}
        On parle parfois de convergence forte pour parler de la convergence usuelle, c'est à dire au sens de la norme $\norme{x_n-x}_E\underset{n\to \infty}{\longrightarrow} 0$.
    \end{re}

    \begin{prop}{}{}
        La convergence forte implique la convergence faible. Pour $(x_n)_\nN$ une suite de $E$ et $x\in E$
        \[\norme{x_n-x}_E \underset{n\to \infty}{\longrightarrow} 0\implies x_n \underset{n\to \infty}{\rightharpoonup} x\]
    \end{prop}

    \begin{proof}
        Si $x_n \to x$ dans $E$, pour tout $l\in E$, par continuité de $l$, $l(x_n)\to l(x)$ dans $\K$. Donc $x_n \underset{n\to \infty}{\rightharpoonup}x$.
    \end{proof}

    \begin{re}{}{}
        La topologie faible sur $E$, noté $\sigma(E,E')$ est la topologie initiale associée à la famille de toutes les fonctions linéaires continues sur $E$: c'est la topologie la moins fine pour laquelle les éléments de $E'$ sont encore continues. Elle est engendrée par les ouverts de la forme $l^{-1}(U)$ où $U$ est un ouvert de $\K$ et $l\in E'$.\\
        Une suite de $E$ converge faiblement \ssi elle converge pour la topologie faible.\\
        La topologie faible sur $E$ est la topologie associée à la famille de semi-norme $|l(\cdot)|$ pour $l\in E'$, elle n'est jamais métrisable.
    \end{re}

    \begin{prop}{}{}
        Si $E$ est de dimension finie, la convergence forte et la convergence faible coïncident.
    \end{prop}
    
    \begin{proof}
        Il faut montrer que si $(x_n)_\nN$ est une suite de $E,x\in E$ et $x_n\underset{n\to \infty}{\longrightarrow} x$ alors $\norme{x_n-x}_E\underset{n\to \infty}{\longrightarrow} 0$. Soit $(e_1,\ldots,e_d)$ une base de $E$ et pour tout $i\in \llbracket 1,d\rrbracket$. Soit $\funto{e_i^*}{E}{\K}$, telle que $e_i^*(\sum_{k=1}^{d}a_k e_k)=a_i$.\\
        Alors $e_i^*\in E'$. Donc $e_i^*(x_n)\underset{n\to \infty}{\longrightarrow} e_i^*(x)$ pour tout $i\in \llbracket 1,d \rrbracket$. On en déduit
        \[x_n=\sum_{k=1}^{d}e_i^*(x_n)e_i\underset{n\to \infty}{\longrightarrow} \sum_{i=1}^{d}e_i^*(x)e_i=x\]
    \end{proof}

    \begin{re}{}{}
        La réciproque est fausse en dimension infinie. Par exemple pour $E=C^1([0,1],\R)$ muni de la norme $\norme{\cdot}_{2}$, la suite $(c_n)_\nN$ donnée par $c_n(t)=\cos(nt)$ converge faiblement vers $0$ mais sa norme ne converge pas vers $0$.
    \end{re}

    \begin{prop}{}{}
        Soit $T\in L(E,F)$, et $(x_n)_\nN$ une suite de $E$ qui converge faiblement vers $x$. Alors $(T(x_n))_\nN$ converge faiblement vers $T(x)$.
    \end{prop}

    \begin{proof}
        Soit $l\in F'$ alors $l\circ T\in E'$ donc
        \[l\circ T(x_n)=l(T(x_n))\underset{n\to\infty}{\longrightarrow}l(T(x))\]
    \end{proof}

    \begin{prop}{}{}
        Soit $(x_n)_\nN$ une suite de $E$ qui converge faiblement vers $x$. Alors la suite $(x_n)_\nN$ est (fortement) bornée et:
        \[\norme{x}_E \le \liminf \norme{x_n}_E\]
    \end{prop}

    \begin{proof}
        Pour $E=\{0_E\}$ c'est vrai. On suppose $E\neq \{0_E\}$. Soit l'injection canonique qui est isométrique:
        \[\fundfto{J}{E}{E''=(E')'}{x}{\fundfto{J(x)}{E'}{\K}{l}{l(x)}}\]
        On note $T=J(x)$ et pour tout $n\in \N$, $T_n=J(x_n)$, ce sont des applications linéaires continues de $E'$ dans $\K$ qui vérifient:
        \[\norme{x}_E=\norme{T}_{E''}\text{ et }\norme{x_n}_E =\norme{T_n}_{E''}\]
        La condition $x_n\rightharpoonup x$ dit précisemment que la suite $(T_n)_\nN$ converge simplement vers $T$.\\
        Pour $l\in B_{E'}(0_{E'},1)$, on a donc:
        \[\left|T(l)\right|=\lim_{n\to \infty}\left|T_n(l)\right|\liminf_{n\to \infty} \left|T_n(l)\right|\overset{(1)}{\le} \liminf_{n\to \infty} \norme{T_n}_{E''}\norme{l}_{E'}\le \liminf_{n\to\infty} \norme{T_n}_{E''}\]
    \end{proof}

    \subsection*{Convergence faible *}

    \begin{df}{}{}
        Soit $(l_n)_\nN$ une suite de $E'$. On dit que $(l_n)_\nN$ converge faiblement vers $l\in E$ si
        \[\forall x\in E, l_n(x)\underset{n\to \infty}{\longrightarrow} l(x)\]
        On note $l_n \underset{n\to \infty}{\rightharpoonup} l$
    \end{df}

    \begin{re}{}{}
        Il s'agit juste de la convergence simple.
    \end{re}

    \begin{prop}{}{}
        \begin{enumerate}
            \item Si une suite $(l_n)_\nN$ de $E'$ converge fortement vers $l\in E'$ (i.e. $\norme{l_n-l}_{E'}\to 0$) alors $(l_n)_\nN$ converge faible * vers $l$.
            \item Soit $E$ un espace de Banach. Si $l_n$(VOIR POUR LA FLECHE) $l$ alors la suite $(l_n)_\nN$ est (fortement) bornée et $\norme{l}_{E'}\le \liminf \norme{l_n}_{E'}.$        
        \end{enumerate}
    \end{prop}

    \begin{proof}
        Pour la $(1)$:
        \[\forall x\in E,\left|l_n(x)-l(x)\right|=\left|(l_n-l)(x)\right|\le \norme{l_n-l}_{E'}\cdot \norme{x}_E\]
        Le théorème de Banach-Steinhaus nous assure que la suite $(l_n)_\nN$ est bornée.\\
        En effet, puisque pour tout $x\in E, (l_n(x))_\nN$ converge en particulier $\{l_n(x),n\in \N\}$ est borné dans $\K$ donc $\{l_n,n\in\N\}$ est borné dans $E'$. Ensuite on procède comme précedemment. Pour $x\in B_E(0_E,1)$, on a:
        \begin{align*}
            \left|l(x)\right|=\lim_{n\to \infty}\left|l_n(x)\right|=\liminf_{n\to \infty} \left|l_n(x)\right|&\le \liminf_{n\to \infty}\norme{l_n}_{E'}\norme{x}_E\\
            &\le \liminf_{n\to \infty} \norme{l_n}_{E'}
        \end{align*}
        Le passage au sup permet de conclure.
    \end{proof}

    \begin{tm}{Théorème de Banach-Alaoglu}{}
        Soit $E$ un espace vectoriel normé séparable.\\
        Toute suite bornée de $E'$ admet une sous-suite qui converge faible * dans $E'$.
    \end{tm}

    \begin{re}{}{}
        On parle de compacité faible *: La boule unité fermée de $E^*$ est *-faiblement compacte.
    \end{re}

    \begin{lm}{}{}
        Soit $E$ un \evn et $D$ une partie dense de $E$.\\
        Soit $(l_n)_\nN$ une suite bornée de $E'$ qui converge simplement (=faible *) vers $l\in E'$ sur $D$. \\
        Alors la suite $(l_n)_\nN$ converge simplement vers $l$ sur $E$
    \end{lm}

    \begin{proof}
        Soit $x\in E$. Par densité de $D$ dans $E$, pour $\epsilon>0$, il existe $y\in D, \norme{x-y_\epsilon}_E\le \epsilon$.\\
        Soit $N\in \N$ tel que $\forall n\in \N,n\ge N, \left|l_n(y_\epsilon)-l(y_\epsilon)\right|\le \epsilon$.\\
        Pour $n\ge N$:\begin{align*}
            \left|l_n(x)-l(x)\right|&\le \left|l_n(x)-l_n(y_\epsilon)\right|+\left|l_n(y_\epsilon)-l(y_\epsilon)\right|+\left|l(y_\epsilon)-l(x)\right|\\
            &\le \norme{l_n}_{E'}\norme{x-y_\epsilon}_E+\epsilon+\norme{l}_{E'}\norme{y_\epsilon-x}_{E}\\
            &\le \left(\norme{l_n}_{E'}+1+\norme{l}_{E'}\right)\epsilon\le C\epsilon
        \end{align*}
        où $C=\sup_{\nN}\norme{l_n}_{E'}+1+\norme{l}_{E'}$
    \end{proof}

    \begin{proof}
        Preuve du théorème de Banach-Alaoglu: Soit $D=\{x_k,k\in \N\}$ une partie dénombrable et dense de $E$. Soit $(l_n)_\nN$ une suite bornée de $E'$. En particulier: Pour tout $k\in \N, (l_n(x_k))_\nN$ est une suite bornée de $\K$.\\
        Il existe $M_k>0$ tel que $\{l_n(x_k),n\in \N\}\subset \overline{B_\K}(0,M_k)$ et $\overline{B_k}(0,M_k)$ est un compact de $\K$. D'après le théorème de Tikhinov métrique dénombrable, le produit
        \[K=\prod_{k\in N} \overline{B_k}(0,M_k)\text{ est compact}\]
        La suite $(a_n)_{\nN}$ de $K$ où pour tout $n\in\N$, $a_n=(l_n(x_k))_{k\in \N}$ admet une sous-suite convergente $(a_{\varphi(n)})_\nN$. Alors la suite extraite $(l_{\varphi(n)})_\nN$ est simplement convergente sur $D$ vers une application $l$ définit sur $D$.\\
        Par linéarité, $l$ s'étend en une forme linéaire, encore notée $l$ sur $F=\text{Vect}(x_k,k\in \N)$.\\
        Par continuité de la somme et de la multiplication extérieure, la suite $(l_n)_\nN$ converge simplement vers $l$ sur $F$. De plus:
        \[\forall x\in F, n\in \N:\left|l_{\varphi(n)}(x)\right|\le \norme{l_{\varphi(n)}}_{E'}\norme{x}_{E}\le \sup_{n\in \N}\norme{l_n}_{E}\norme{x}_E\]
        Donc
        \[\left|l(x)\right|=\lim_{n\to \infty}\left|l_{\varphi(n)}(x)\right|\le \sup_{\nN}\norme{l_n}_{E'}\norme{x}_E\]
        Donc $l$ est continue sur $F$ qui est un \sev dense de $E$. Par le théorème de prolongement des applications linéaires continues, $l$ peut être étendue à $\tilde{l}\in E'$. Le lemme précédent permet de conclure.
    \end{proof}

    \section{Application de la dualité}

    Dans le cas des \evn, on s'interesse à la dualité induite par la forme bilinéaire canonique:
    \[\fundfto{B}{E'\times E}{\K}{(l,x)}{l(x)}\]
    On note parfois $B(l,x)=\sca{l,x}$\\
    Remarquons:\begin{itemize}
        \item Pour $x\in E$, si $\forall l\in E',l(x)=0$ d'après un corollaire du Théorème de prolongement de Hahn-Banach, $x=0$
        \item Pour $l\in E'$, si $\forall x\in E$, $l(x)=0$ alors $l=0_{E'}$. On dit que $B$ est non-dégénérée. De plus,
        \[ \forall l\in E',\forall x\in E |l(x)|\le \norme{l}_{E'} \norme{x}_E\]
        Donc $B$ est continue et de norme $1$ par le corollaire du théorème de Hahn-Banach
    \end{itemize}

    \subsection*{Les espaces de Hilbert}

    Dans la suite, $\Hr$ est un espace de Hilbert sur $\K=\R$ ou $\C$, c'est à dire que $\Hr$ est un espace de Banach sur $\K$ dont la norme $\norme{\cdot}_{\Hr}$ provient d'un produit scalaire $\sca{\cdot,\cdot}$ (réel ou hermitien).\\
    On peut penser $l^2(\N,\K)$ et $L^2$

    \begin{prop}{}{}
        Pour $y\in \Hr$, la forme linéaire $\funto{l_y}{x\in H}{\sca{x,y}\in \K}$ est continue.\\
        L'application $\fundfto{\tau}{\Hr}{\Hr'}{y}{l_y}$ est antilinéaire et isométrique de $\Hr$ dans $\Hr'$.
    \end{prop}

    \begin{proof}
        Pour $x,y$ dans $\Hr$, on a par l'inégalité de Cauchy-Schwarz
        \[|l_y(x)|=|\sca{x,y}|\le \norme{x}_{\Hr}\norme{y}_{\Hr}\]
        On a évidemment que $l_y$ est linéaire. Donc $l_y\in \Hr'$ et $\norme{l_y}_{\Hr'}\norme{y}_{\Hr}$. De plus:
        \[\norme{y}_\Hr^2=\sca{y,y}=l_y(y)\le \norme{l_y}_{\Hr'}\norme{y}_\Hr\]
        D'où $\norme{y}_\Hr\le \norme{l_y}_{\Hr'}$ et donc $\norme{y}_\Hr=\norme{l_y}_{\Hr'}$\\
        Pour $y,z\in \Hr$ et $\lambda\in \C$, on a:
        \begin{align*}
            \forall x\in \Hr, \tau(y+\lambda z)(x)&=\sca{x,y+\lambda z}=\sca{x,y}+ \overline{\lambda}\sca{x,z}\\
            &=\tau(y)(x)+\overline{\lambda}\tau(z)(x)\\
            &=(\tau(y)+\overline{\lambda}\tau(z))(x)
        \end{align*}
        Donc $\tau(y+\lambda z)=\tau(y)+\overline{\lambda}\tau (z)$. C'est à dire que $\tau$ est anti-linéaire.
    \end{proof}

    \begin{df}{}{}
        \begin{itemize}
            \item Pour $x,y\in \Hr$ $x$ est orthogonal à $y$ si $\sca{x,y}=0$
            \item Pour $A,B\in \Pr(\Hr)$, $A$ est orthogonale à $B$ si
            \[\forall x\in A,\forall y\in B, \sca{x,y}=0\]
            \item Pour $A\subset \Hr$, l'orthogonal de $A$ est
            \[A^{\perp}=\{x\in H,\forall y\in A, \sca{x,y}=0\}\]
        \end{itemize}
    \end{df}

    \begin{re}{}{}
        On a:
        \begin{align*}
            A^{\perp}&=\{x\in H,\forall y\in A,\sca{x,y}=0\}=\bigcap_{y\in A}\ker l_y\\
            &=\text{Vect}(A)^{\perp}\overset{(*)}{=}\overline{\text{Vect}}(A)^\perp=\overline{A}^\perp
        \end{align*}
        En particulier, $A^\perp$ est un \sev fermé de $\Hr$\\
        Pour l'égalité $(*)$, puisque $A\subset \overline{A}$, on a clairement $\overline{A}^\perp \subset A^\perp$. Soit $y\in A^\perp,x\in \overline{A}$ et $(x_n)_\nN$, une suite de $A$ qui converge vers $x$. On a:
        \[\left|\sca{x,y}\right|=\left|\sca{x-x_n,y}+\underset{=0}{\sca{x_n,y}}\right|\le \norme{x-x_n}_\Hr\norme{y}_\Hr\]
        Par passage à la limite quand $n$ tend vers $+\infty$, $\sca{x,y}=0$. On a donc $A^\perp \subset \overline{A}^\perp$. De plus, $p_F$ est linéaire et continue.
    \end{re}

    \begin{tm}{Théorème de projection}{}
        Soit $\Hr$ un espace de Hilbert et $C$ une partie convexe fermée de $\Hr$.
        \begin{enumerate}
            \item $\forall x\in \Hr,\exists ! x_0\in C, d(x,C)=\norme{x-x_0}_\Hr$. On dit alors que $x_0$ est la projection orthogonale de $x$ sur $C$
            \item Le point $x_0$ est caractérisé par:
            \[ x_0\in C\text{ et }\Re(\sca{x_0-x,x_0-y})\le 0,\forall y\in \Cr\]
        \end{enumerate}
    \end{tm}

    \begin{prop}{}{}
        Soit $\Hr$ un espace de Hilbert, $F$ un \sev fermé de $\Hr$. On a
        \[\Hr=F\oplus F^{\perp}\]
        La projection orhogonale $\funto{p_F}{\Hr}{F}$ (Parallèlement à $F$) est la projection orthogonale au sens du théorème de projection.\\
        De plus, $p_F$ est linéaire et continue
    \end{prop}

    \begin{re}{}{}
        Pour $x\in \Hr$, on a:
        \[x=p_F(x)\iff x\in F\text{ et }x-x_0\in F^\perp\iff d(x,F)=\norme{x-x_0}_\Hr\]
    \end{re}

    \begin{re}{}{}
        On en déduit que si $\dim F$ est finie et $(e_1,\ldots,e_n)$ est une base orthonormée de $F$, pour tout $x\in \Hr$: 
        \[p_F(x)=\sum_{k=0}^{n}\sca{x,e_i}e_i\]
    \end{re}

    \begin{proof}
        Preuve de $p_F$ est continue: Pour tout $x\in F$ et $y\in F^\perp$, on a:
        \[\norme{x+y}_\Hr^2=\norme{x}_\Hr^2+\norme{y}_\Hr^2\ge \norme{x}_\Hr^2=\norme{p_F(x+y)}^2\]
        Donc $p_F\in L(\Hr,\Hr)$ et $\norme{p_F}_{L(\Hr,\Hr)}\le 1$. De plus pour $x\in F, p_F(x)=x$ donc $\norme{p_F(x)}_\Hr=\norme{x}_\Hr$ donc si $F=\{0_\Hr\},\norme{p_F}_{L(\Hr,\Hr)}=1$
    \end{proof}

    \begin{cor}{}{}
        Soit $H$ un espace de Hilbert
        \begin{enumerate}
            \item Si $A\subset H,(A^\perp)^\perp$ est le plus petit \sev fermé de $H$ qui contient $A$.
            \item Si $F$ est un \sev de $H,(F^\perp)^\perp=\overline{F}$
            \item Si $F$ est un \sev de $H$, $F$ est dense dans $H$ \ssi $F^\perp=\{0_H\}$
        \end{enumerate}
    \end{cor}

    \begin{tm}{Théorème de représentation de Riesz}{}
        Soit $H$ un espace de Hilbert. L'application isométrique antilinéaire
        \[\fundfto{\tau}{H}{H'}{y}{\left(\fundfto{l_y}{H}{\K}{x}{\sca{x,y}}\right)}\]
        est une bijection isométrique de $H$ sur $H'$.
    \end{tm}

    \begin{re}{}{}
        Puisque $H'$ est isométrique à $H$, la norme de $H'$ est issue d'un produit scalaire. On a $\sca{l_x,l_y}_{H'}=\sca{y,x}_H$. Donc $H'$ est un espace de Hilbert
    \end{re}

    \begin{proof}
        On sait déjà que $\tau$ est antilinéaire et isométrique. Donc pour tout $x\in H,\norme{\varphi_x}_{H'}=\norme{x}_H$ donc $\tau $ est injective. Soit $f\in H'$ et $F=\ker f$.\\
        Si $f\neq 0_{H'}, F\neq H$ et puisque $H=F\oplus F^\perp,F^\perp\neq \{0_H\}$.\\
        Soit $x\in F^\perp,x\neq 0_H$, on a:
        \begin{itemize}
            \item $F$ est un hyperplan et $x\notin F$ dans $H=F\oplus x\K$
            \item $l_x(x)=\norme{x}_H^2\neq 0$, soit $\lambda \in \K$ tel que $f(x)=\lambda l_x(x)$.
            \item Comme $x\in F^\perp=\{y\in H,\forall z\in F, \sca{y,z}=0\}$, la forme $l_x$ s'annule sur $F$.
        \end{itemize}
        Finalement, $f$ et $\lambda l_x$ coïncident sur $F$ (elles sont nulles) et en $x$, elles sont donc égales. D'où $f=\lambda l_x=\tau(\overline{\lambda}x)$.\\
        Donc $H'=\Im (\tau)$
    \end{proof}

    \begin{cor}{}{}
        Tout espace de Hilbert est réflexif. C'est à dire que
        \[\fundfto{J}{H}{H''}{x}{\fundfto{J(x)}{H'}{\K}{l}{l(x)}}\]
        est une bijection isométrique
    \end{cor}

    \begin{proof}
        On sait déjà que $J$ est isométrique. Soit $\varphi\in H''$. L'application $\fundf{H}{\K}{x}{\overline{\varphi(l_x)}}$ est une forme linéaire continue sur $H$. Par le théorème de Riesz, il existe un unique $y\in H$ tel que:
        \[\forall x\in H, \varphi(l_x)=\overline{l_y(x)}=\sca{y,x}=l_x(y)\]
        Donc:
        \[\forall f\in H'\text{ on a }\varphi(f)=f(y)\]
        Donc $\varphi=J(y)$. Donc $J$ est surjective
    \end{proof}

    \begin{cor}{Convergence faible et faible* dans un Hilbert}{}
        Soit $H$ un espace de Hilbert, $(x_n)_\nN$ une suite de $H$ et $x\in H$. Sont équivalents:
        \begin{enumerate}
            \item $(x_n)_\nN$ converge faiblement vers $x$.
            \item $(l_{x_n})_\nN$ converge faible* vers $l_x$.
            \item $\forall y\in H$, la suite $\left(\sca{x_n,y}\right)_\nN$ converge vers $\sca{x,y}$.
        \end{enumerate}
    \end{cor}

    \begin{proof}
        Laisser en exercice
    \end{proof}

    \begin{cor}{Corollaire de Riesz + Banach-Alaoglu}{}
        Toute suite bornée dans un espace de Hilbert admet une sous-suite qui converge faiblement.
    \end{cor}

    \subsection*{Les espaces $l^p$ et $L^p$}

    \subsubsection*{Inégalités de Hôlder et Minkowski}

    Soit $(X,\Ar,\mu)$ un espace mesuré et $p,q\in ]1,+\infty[$ conjugués:
    \[\frac{1}{q}+\frac{1}{p}=1\]
    Pour tout $\funto{f,g}{X}{\R_+}$ mesurables, on a: Inégalité de Hölder:
    \[\int_X fg\,d\mu \le \left(\int_X f^p \,d\mu\right)^{\frac{1}{p}}\left(\int_X g^p\,d\mu\right)^{\frac{1}{q}}\]
    Inégalité de Minkowski:
    \[\left(\int_X (f+g)^p\,d\mu\right)^{\frac{1}{p}}\le \left(\int_X f^p\,d\mu\right)^{\frac{1}{p}}\left(\int_X g^p\,d\mu\right)^{\frac{1}{p}}\]

    \begin{proof}
        cf poly du cours de M1 2024/2025
    \end{proof}

    \subsubsection*{Les espaces $l^p$}

    Soit $x=(x_n)_n\in \K^\N$ et $p\in [1,\infty]$. On définit $\norme{x}_p\in [0,\infty]$ par:
    \[\norme{x}_p=\cho{\sup \{|x_n|,n\in \N\}\text{ pour $p=+\infty$}\\ \left(\sum_{k=0}^{n}|x_k|^p\right)^{\frac{1}{p}}\text{ pour $p<+\infty$}}\]
    On note $l^p(\N,\K)$ ou juste $l^p$ le sous-ensemble des suites $x=(x_n)_\nN$ tel que $\norme{x}_p<+\infty$.\\
    Les inégalités de Hölder et Minkowski, appliquées à la mesure de comptage (en prolongeant à $p=1$ et $p=\infty$) donne, pour $p,q\in [1,+\infty],\frac{1}{p}+\frac{1}{q}=1$; $x,z\in l^p$ et $y\in l^q$\\
    Hölder: $x \cdot y\in l^1$ et $\norme{x\cdot y}_{l^1}\le \norme{x}_p \norme{y}_q$\\
    Minkowski: $\norme{x+z}_p\le \norme{x}_p +\norme{z}_p$

    \begin{tm}{}{}
        (Rappel) Pour $p\in [0,+\infty]$
        \begin{enumerate}
            \item $l^p(\N,\K)$ un \sev de $\K^\N$
            \item L'application $x\mapsto \norme{x}_p$ est une norme sur $l^p(\N,\K)$
            \item $l^p(\N,\K)$ muni de la norme $\norme{\cdot}_p$ est un espace de Banach.
        \end{enumerate}
    \end{tm}

    \begin{proof}
        Laisser en exercice
    \end{proof}

    \begin{prop}{}{}
        Soit $a,b\in [1,\infty]$ tels que $a<b$ alors $l^a(\N,\K)\subset l^b(\N,\K)$ et pour tout $x\in l^a(\N,\K)$, on a $\norme{x}_b\le \norme{x}_a$. Donc l'inclusion est continue et de norme $1$.
    \end{prop}

    \begin{proof}
        C'est vraie dans $\K^n$, ou de façon équivalente pour $x\in \K^\N$ dont tous les termes sont nuls à partir du rang $n-1$. En effet, notons $(l_k)_{k\in \llbracket 1,n\rrbracket}$ la base canonique de $\K^n$:
        \begin{itemize}
            \item Si $a=1,x=\sum_{k=1}^{n}x_k l_k,\norme{x}_b\le \sum_{k=1}^{n}|x_k|\norme{l_k}_b=\norme{x}_1$ puisque $\norme{l_k}_b=1$
            \item En général, soit $c=\frac{b}{a}\in ]1,+\infty[$ et $y=(|x_1|^\frac{1}{k},\ldots,|x_k|^\frac{1}{k})$ de sorte que $\norme{y}_a^a=\norme{x}_1$ et $\norme{y}_b^a=\norme{x}_c$. Le résultat précédent assure que $\norme{x}_c\le \norme{x}_1$ et permet de conclure.
        \end{itemize}
        Le cas de $\K^\N$: On considère $\fundfto{p_n}{\K^\N}{\K^n}{(x_k)_{k\in \N}}{(x_0,\ldots,x_{n-1})}$.\\
        Pour $x\in \K^\N$ et $\nN, \norme{\mu_n(x)}_b\le \norme{p_n(x)}_a$. En passant à la limite quand $n\to \infty$, on trouve $\norme{x}_p\le \norme{x}_a$.\\
        De plus pour la suite dont tous les termes sont nuls sauf le premier, on a: $\norme{l_0}_a=1=\norme{l_0}_b$. Donc l'inclusion $l^a\subset l^b$ est continue de norme $1$.
    \end{proof}

    On note $K(\N,\K)$ le \sev de $\K^\N$ constitué des suites dont tous les termes sont nuls à partir d'un certain rang et $\K_0^\N$ le \sev des suites qui convergent vers $0$.

    \begin{prop}{}{}
        Pour tout $p\in [1,+\infty[,K(\N,\K)$ est dense dans $l^p$ et l'adhérence de $K(\N,\K)$ dans $l^\infty(\N,\K)$ est $\K_0^\N$.
    \end{prop}

    \begin{proof}
        Laisser en exercice
    \end{proof}

    \begin{re}{}{}
        Une conséquence immédiate de cette proposition est que $l^p(\N,\K)$ est séparable si $p<+\infty$: Le sous-ensemble de $K(\N,\K)$ constitué des suites dont tous les termes sont dans $\Q+i \Q$ (Si $\K=\C$) ou $Q$ (Si $\K=\R$) est dénombrable et dense dans $l^p$.\\
        On peut montrer que $l^\infty$ n'est pas séparable (Exo 13 TD7)
    \end{re}

    \begin{re}{}{}
        Pour $L^p$, on a la séparabilité si $\mu$ est $\sigma$-finie.
    \end{re}

    \begin{prop}{}{}
        Soit $1\le p\le +\infty$ et $q$ son exposant conjugué ($\frac{1}{p}+\frac{1}{q}=1$) pour $y=(y_n)_\nN\in l^q$, la forme linéaire $\funto{\varphi_y}{x=(x_n)_\nN\in l^p}{\sum_{n=0}^{+\infty}x_n y_n}$ est continue.\\
        L'application $\fundfto{\tau}{l^q}{(l^p)'}{y}{\varphi_y}$ est linéaire et isométrique $\norme{\varphi_y}_{(l^p)'}=\norme{y}_q$
    \end{prop}

    \begin{re}{}{}
        On a le même théorème dans $L^p(X,\mu)$ où pour $g\in L^q(X)$ et $f\in L^p(X)$,
        \[\varphi_g(f)=\int_x fg\, d\mu\]
    \end{re}

    \begin{proof}
        Soit $y=(y_n)_\nN\in l^q$ et $x=(x_n)_\nN\in l^p$, on a
        \[\varphi_y(x)=\left|\sum_{n=0}^{+\infty}x_n y_n\right|\le \sum_{n=0}^{+\infty}\left| x_n y_n\right|=\norme{xy}_1\le \norme{x}_p \norme{y}_q\]
        En particulier, $\sum x_n y_n$ est normalement convergente donc convergente dans $\K$, $\varphi_y$ est continue et $\norme{\varphi_y}_{(l^p)'}\le \norme{y}_q$.\\
        L'application $\tau$ est évidemment linéaire. Montrons que qu'elle est isométrique.\\
        On suppose $y\neq 0_A$ et que $p\in ]1,+\infty[$ (Les cas extrêmes se font de façon analogue).\\
        Soit $z=(z_n)_\nN$ donnée par: $z_n=\cho{\frac{\overline{y_n}}{\left| y_n \right| }\left| y_n\right|^{q-1}\text{ si $y_n\neq 0$}\\ 0\text{ si }y_n =0}$\\
        Puisque $q=p(q-1)$ pour tout $n\in \N, \left|z_n \right|^p=\left| y_n \right|^q$ donc $z\in l^p$ et $\norme{z}_p=\norme{y}^{\frac{q}{p}_q}$. \\ 
        De plus, $\left| \varphi_y(z)\right|=\sum_{n=0}^{+\infty} \frac{\overline{y_n}y_n}{\left|y_n\right|}\left|y_n\right|^{q-1}=\sum_{n=0}^{\infty} \left| y_n\right|^q$.\\
        Comme $\left|\varphi_y (z)\right| =\norme{\varphi_y}_{(l^p)'}=\norme{z}_p$, on a
        \[\norme{y}_p^q\le \norme{\varphi_y}_{(l^p)'}\norme{y}_q^{\frac{q}{p}}\text{ et donc }\norme{\varphi_y}_{(l^p)'}\ge \norme{y}_q^{q-\frac{q}{p}}=\norme{y}_a\] 
    \end{proof}

    \begin{tm}{Théorème de Dualité}{}
        Soit $1\le p<+\infty$ et $q$ l'exposant conjugué $\frac{1}{p}+ \frac{1}{q}=1$. L'application $\fundfto{\tau}{l^q}{(l^p)'}{y}{\varphi_y}$ est un isomorphisme isométrique
    \end{tm}

    \begin{re}{}{}
        On a le même résultat dans $L^p(X)$ à condition que $\mu$ soit $\sigma$-finie.
    \end{re}

    \begin{re}{}{}
        Le théorème est faux pour $p=+\infty$. Dans ce cas, l'image de $l^1$ par $\tau$ est strictement plus petite que $(l^\infty)'$
    \end{re}

    \begin{proof}
        Il reste à voir que $\tau$ est surjective auquel cas elle sera bi-continue. On se place dans le cas $p>1$.\\
        Soit $(l^k)_{k\in \N}$ la base canonique de $K(\N,\K)$. $l^k=(\delta_n^k)_{\nN}$. Soit $f\in (l^n)'$ non-nulle, si il existe $y\in l^q$ telle que $f=\varphi_y$, alors nécessairement $\forall k\in \N, f(l^k)=\varphi_y(l^k)=y_k$. On considère $y=(y_k)_{k\in \N}$ ainsi définie. Montrons que $y\in l^q$.\\
        Pour $j\in \N$, soit $z^{j}=(z_n^j)_\nN, K(\N,\K), z^j=\sum_{n=0}^{j}\left| y_n\right|^{q_2}, \overline{y_n}\cdot l^n$.\\
        On a \begin{align*}
            \varphi(z^j)&=\varphi (\sum_{n=0}^{+\infty}\left|y_n\right|^{q-2}\overline{y_n}\cdot l^n)\\
            &=\sum_{n=0}^{+\infty}\left|y_n\right|^{q-2}\overline{y_n}\varphi(l^n)=\sum_{n=0}^{+\infty}\left|y_n\right|^q
        \end{align*}
        \[\norme{z}_p= \left(\sum_{n=0}^{p}\left|y_n\right|^{(q-2)}\left|y_n\right|^p\right)^{\frac{1}{p}}=\left(\sum_{n=0}^{+\infty} \left|y_n\right|^q\right)^{\frac{1}{p}}\]
        On en déduit
        \[\left|\varphi(z^p)\right|=\sum_{n=0}^{+\infty}\left|y_n\right|\le \norme{\varphi}_p \norme{y^{q}}_p=\norme{\varphi}_p\left(\sum_{n=0}^{+\infty}\left|y_n\right|^q\right)^{\frac{1}{p}}\]
        D'où
        \[\norme{\varphi}_{(l^p)'}\le \left( \sum_{n=0}^{+\infty}\right)^{\frac{1}{q}}\]
        On fait tendre $j\to \infty$ et on obtient $\norme{y}_q\le \norme{\varphi}_{(l^p)'}$ donc $y\in l^q$.\\
        Puisque $\varphi$ et $\varphi_y$ coïncident sur les $l^k$ pour tout $k\in \N$, par linéarité elles coïncident sur $K(\N,\K)$ qui est dense dans $l^p$ puisque $p\neq \infty$. Par continuité, $\varphi=\varphi_y$
    \end{proof}

    \begin{cor}{}{}
        Pour $1<p<+\infty$, $l^p$ est réflexif. Si $q$ est l'exposant conjugué, $1<q<\infty$:
        \[\fundfto{J}{l^p}{(l^p)''}{x}{\fundfto{J(x)}{(l^p)'\cong l^q}{\K}{y}{\varphi_y(x)=\varphi_x(y)=\sum_{n=0}^{+\infty} x_n y_n}}\]
        est un isomorphisme isométrique.
    \end{cor}

    \begin{proof}
        Laisser en exercice
    \end{proof}

    \begin{cor}{}{}
        Toute suite bornée dans $l^p$ pour $1<p<+\infty$ admet une sous-suite qui converge faiblement
    \end{cor}

    \begin{proof}
        Laisser en exercice
    \end{proof}

    \begin{cor}{}{}
        Soit $1\le p\le +\infty$ et $q$ l'exposant conjugué. Soit $(x^n)_{\nN}$ une suite de $l^p$ et $x\in l^p$. On considère les ennoncés:
        \begin{enumerate}
            \item $x_n \rightharpoonup x$
            \item $\forall y\in l^q, \sum_{k=0}^{n} x_k y_k \underset{k\to \infty}{\longrightarrow} \sum_{k=0}^{\infty} x_k y_k$
            \item $(\varphi_{x^n})\overset{*}{\longrightarrow} \varphi_x$
        \end{enumerate}
        On a:
        \begin{itemize}
            \item Pour $1<p<+\infty$: $(1)\iff (2)\iff (3)$
            \item Pour $p=1$: $(1)\iff (2)$
            \item Pour $p=+\infty$: $(2)\iff (3)$
        \end{itemize}
    \end{cor}

    \begin{proof}
        Laisser en exercice
    \end{proof}
\end{document}